%%% element catalog, generated on Mon May 18 22:06:16 CEST 2026

%%% WARNING %%%
%%% This file was automatically generated by SAPlugin,
%%% and will be overwritten next time.

\chapter{Element catalog}\label{sec:catalog}
% COMPONENTS
\section{Components}\label{sec:components}
\subsection{Av2AcknowledgementService}\label{comp:ComponentsAv2AcknowledgementService}
	\begin{description}[noitemsep,nolistsep]
		\item[Responsibility:]~After RoutingManager confirms downstream delivery (e.g.,
MessageBroker accepts the message), DeliveryAcknowledger emits an
acknowledgement back to the originating PatientGateway, allowing the
gateway's RetryAndSynchronizationService to clear delivered entries from
LocalBufferingRepository.
		\item[Super-components:]~None
		\item[Sub-components:]~None
		\item[Provided interfaces:]~\iconprovided{}~\vpett{\nameref{int:InterfacesAv2DeliveryAck}}
		\item[Required interfaces:]~\iconrequired{}~\vpett{\nameref{int:InterfacesAv2BrokerCommitNotification}}
		\item[Deployed on:]~\faSquareO~\vpett{\nameref{node:DeploymenteHealthPlatform}}, \faSquareO~\vpett{\nameref{node:DeploymentBackendGatewayNode}}, \faSquareO~\vpett{\nameref{node:DeploymentBackendGatewayNodePilotDeployment}}, \faSquareO~\vpett{\nameref{node:DeploymentDevBackendNodeDevDeployment}}
		\item[Visible on diagrams:]~\cref{diag:Component:Primarydiagramoftheclientserverview,diag:Deployment:Contextdiagramforthedeploymentview,diag:Deployment:DevelopmentTestDeployment20patients3clinmodels,diag:Deployment:PilotDeployment200patients5clinmodels,diag:Deployment:Primarydeploymentdiagram}		
	\end{description}

\subsection{Av2AdminPortal}\label{comp:ComponentsAv2AdminPortal}
	\begin{description}[noitemsep,nolistsep]
		\item[Responsibility:]~Provides the system administrator with a web-based interface for monitoring the operational status of the PMS, receiving failure notifications, and triggering redeployment or rollback of failed internal subsystems. It consumes the RedeploymentControl interface to initiate and monitor redeployment jobs, and receives notifications from the MonitoringService via the NotificationDispatch interface. It is the only component through which a system administrator can interact with the RedeploymentCoordinator.
		\item[Super-components:]~None
		\item[Sub-components:]~None
		\item[Provided interfaces:]~None
		\item[Required interfaces:]~\iconrequired{}~\vpett{\nameref{int:InterfacesAv2NotificationDispatch}}, \iconrequired{}~\vpett{\nameref{int:InterfacesAv2RedeploymentControl}}, \iconrequired{}~\vpett{\nameref{int:InterfacesAv2SLAReporting}}
		\item[Deployed on:]~\faSquareO~\vpett{\nameref{node:DeploymentAdminWorkstation}}
		\item[Visible on diagrams:]~\cref{diag:Component:Primarydiagramoftheclientserverview,diag:Deployment:Contextdiagramforthedeploymentview}		
	\end{description}

\subsection{Av2CommunicationGateway}\label{comp:ComponentsAv2CommunicationGateway}
	\begin{description}[noitemsep,nolistsep]
		\item[Responsibility:]~The single backend ingress point for all primary-channel
patient gateway traffic. Authenticates incoming transmissions, routes
them to downstream services (MessageBroker, EmergencyNotificationService),
detects communication failures, coordinates retries with exponential
backoff, tracks per-gateway heartbeats for missing-data detection, and
acknowledges successfully dispatched messages back to the originating
gateway.
		\item[Super-components:]~None
		\item[Sub-components:]~None
		\item[Provided interfaces:]~\iconprovided{}~\vpett{\nameref{int:InterfacesAv2AvailabilityMonitoring}}, \iconprovided{}~\vpett{\nameref{int:InterfacesAv2BufferedDataDispatch}}, \iconprovided{}~\vpett{\nameref{int:InterfacesAv2CommunicationGatewayHealthCheck}}
		\item[Required interfaces:]~\iconrequired{}~\vpett{\nameref{int:InterfacesAv2GatewayDataIngress}}, \iconrequired{}~\vpett{\nameref{int:InterfacesAv2PatientStatusOverview}}
		\item[Deployed on:]~\faSquareO~\vpett{\nameref{node:DeploymenteHealthPlatform}}, \faSquareO~\vpett{\nameref{node:DeploymentBackendGatewayNode}}, \faSquareO~\vpett{\nameref{node:DeploymentBackendGatewayNodePilotDeployment}}, \faSquareO~\vpett{\nameref{node:DeploymentDevBackendNodeDevDeployment}}
		\item[Visible on diagrams:]~\cref{diag:Component:DecompositionoftheAv2CommunicationGateway,diag:Component:Primarydiagramoftheclientserverview,diag:Deployment:Contextdiagramforthedeploymentview,diag:Deployment:DevelopmentTestDeployment20patients3clinmodels,diag:Deployment:PilotDeployment200patients5clinmodels,diag:Deployment:Primarydeploymentdiagram}		
	\end{description}

\subsection{Av2DataIngestionService}\label{comp:ComponentsAv2DataIngestionService}
	\begin{description}[noitemsep,nolistsep]
		\item[Responsibility:]~Consumes validated sensor data and emergency notifications from the MessageBroker and persists them through the SensorDataRepository. Acts as the single write path for incoming patient data into the backend persistence layer, decoupling the ingestion rate from the downstream risk-estimation pipeline.
		\item[Super-components:]~None
		\item[Sub-components:]~None
		\item[Provided interfaces:]~None
		\item[Required interfaces:]~\iconrequired{}~\vpett{\nameref{int:InterfacesAv2QueuedSensorDataMgmt}}, \iconrequired{}~\vpett{\nameref{int:InterfacesSensorDataMgmt}}
		\item[Deployed on:]~\faSquareO~\vpett{\nameref{node:DeploymentBackendGatewayNode}}, \faSquareO~\vpett{\nameref{node:DeploymentBackendGatewayNodePilotDeployment}}, \faSquareO~\vpett{\nameref{node:DeploymentDevBackendNodeDevDeployment}}
		\item[Visible on diagrams:]~\cref{diag:Component:Primarydiagramoftheclientserverview,diag:Deployment:DevelopmentTestDeployment20patients3clinmodels,diag:Deployment:PilotDeployment200patients5clinmodels,diag:Deployment:Primarydeploymentdiagram}		
	\end{description}

\subsection{Av2EmergencyBackupReceiver}\label{comp:ComponentsAv2EmergencyBackupReceiver}
	\begin{description}[noitemsep,nolistsep]
		\item[Responsibility:]~Receives emergency notifications transmitted via the backup
SMS gateway when the primary channel is unavailable. Authenticates the
sending patient gateway by a pre-registered device token embedded in the
SMS payload, and forwards validated emergencies to the MessageBroker for
the same downstream processing as primary-channel emergencies.
		\item[Super-components:]~None
		\item[Sub-components:]~None
		\item[Provided interfaces:]~\iconprovided{}~\vpett{\nameref{int:InterfacesAv2BackupEmergencyIngress}}
		\item[Required interfaces:]~\iconrequired{}~\vpett{\nameref{int:InterfacesAv2QueuedSensorDataMgmt}}
		\item[Deployed on:]~\faSquareO~\vpett{\nameref{node:DeploymenteHealthPlatform}}, \faSquareO~\vpett{\nameref{node:DeploymentBackendGatewayNode}}, \faSquareO~\vpett{\nameref{node:DeploymentBackendGatewayNodePilotDeployment}}, \faSquareO~\vpett{\nameref{node:DeploymentDevBackendNodeDevDeployment}}
		\item[Visible on diagrams:]~\cref{diag:Component:Primarydiagramoftheclientserverview,diag:Deployment:Contextdiagramforthedeploymentview,diag:Deployment:DevelopmentTestDeployment20patients3clinmodels,diag:Deployment:PilotDeployment200patients5clinmodels,diag:Deployment:Primarydeploymentdiagram}		
	\end{description}

\subsection{Av2MessageBroker}\label{comp:ComponentsAv2MessageBroker}
	\begin{description}[noitemsep,nolistsep]
		\item[Responsibility:]~Buffers validated sensor-data packages and emergency
notifications between the CommunicationGateway and the
DataIngestionService. Provides durability and decoupling so that
transient failures of downstream services do not lead to data loss, and
exposes a health-check endpoint for the InternalFailureDetector.
		\item[Super-components:]~None
		\item[Sub-components:]~None
		\item[Provided interfaces:]~\iconprovided{}~\vpett{\nameref{int:InterfacesAv2BrokerCommitNotification}}, \iconprovided{}~\vpett{\nameref{int:InterfacesAv2MessageBrokerHealthCheck}}, \iconprovided{}~\vpett{\nameref{int:InterfacesAv2QueuedSensorDataMgmt}}
		\item[Required interfaces:]~\iconrequired{}~\vpett{\nameref{int:InterfacesAv2BufferedDataDispatch}}
		\item[Deployed on:]~\faSquareO~\vpett{\nameref{node:DeploymenteHealthPlatform}}, \faSquareO~\vpett{\nameref{node:DeploymentBackendGatewayNode}}, \faSquareO~\vpett{\nameref{node:DeploymentBackendGatewayNodePilotDeployment}}, \faSquareO~\vpett{\nameref{node:DeploymentDevBackendNodeDevDeployment}}
		\item[Visible on diagrams:]~\cref{diag:Component:Primarydiagramoftheclientserverview,diag:Deployment:Contextdiagramforthedeploymentview,diag:Deployment:DevelopmentTestDeployment20patients3clinmodels,diag:Deployment:PilotDeployment200patients5clinmodels,diag:Deployment:Primarydeploymentdiagram}		
	\end{description}

\subsection{Av2MonitoringService}\label{comp:ComponentsAv2MonitoringService}
	\begin{description}[noitemsep,nolistsep]
		\item[Responsibility:]~Central observability and recovery-orchestration component.
Detects (a) missing sensor data, (b) internal subsystem crashes, (c) SLA
violations of the telecom channel; classifies failures by affected patient
risk level; notifies administrators and patient contacts within Av2
deadlines; records communication-gap durations; and coordinates recovery
with the RecoverySyncService and RedeploymentCoordinator.
		\item[Super-components:]~None
		\item[Sub-components:]~None
		\item[Provided interfaces:]~\iconprovided{}~\vpett{\nameref{int:InterfacesAv2NotificationDispatch}}, \iconprovided{}~\vpett{\nameref{int:InterfacesAv2RecoveryCoordination}}, \iconprovided{}~\vpett{\nameref{int:InterfacesAv2SLAReporting}}
		\item[Required interfaces:]~\iconrequired{}~\vpett{\nameref{int:InterfacesAv2AvailabilityMonitoring}}, \iconrequired{}~\vpett{\nameref{int:InterfacesAv2CommunicationGatewayHealthCheck}}, \iconrequired{}~\vpett{\nameref{int:InterfacesAv2MessageBrokerHealthCheck}}, \iconrequired{}~\vpett{\nameref{int:InterfacesAv2PatientContactLookup}}
		\item[Deployed on:]~\faSquareO~\vpett{\nameref{node:DeploymenteHealthPlatform}}, \faSquareO~\vpett{\nameref{node:DeploymentMonitoringNode}}, \faSquareO~\vpett{\nameref{node:DeploymentDevBackendNodeDevDeployment}}, \faSquareO~\vpett{\nameref{node:DeploymentMonitoringNodePilotDeployment}}
		\item[Visible on diagrams:]~\cref{diag:Component:DecompositionoftheAv2MonitoringService,diag:Component:Primarydiagramoftheclientserverview,diag:Deployment:Contextdiagramforthedeploymentview,diag:Deployment:DevelopmentTestDeployment20patients3clinmodels,diag:Deployment:PilotDeployment200patients5clinmodels,diag:Deployment:Primarydeploymentdiagram}		
	\end{description}

\subsection{Av2PatientGateway}\label{comp:ComponentsAv2PatientGateway}
	\begin{description}[noitemsep,nolistsep]
		\item[Responsibility:]~The patient-facing edge component running on the wearable/handheld device. Acquires raw sensor readings, transmits them to the backend CommunicationGateway over the primary telecom channel, monitors device battery and network health, manages local buffering during connectivity loss, and dispatches emergency notifications via the backup channel when the primary channel is unavailable. It is the only component that physically resides with the patient.
		\item[Super-components:]~None
		\item[Sub-components:]~None
		\item[Provided interfaces:]~\iconprovided{}~\vpett{\nameref{int:InterfacesAv2EmergencyDispatch}}, \iconprovided{}~\vpett{\nameref{int:InterfacesAv2GatewayDataIngress}}
		\item[Required interfaces:]~\iconrequired{}~\vpett{\nameref{int:InterfacesAv2BackupEmergencyIngress}}, \iconrequired{}~\vpett{\nameref{int:InterfacesAv2DeliveryAck}}, \iconrequired{}~\vpett{\nameref{int:InterfacesAv2GatewayResyncCommand}}
		\item[Deployed on:]~\faSquareO~\vpett{\nameref{node:DeploymentGatewayNode}}, \faSquareO~\vpett{\nameref{node:DeploymentDevGatewayNodeDevDeployment}}, \faSquareO~\vpett{\nameref{node:DeploymentGatewayNodePilotDeployment}}, \faSquareO~\vpett{\nameref{node:DeploymentPatientGatewayDevice}}
		\item[Visible on diagrams:]~\cref{diag:Component:DecompositionoftheAv2PatientGateway,diag:Component:Primarydiagramoftheclientserverview,diag:Deployment:Contextdiagramforthedeploymentview,diag:Deployment:DevelopmentTestDeployment20patients3clinmodels,diag:Deployment:PilotDeployment200patients5clinmodels,diag:Deployment:Primarydeploymentdiagram}		
	\end{description}

\subsection{Av2RecoverySyncService}\label{comp:ComponentsAv2RecoverySyncService}
	\begin{description}[noitemsep,nolistsep]
		\item[Responsibility:]~Coordinates post-failure recovery. When RecoveryTracker
signals a failure has been resolved, RecoverySyncService (a) drains any
sensor data queued in the MessageBroker that arrived during degraded mode,
and (b) instructs the affected PatientGateway(s) to flush their
LocalBufferingRepository through a GatewayResyncCommand.
		\item[Super-components:]~None
		\item[Sub-components:]~None
		\item[Provided interfaces:]~\iconprovided{}~\vpett{\nameref{int:InterfacesAv2GatewayResyncCommand}}
		\item[Required interfaces:]~\iconrequired{}~\vpett{\nameref{int:InterfacesAv2QueuedSensorDataMgmt}}, \iconrequired{}~\vpett{\nameref{int:InterfacesAv2RecoveryCoordination}}
		\item[Deployed on:]~\faSquareO~\vpett{\nameref{node:DeploymenteHealthPlatform}}, \faSquareO~\vpett{\nameref{node:DeploymentMonitoringNode}}, \faSquareO~\vpett{\nameref{node:DeploymentDevBackendNodeDevDeployment}}, \faSquareO~\vpett{\nameref{node:DeploymentMonitoringNodePilotDeployment}}
		\item[Visible on diagrams:]~\cref{diag:Component:Primarydiagramoftheclientserverview,diag:Deployment:Contextdiagramforthedeploymentview,diag:Deployment:DevelopmentTestDeployment20patients3clinmodels,diag:Deployment:PilotDeployment200patients5clinmodels,diag:Deployment:Primarydeploymentdiagram}		
	\end{description}

\subsection{Av2RedeploymentCoordinator}\label{comp:ComponentsAv2RedeploymentCoordinator}
	\begin{description}[noitemsep,nolistsep]
		\item[Responsibility:]~Enables the system administrator to trigger redeployment or rollback of the CommunicationGateway within 5 minutes. Monitors rollback progress and confirms success.
		\item[Super-components:]~None
		\item[Sub-components:]~None
		\item[Provided interfaces:]~\iconprovided{}~\vpett{\nameref{int:InterfacesAv2RedeploymentControl}}
		\item[Required interfaces:]~\iconrequired{}~\vpett{\nameref{int:InterfacesAv2RecoveryCoordination}}
		\item[Deployed on:]~\faSquareO~\vpett{\nameref{node:DeploymenteHealthPlatform}}, \faSquareO~\vpett{\nameref{node:DeploymentAdminNode}}, \faSquareO~\vpett{\nameref{node:DeploymentAdminNodePilotDeployment}}, \faSquareO~\vpett{\nameref{node:DeploymentDevBackendNodeDevDeployment}}
		\item[Visible on diagrams:]~\cref{diag:Component:Primarydiagramoftheclientserverview,diag:Deployment:Contextdiagramforthedeploymentview,diag:Deployment:DevelopmentTestDeployment20patients3clinmodels,diag:Deployment:PilotDeployment200patients5clinmodels,diag:Deployment:Primarydeploymentdiagram}		
	\end{description}

\subsection{Av2SensorDataRepository}\label{comp:ComponentsAv2SensorDataRepository}
	\begin{description}[noitemsep,nolistsep]
		\item[Responsibility:]~Persistent store for incoming sensor data packages and their arrival timestamps. Provides the SensorDataMgmt interface for appending new readings and for querying historical sensor data needed by clinical models that require historical context.
		\item[Super-components:]~None
		\item[Sub-components:]~None
		\item[Provided interfaces:]~\iconprovided{}~\vpett{\nameref{int:InterfacesSensorDataMgmt}}
		\item[Required interfaces:]~None
		\item[Deployed on:]~\faSquareO~\vpett{\nameref{node:DeploymentDevPatientRecordsNodeDevDeployment}}, \faSquareO~\vpett{\nameref{node:DeploymentPatientRecordsNodePilotDeployment}}, \faSquareO~\vpett{\nameref{node:DeploymentPatientRecordsNode}}
		\item[Visible on diagrams:]~\cref{diag:Component:Primarydiagramoftheclientserverview,diag:Deployment:DevelopmentTestDeployment20patients3clinmodels,diag:Deployment:PilotDeployment200patients5clinmodels,diag:Deployment:Primarydeploymentdiagram}		
	\end{description}

\subsection{ClinicalJobCreator}\label{comp:ComponentseHealthPlatformClinicalJobCreator}
	\begin{description}[noitemsep,nolistsep]
		\item[Responsibility:]~The \vpett{\nameref{comp:ComponentseHealthPlatformClinicalJobCreator}} is responsible for constructing the \ref{data:DatatypesClinicalModelJob} objects and populating them with the necessary patient health data for execution by the \vpett{\nameref{comp:ComponentseHealthPlatformRiskEstimationProcessor}}s.
		\item[Super-components:]~\iconcomponent{}~\vpett{\nameref{comp:ComponentseHealthPlatform}}
		\item[Sub-components:]~None
		\item[Provided interfaces:]~\iconprovided{}~\vpett{\nameref{int:InterfacesClinicalJobMgmt}}
		\item[Required interfaces:]~\iconrequired{}~\vpett{\nameref{int:InterfacesFetchClinicalModels}}, \iconrequired{}~\vpett{\nameref{int:InterfacesOtherDataMgmt}}, \iconrequired{}~\vpett{\nameref{int:InterfacesSensorDataMgmt}}
		\item[Deployed on:]~\faSquareO~\vpett{\nameref{node:DeploymentRiskEstimationManagementNodePilotDeployment}}, \faSquareO~\vpett{\nameref{node:DeploymentRiskEstimationManagmentNodeDevDeployment}}, \faSquareO~\vpett{\nameref{node:DeploymenteHealthPlatformRiskEstimationMgmtNode}}
		\item[Visible on diagrams:]~\cref{diag:Component:Primarydiagramoftheclientserverview,diag:Deployment:DevelopmentTestDeployment20patients3clinmodels,diag:Deployment:PilotDeployment200patients5clinmodels,diag:Deployment:Primarydeploymentdiagram,diag:Interaction:Riskestimationprocess}		
	\end{description}

\subsection{ClinicalModelCache}\label{comp:ComponentseHealthPlatformClinicalModelCache}
	\begin{description}[noitemsep,nolistsep]
		\item[Responsibility:]~The \vpett{\nameref{comp:ComponentseHealthPlatformClinicalModelCache}} is a read-through cache responsible for caching the \ref{data:DatatypesClinicalModel}s that should be evaluated for a certain patient and their configurations. This cache is located close to the \vpett{\nameref{comp:ComponentseHealthPlatformRiskEstimationProcessor}} and \vpett{\nameref{comp:ComponentseHealthPlatformRiskEstimationCombiner}} in order to improve the latency of the whole risk estimation flow. The items in the \vpett{\nameref{comp:ComponentseHealthPlatformClinicalModelCache}} do not expire over time, but should be invalided explicitly if needed.
Cached models remain available for the \vpett{\nameref{comp:ComponentseHealthPlatformRiskEstimationProcessor}} while updates are made in the \vpett{\nameref{comp:ComponentseHealthPlatformClinicalModelDB}}.
		\item[Super-components:]~\iconcomponent{}~\vpett{\nameref{comp:ComponentseHealthPlatform}}
		\item[Sub-components:]~None
		\item[Provided interfaces:]~\iconprovided{}~\vpett{\nameref{int:InterfacesClinicalModelCacheMgmt}}, \iconprovided{}~\vpett{\nameref{int:InterfacesFetchClinicalModels}}
		\item[Required interfaces:]~\iconrequired{}~\vpett{\nameref{int:InterfacesClinicalModelStorage}}, \iconrequired{}~\vpett{\nameref{int:InterfacesOtherDataMgmt}}
		\item[Deployed on:]~\faSquareO~\vpett{\nameref{node:DeploymentRiskEstimationManagementNodePilotDeployment}}, \faSquareO~\vpett{\nameref{node:DeploymentRiskEstimationManagmentNodeDevDeployment}}, \faSquareO~\vpett{\nameref{node:DeploymenteHealthPlatformRiskEstimationMgmtNode}}
		\item[Visible on diagrams:]~\cref{diag:Component:Primarydiagramoftheclientserverview,diag:Deployment:DevelopmentTestDeployment20patients3clinmodels,diag:Deployment:PilotDeployment200patients5clinmodels,diag:Deployment:Primarydeploymentdiagram,diag:Interaction:Riskestimationprocess}		
	\end{description}

\subsection{ClinicalModelDB}\label{comp:ComponentseHealthPlatformClinicalModelDB}
	\begin{description}[noitemsep,nolistsep]
		\item[Responsibility:]~The \vpett{\nameref{comp:ComponentseHealthPlatformClinicalModelDB}} stores all the \ref{data:DatatypesClinicalModel}s and the configurations of these \ref{data:DatatypesClinicalModel}s for the different patients separately from other data.
It only allows appending new \ref{data:DatatypesClinicalModel}s.
		\item[Super-components:]~\iconcomponent{}~\vpett{\nameref{comp:ComponentseHealthPlatform}}
		\item[Sub-components:]~None
		\item[Provided interfaces:]~\iconprovided{}~\vpett{\nameref{int:InterfacesClinicalModelStorage}}, \iconprovided{}~\vpett{\nameref{int:InterfacesP1ClinicalConfigMgmt}}
		\item[Required interfaces:]~None
		\item[Deployed on:]~\faSquareO~\vpett{\nameref{node:DeploymentTODONodePilotDeployment}}, \faSquareO~\vpett{\nameref{node:DeploymentRiskEstimationManagmentNodeDevDeployment}}, \faSquareO~\vpett{\nameref{node:DeploymenteHealthPlatformTODONode}}
		\item[Visible on diagrams:]~\cref{diag:Component:Primarydiagramoftheclientserverview,diag:Deployment:DevelopmentTestDeployment20patients3clinmodels,diag:Deployment:PilotDeployment200patients5clinmodels,diag:Deployment:Primarydeploymentdiagram,diag:Interaction:Riskestimationprocess}		
	\end{description}

\subsection{eHealth Platform}\label{comp:ComponentseHealthPlatform}
	\begin{description}[noitemsep,nolistsep]
		\item[Responsibility:]~The parent \vpett{\nameref{comp:ComponentseHealthPlatform}} component that represents the context boundary in the client-server view.
		\item[Super-components:]~None
		\item[Sub-components:]~\iconcomponent{}~\vpett{\nameref{comp:ComponentseHealthPlatformOtherFunctionality}}, \iconcomponent{}~\vpett{\nameref{comp:ComponentseHealthPlatformMLModelManager}}, \iconcomponent{}~\vpett{\nameref{comp:ComponentseHealthPlatformClinicalModelCache}}, \iconcomponent{}~\vpett{\nameref{comp:ComponentseHealthPlatformRiskEstimationCombiner}}, \iconcomponent{}~\vpett{\nameref{comp:ComponentseHealthPlatformRiskEstimationProcessor}}, \iconcomponent{}~\vpett{\nameref{comp:ComponentseHealthPlatformRiskEstimationScheduler}}, \iconcomponent{}~\vpett{\nameref{comp:ComponentseHealthPlatformClinicalJobCreator}}, \iconcomponent{}~\vpett{\nameref{comp:ComponentseHealthPlatformClinicalModelDB}}
		\item[Provided interfaces:]~None
		\item[Required interfaces:]~\iconrequired{}~\vpett{\nameref{int:InterfacesKVStore}}, \iconrequired{}~\vpett{\nameref{int:InterfacesMLaaSService}}
		\item[Deployed on:]~\faSquareO~\vpett{\nameref{node:DeploymentTODONodePilotDeployment}}, \faSquareO~\vpett{\nameref{node:DeploymentRiskEstimationManagementNodePilotDeployment}}, \faSquareO~\vpett{\nameref{node:DeploymentRiskEstimationProcessorNodePilotDeployment}}, \faSquareO~\vpett{\nameref{node:DeploymentMLaaSConnPilotDeployment}}, \faSquareO~\vpett{\nameref{node:DeploymentRiskEstimationManagmentNodeDevDeployment}}, \faSquareO~\vpett{\nameref{node:DeploymentRiskEstimationProcessorNodeDevDeployment}}, \faSquareO~\vpett{\nameref{node:DeploymenteHealthPlatformMLaaSConnectorNode}}, \faSquareO~\vpett{\nameref{node:DeploymenteHealthPlatformRiskEstimationMgmtNode}}, \faSquareO~\vpett{\nameref{node:DeploymenteHealthPlatformRiskEstimationProcessorNode}}, \faSquareO~\vpett{\nameref{node:DeploymenteHealthPlatformTODONode}}
		\item[Visible on diagrams:]~\cref{diag:Component:Contextdiagramfortheclientserverview,diag:Component:Contextdiagramforthedecompositionview}		
	\end{description}

\subsection{HIS}\label{comp:ComponentsHIS}
	\begin{description}[noitemsep,nolistsep]
		\item[Responsibility:]~Hospital Information System
		\item[Super-components:]~None
		\item[Sub-components:]~None
		\item[Provided interfaces:]~\iconprovided{}~\vpett{\nameref{int:InterfaceshealthAPI}}
		\item[Required interfaces:]~None
		\item[Deployed on:]~\faSquareO~\vpett{\nameref{node:DeploymentHISServer}}
		\item[Visible on diagrams:]~\cref{diag:Deployment:Contextdiagramforthedeploymentview}		
	\end{description}

\subsection{KVStorageService}\label{comp:ComponentsKVStorageService}
	\begin{description}[noitemsep,nolistsep]
		\item[Responsibility:]~Service for key-value storage of data.
		\item[Super-components:]~None
		\item[Sub-components:]~None
		\item[Provided interfaces:]~\iconprovided{}~\vpett{\nameref{int:InterfacesKVStore}}
		\item[Required interfaces:]~None
		\item[Deployed on:]~\faSquareO~\vpett{\nameref{node:DeploymentRiskEstimationProcessorNodePilotDeployment}}, \faSquareO~\vpett{\nameref{node:DeploymentRiskEstimationProcessorNodeDevDeployment}}, \faSquareO~\vpett{\nameref{node:DeploymenteHealthPlatformRiskEstimationProcessorNode}}
		\item[Visible on diagrams:]~\cref{diag:Component:Contextdiagramfortheclientserverview,diag:Deployment:DevelopmentTestDeployment20patients3clinmodels,diag:Deployment:PilotDeployment200patients5clinmodels,diag:Deployment:Primarydeploymentdiagram}		
	\end{description}

\subsection{MLaaS}\label{comp:ComponentsMLaaS}
	\begin{description}[noitemsep,nolistsep]
		\item[Responsibility:]~MachineLearning-as-a-Service cloud provider. This cloud service provides APIs for managing different \ref{data:DatatypesMLModel}s and exchanging \ref{data:DatatypesMLModel} updates with other parties.
		\item[Super-components:]~None
		\item[Sub-components:]~None
		\item[Provided interfaces:]~\iconprovided{}~\vpett{\nameref{int:InterfacesMLaaSService}}
		\item[Required interfaces:]~None
		\item[Deployed on:]~\faSquareO~\vpett{\nameref{node:DeploymentMLaaSCloud}}, \faSquareO~\vpett{\nameref{node:DeploymentMLaaSServicePilotDeployment}}, \faSquareO~\vpett{\nameref{node:DeploymentMLaaSServiceDevDeployment}}
		\item[Visible on diagrams:]~\cref{diag:Component:Contextdiagramfortheclientserverview,diag:Deployment:Contextdiagramforthedeploymentview,diag:Deployment:DevelopmentTestDeployment20patients3clinmodels,diag:Deployment:PilotDeployment200patients5clinmodels}		
	\end{description}

\subsection{MLModelManager}\label{comp:ComponentseHealthPlatformMLModelManager}
	\begin{description}[noitemsep,nolistsep]
		\item[Responsibility:]~This is component is responsible for managing the local \ref{data:DatatypesMLModel}s for making predictions for malignant events based on the incoming sensor data.
It keeps track of the available \ref{data:DatatypesMLModel}s that can be used and supports exchanging \ref{data:DatatypesMLModel} updates with other systems using the APIs of \vpett{\nameref{comp:ComponentsMLaaS}} providers.
		\item[Super-components:]~\iconcomponent{}~\vpett{\nameref{comp:ComponentseHealthPlatform}}
		\item[Sub-components:]~None
		\item[Provided interfaces:]~\iconprovided{}~\vpett{\nameref{int:InterfacesMLModelMgmt}}
		\item[Required interfaces:]~\iconrequired{}~\vpett{\nameref{int:InterfacesMLaaSService}}
		\item[Deployed on:]~\faSquareO~\vpett{\nameref{node:DeploymentMLaaSConnPilotDeployment}}, \faSquareO~\vpett{\nameref{node:DeploymentRiskEstimationManagmentNodeDevDeployment}}, \faSquareO~\vpett{\nameref{node:DeploymenteHealthPlatformMLaaSConnectorNode}}
		\item[Visible on diagrams:]~\cref{diag:Component:Contextdiagramfortheclientserverview,diag:Component:DecompositionoftheMLModelManager,diag:Component:Primarydiagramoftheclientserverview,diag:Deployment:Contextdiagramforthedeploymentview,diag:Deployment:DevelopmentTestDeployment20patients3clinmodels,diag:Deployment:PilotDeployment200patients5clinmodels,diag:Deployment:Primarydeploymentdiagram}		
	\end{description}

\subsection{OtherFunctionality}\label{comp:ComponentseHealthPlatformOtherFunctionality}
	\begin{description}[noitemsep,nolistsep]
		\item[Responsibility:]~Other functionality that still has to be worked out.
		\item[Super-components:]~\iconcomponent{}~\vpett{\nameref{comp:ComponentseHealthPlatform}}
		\item[Sub-components:]~None
		\item[Provided interfaces:]~None
		\item[Required interfaces:]~None
		\item[Deployed on:]~\faSquareO~\vpett{\nameref{node:DeploymentTODONodePilotDeployment}}, \faSquareO~\vpett{\nameref{node:DeploymentRiskEstimationManagmentNodeDevDeployment}}, \faSquareO~\vpett{\nameref{node:DeploymenteHealthPlatformTODONode}}
		\item[Visible on diagrams:]~\cref{diag:Deployment:DevelopmentTestDeployment20patients3clinmodels,diag:Deployment:PilotDeployment200patients5clinmodels,diag:Deployment:Primarydeploymentdiagram,diag:Interaction:CombiningClinicalModelJobresults,diag:Interaction:Incomingsensordata,diag:Interaction:Riskestimationprocess}		
	\end{description}

\subsection{P1HISAdapter}\label{comp:ComponentsP1HISAdapter}
	\begin{description}[noitemsep,nolistsep]
		\item[Responsibility:]~Outbound adapter for the external \vpett{\nameref{comp:ComponentsHIS}} healthAPI. Used by \vpett{\nameref{mod:ComponentsP1PatientDataServiceP1EHRProxyModule}} for UC16 reads and UC17 writes.
		\item[Super-components:]~None
		\item[Sub-components:]~None
		\item[Provided interfaces:]~\iconprovided{}~\vpett{\nameref{int:InterfacesP1IHISAccess}}
		\item[Required interfaces:]~\iconrequired{}~\vpett{\nameref{int:InterfaceshealthAPI}}
		\item[Deployed on:]~\faSquareO~\vpett{\nameref{node:DeploymentPatientDataNode}}, \faSquareO~\vpett{\nameref{node:DeploymentPatientDataNodeDevDeployment}}, \faSquareO~\vpett{\nameref{node:DeploymentPatientDataNodePilotDeployment}}
		\item[Visible on diagrams:]~\cref{diag:Component:Primarydiagramoftheclientserverview,diag:Deployment:DevelopmentTestDeployment20patients3clinmodels,diag:Deployment:PilotDeployment200patients5clinmodels,diag:Deployment:Primarydeploymentdiagram}		
	\end{description}

\subsection{P1NotificationDispatcher}\label{comp:ComponentsP1NotificationDispatcher}
	\begin{description}[noitemsep,nolistsep]
		\item[Responsibility:]~UC5 dispatcher. Subscribes to P1IRiskEvents and fans matching events out to the registered physicians. \ref{data:DatatypesPriority} queue orders red over yellow over green. Drops green first when over 100 per minute. Persists red and yellow to a durable log so they survive a crash.
		\item[Super-components:]~None
		\item[Sub-components:]~None
		\item[Provided interfaces:]~\iconprovided{}~\vpett{\nameref{int:InterfacesP1INotificationInbox}}, \iconprovided{}~\vpett{\nameref{int:InterfacesP1IRiskEventListener}}
		\item[Required interfaces:]~\iconrequired{}~\vpett{\nameref{int:InterfacesAv2EmergencyDispatch}}, \iconrequired{}~\vpett{\nameref{int:InterfacesP1IRiskEvents}}
		\item[Deployed on:]~\faSquareO~\vpett{\nameref{node:DeploymenteHealthPlatform}}, \faSquareO~\vpett{\nameref{node:DeploymentDevBackendNodeDevDeployment}}, \faSquareO~\vpett{\nameref{node:DeploymentPhysicianAccessNode}}, \faSquareO~\vpett{\nameref{node:DeploymentPhysicianNodePilotDeployment}}
		\item[Visible on diagrams:]~\cref{diag:Component:DecompositionoftheP1NotificationDispatcher,diag:Component:Primarydiagramoftheclientserverview,diag:Deployment:Contextdiagramforthedeploymentview,diag:Deployment:DevelopmentTestDeployment20patients3clinmodels,diag:Deployment:PilotDeployment200patients5clinmodels,diag:Deployment:Primarydeploymentdiagram}		
	\end{description}

\subsection{P1PatientDataService}\label{comp:ComponentsP1PatientDataService}
	\begin{description}[noitemsep,nolistsep]
		\item[Responsibility:]~Owns raw sensor data, patient status, and the \vpett{\nameref{comp:ComponentsHIS}} EHR proxy. Takes over the three interfaces \vpett{\nameref{comp:ComponentseHealthPlatformOtherFunctionality}} exposed.
		\item[Super-components:]~None
		\item[Sub-components:]~None
		\item[Provided interfaces:]~\iconprovided{}~\vpett{\nameref{int:InterfacesOtherDataMgmt}}, \iconprovided{}~\vpett{\nameref{int:InterfacesP1IRiskEvents}}, \iconprovided{}~\vpett{\nameref{int:InterfacesPatientRecordMgmt}}, \iconprovided{}~\vpett{\nameref{int:InterfacesSensorDataMgmt}}
		\item[Required interfaces:]~\iconrequired{}~\vpett{\nameref{int:InterfacesLaunchRiskEstimation}}, \iconrequired{}~\vpett{\nameref{int:InterfacesP1IHISAccess}}, \iconrequired{}~\vpett{\nameref{int:InterfacesP1IRiskEventListener}}
		\item[Deployed on:]~\faSquareO~\vpett{\nameref{node:DeploymentPatientDataNode}}, \faSquareO~\vpett{\nameref{node:DeploymentPatientDataNodeDevDeployment}}, \faSquareO~\vpett{\nameref{node:DeploymentPatientDataNodePilotDeployment}}
		\item[Visible on diagrams:]~\cref{diag:Component:DecompositionoftheP1PatientDataService,diag:Component:Primarydiagramoftheclientserverview,diag:Deployment:DevelopmentTestDeployment20patients3clinmodels,diag:Deployment:PilotDeployment200patients5clinmodels,diag:Deployment:Primarydeploymentdiagram}		
	\end{description}

\subsection{P1PatientGatewayCommander}\label{comp:ComponentsP1PatientGatewayCommander}
	\begin{description}[noitemsep,nolistsep]
		\item[Responsibility:]~Outbound channel from PMS to patient gateways for UC9 on-demand sensor-data fetches. Synchronous with a bounded timeout.
		\item[Super-components:]~None
		\item[Sub-components:]~None
		\item[Provided interfaces:]~\iconprovided{}~\vpett{\nameref{int:InterfacesP1IOnDemandSensorFetch}}
		\item[Required interfaces:]~None
		\item[Deployed on:]~\faSquareO~\vpett{\nameref{node:DeploymentDevBackendNodeDevDeployment}}, \faSquareO~\vpett{\nameref{node:DeploymentPhysicianAccessNode}}, \faSquareO~\vpett{\nameref{node:DeploymentPhysicianNodePilotDeployment}}
		\item[Visible on diagrams:]~\cref{diag:Component:Primarydiagramoftheclientserverview,diag:Deployment:DevelopmentTestDeployment20patients3clinmodels,diag:Deployment:PilotDeployment200patients5clinmodels,diag:Deployment:Primarydeploymentdiagram}		
	\end{description}

\subsection{P1PatientQueryService}\label{comp:ComponentsP1PatientQueryService}
	\begin{description}[noitemsep,nolistsep]
		\item[Responsibility:]~Handles UC6 patient-status reads. Tiered priority queue (high, medium, low) feeds the 2 / 5 / 10 s SLA tiers. The tier for an incoming query comes from an in-memory patient-to-tier index kept current by subscribing to P1IRiskEvents, so the tier decision is O(1) on the hot path. Unknown patients default to the high tier so the SLA fails safe.
		\item[Super-components:]~None
		\item[Sub-components:]~None
		\item[Provided interfaces:]~\iconprovided{}~\vpett{\nameref{int:InterfacesP1IPatientQuery}}
		\item[Required interfaces:]~\iconrequired{}~\vpett{\nameref{int:InterfacesP1IPatientStatusRead}}, \iconrequired{}~\vpett{\nameref{int:InterfacesP1IRiskEvents}}
		\item[Deployed on:]~\faSquareO~\vpett{\nameref{node:DeploymenteHealthPlatform}}, \faSquareO~\vpett{\nameref{node:DeploymentDevBackendNodeDevDeployment}}, \faSquareO~\vpett{\nameref{node:DeploymentPhysicianAccessNode}}, \faSquareO~\vpett{\nameref{node:DeploymentPhysicianNodePilotDeployment}}
		\item[Visible on diagrams:]~\cref{diag:Component:Primarydiagramoftheclientserverview,diag:Deployment:Contextdiagramforthedeploymentview,diag:Deployment:DevelopmentTestDeployment20patients3clinmodels,diag:Deployment:PilotDeployment200patients5clinmodels,diag:Deployment:Primarydeploymentdiagram}		
	\end{description}

\subsection{P1PatientStatusCache}\label{comp:ComponentsP1PatientStatusCache}
	\begin{description}[noitemsep,nolistsep]
		\item[Responsibility:]~Read-through cache for patient status. Pins high-risk patients to keep status reads off the slow path.
		\item[Super-components:]~None
		\item[Sub-components:]~None
		\item[Provided interfaces:]~\iconprovided{}~\vpett{\nameref{int:InterfacesP1IPatientStatusRead}}
		\item[Required interfaces:]~\iconrequired{}~\vpett{\nameref{int:InterfacesOtherDataMgmt}}
		\item[Deployed on:]~\faSquareO~\vpett{\nameref{node:DeploymentPhysicianAccessNode}}, \faSquareO~\vpett{\nameref{node:DeploymentPatientDataNodeDevDeployment}}, \faSquareO~\vpett{\nameref{node:DeploymentPatientDataNodePilotDeployment}}
		\item[Visible on diagrams:]~\cref{diag:Component:Primarydiagramoftheclientserverview,diag:Deployment:DevelopmentTestDeployment20patients3clinmodels,diag:Deployment:PilotDeployment200patients5clinmodels,diag:Deployment:Primarydeploymentdiagram}		
	\end{description}

\subsection{P1PhysicianCommandService}\label{comp:ComponentsP1PhysicianCommandService}
	\begin{description}[noitemsep,nolistsep]
		\item[Responsibility:]~Handles physician write and command flows for UC7 (configure risk assessment), UC8 (update risk level), and UC9 (on-demand consultation).
		\item[Super-components:]~None
		\item[Sub-components:]~None
		\item[Provided interfaces:]~\iconprovided{}~\vpett{\nameref{int:InterfacesP1IPhysicianCommand}}
		\item[Required interfaces:]~\iconrequired{}~\vpett{\nameref{int:InterfacesClinicalModelCacheMgmt}}, \iconrequired{}~\vpett{\nameref{int:InterfacesLaunchRiskEstimation}}, \iconrequired{}~\vpett{\nameref{int:InterfacesP1ClinicalConfigMgmt}}, \iconrequired{}~\vpett{\nameref{int:InterfacesP1IOnDemandSensorFetch}}, \iconrequired{}~\vpett{\nameref{int:InterfacesP1IRiskEvents}}, \iconrequired{}~\vpett{\nameref{int:InterfacesPatientRecordMgmt}}, \iconrequired{}~\vpett{\nameref{int:InterfacesSensorDataMgmt}}
		\item[Deployed on:]~\faSquareO~\vpett{\nameref{node:DeploymenteHealthPlatform}}, \faSquareO~\vpett{\nameref{node:DeploymentDevBackendNodeDevDeployment}}, \faSquareO~\vpett{\nameref{node:DeploymentPhysicianAccessNode}}, \faSquareO~\vpett{\nameref{node:DeploymentPhysicianNodePilotDeployment}}
		\item[Visible on diagrams:]~\cref{diag:Component:DecompositionoftheP1PhysicianCommandService,diag:Component:Primarydiagramoftheclientserverview,diag:Deployment:Contextdiagramforthedeploymentview,diag:Deployment:DevelopmentTestDeployment20patients3clinmodels,diag:Deployment:PilotDeployment200patients5clinmodels,diag:Deployment:Primarydeploymentdiagram}		
	\end{description}

\subsection{P1PhysicianGateway}\label{comp:ComponentsP1PhysicianGateway}
	\begin{description}[noitemsep,nolistsep]
		\item[Responsibility:]~Single external entry point for physicians. Thin router with no business logic. Also dispatches UC9 result pushes over the physician's already-open client channel.
		\item[Super-components:]~None
		\item[Sub-components:]~None
		\item[Provided interfaces:]~\iconprovided{}~\vpett{\nameref{int:InterfacesP1IPhysicianAPI}}
		\item[Required interfaces:]~\iconrequired{}~\vpett{\nameref{int:InterfacesP1INotificationInbox}}, \iconrequired{}~\vpett{\nameref{int:InterfacesP1IPatientQuery}}, \iconrequired{}~\vpett{\nameref{int:InterfacesP1IPhysicianCommand}}
		\item[Deployed on:]~\faSquareO~\vpett{\nameref{node:DeploymentDevBackendNodeDevDeployment}}, \faSquareO~\vpett{\nameref{node:DeploymentPhysicianAccessNode}}, \faSquareO~\vpett{\nameref{node:DeploymentPhysicianWorkstation}}, \faSquareO~\vpett{\nameref{node:DeploymentPhysicianNodePilotDeployment}}
		\item[Visible on diagrams:]~\cref{diag:Component:Primarydiagramoftheclientserverview,diag:Deployment:Contextdiagramforthedeploymentview,diag:Deployment:DevelopmentTestDeployment20patients3clinmodels,diag:Deployment:PilotDeployment200patients5clinmodels,diag:Deployment:Primarydeploymentdiagram}		
	\end{description}

\subsection{PatientManagementService}\label{comp:ComponentsPatientManagementService}
	\begin{description}[noitemsep,nolistsep]
		\item[Responsibility:]~Provides administrative operations for managing patient records (enrolment, demographic updates, assigned clinical models, risk-level changes). It is the write-side service over the \vpett{\nameref{comp:ComponentsPatientRecordRepository}} and the source of truth for patient-related configuration consumed by other backend components.
		\item[Super-components:]~None
		\item[Sub-components:]~None
		\item[Provided interfaces:]~None
		\item[Required interfaces:]~None
		\item[Deployed on:]~\faSquareO~\vpett{\nameref{node:DeploymentDevPatientRecordsNodeDevDeployment}}, \faSquareO~\vpett{\nameref{node:DeploymentPatientRecordsNodePilotDeployment}}, \faSquareO~\vpett{\nameref{node:DeploymentPatientRecordsNode}}
		\item[Visible on diagrams:]~\cref{diag:Component:Primarydiagramoftheclientserverview,diag:Deployment:DevelopmentTestDeployment20patients3clinmodels,diag:Deployment:PilotDeployment200patients5clinmodels,diag:Deployment:Primarydeploymentdiagram}		
	\end{description}

\subsection{PatientRecordRepository}\label{comp:ComponentsPatientRecordRepository}
	\begin{description}[noitemsep,nolistsep]
		\item[Responsibility:]~Persistent store for patient demographic data, electronic health record references, assigned clinical models, current risk level, and authentication credentials. Provides the PatientRecordMgmt, PatientStatusOverview, and UserAuthentication interfaces. Acts as the authoritative source of patient identity and risk classification within the PMS backend.
		\item[Super-components:]~None
		\item[Sub-components:]~None
		\item[Provided interfaces:]~\iconprovided{}~\vpett{\nameref{int:InterfacesAv2PatientContactLookup}}, \iconprovided{}~\vpett{\nameref{int:InterfacesAv2PatientStatusOverview}}, \iconprovided{}~\vpett{\nameref{int:InterfacesUserAuthentication}}
		\item[Required interfaces:]~None
		\item[Deployed on:]~\faSquareO~\vpett{\nameref{node:DeploymenteHealthPlatform}}, \faSquareO~\vpett{\nameref{node:DeploymentDevPatientRecordsNodeDevDeployment}}, \faSquareO~\vpett{\nameref{node:DeploymentPatientRecordsNodePilotDeployment}}, \faSquareO~\vpett{\nameref{node:DeploymentPatientRecordsNode}}
		\item[Visible on diagrams:]~\cref{diag:Component:Primarydiagramoftheclientserverview,diag:Deployment:Contextdiagramforthedeploymentview,diag:Deployment:DevelopmentTestDeployment20patients3clinmodels,diag:Deployment:PilotDeployment200patients5clinmodels,diag:Deployment:Primarydeploymentdiagram}		
	\end{description}

\subsection{PatientStatusService}\label{comp:ComponentsPatientStatusService}
	\begin{description}[noitemsep,nolistsep]
		\item[Responsibility:]~Exposes read-only patient status information (current risk level, recent risk-assessment outcomes) to other services and user-facing components. Reads from the \vpett{\nameref{comp:ComponentsPatientRecordRepository}} via the PatientStatusOverview interface.
		\item[Super-components:]~None
		\item[Sub-components:]~None
		\item[Provided interfaces:]~None
		\item[Required interfaces:]~\iconrequired{}~\vpett{\nameref{int:InterfacesAv2PatientStatusOverview}}
		\item[Deployed on:]~\faSquareO~\vpett{\nameref{node:DeploymentPatientRecordsNode}}
		\item[Visible on diagrams:]~\cref{diag:Component:Primarydiagramoftheclientserverview,diag:Deployment:Primarydeploymentdiagram}		
	\end{description}

\subsection{RiskEstimationCombiner}\label{comp:ComponentseHealthPlatformRiskEstimationCombiner}
	\begin{description}[noitemsep,nolistsep]
		\item[Responsibility:]~The \vpett{\nameref{comp:ComponentseHealthPlatformRiskEstimationCombiner}} is responsible for combining the results of the \ref{data:DatatypesClinicalModelJob}s belonging to a patient risk estimation.
More precisely, the \vpett{\nameref{comp:ComponentseHealthPlatformRiskEstimationScheduler}} passes the set of the scheduled jobs for a risk estimation to the \vpett{\nameref{comp:ComponentseHealthPlatformRiskEstimationCombiner}} before scheduling them.
The \vpett{\nameref{comp:ComponentseHealthPlatformRiskEstimationCombiner}} then waits for all results to arrive, combines them, and propagates the new sensor data and the results of the risk estimation to the patient record if needed.
		\item[Super-components:]~\iconcomponent{}~\vpett{\nameref{comp:ComponentseHealthPlatform}}
		\item[Sub-components:]~None
		\item[Provided interfaces:]~\iconprovided{}~\vpett{\nameref{int:InterfacesJobMgmt}}, \iconprovided{}~\vpett{\nameref{int:InterfacesResults}}
		\item[Required interfaces:]~\iconrequired{}~\vpett{\nameref{int:InterfacesOtherDataMgmt}}, \iconrequired{}~\vpett{\nameref{int:InterfacesPatientRecordMgmt}}
		\item[Deployed on:]~\faSquareO~\vpett{\nameref{node:DeploymentRiskEstimationManagementNodePilotDeployment}}, \faSquareO~\vpett{\nameref{node:DeploymentRiskEstimationManagmentNodeDevDeployment}}, \faSquareO~\vpett{\nameref{node:DeploymenteHealthPlatformRiskEstimationMgmtNode}}
		\item[Visible on diagrams:]~\cref{diag:Component:Primarydiagramoftheclientserverview,diag:Deployment:DevelopmentTestDeployment20patients3clinmodels,diag:Deployment:PilotDeployment200patients5clinmodels,diag:Deployment:Primarydeploymentdiagram,diag:Interaction:CombiningClinicalModelJobresults,diag:Interaction:Incomingsensordata,diag:Interaction:Riskestimationprocess}		
	\end{description}

\subsection{RiskEstimationProcessor}\label{comp:ComponentseHealthPlatformRiskEstimationProcessor}
	\begin{description}[noitemsep,nolistsep]
		\item[Responsibility:]~The \vpett{\nameref{comp:ComponentseHealthPlatformRiskEstimationProcessor}} is responsible for computing \ref{data:DatatypesClinicalModelJob}s. The
\vpett{\nameref{comp:ComponentseHealthPlatformRiskEstimationProcessor}} fetches new \ref{data:DatatypesClinicalModelJob}s from the \vpett{\nameref{comp:ComponentseHealthPlatformRiskEstimationScheduler}}, uses the specified \ref{data:DatatypesMLModel} to calculate the risk prediction, and passes the result to the \vpett{\nameref{comp:ComponentseHealthPlatformRiskEstimationCombiner}}.
Multiple instances of the \vpett{\nameref{comp:ComponentseHealthPlatformRiskEstimationProcessor}} run in parallel to improve the throughput.
		\item[Super-components:]~\iconcomponent{}~\vpett{\nameref{comp:ComponentseHealthPlatform}}
		\item[Sub-components:]~None
		\item[Provided interfaces:]~None
		\item[Required interfaces:]~\iconrequired{}~\vpett{\nameref{int:InterfacesFetchJobs}}, \iconrequired{}~\vpett{\nameref{int:InterfacesKVStore}}, \iconrequired{}~\vpett{\nameref{int:InterfacesMLModelMgmt}}, \iconrequired{}~\vpett{\nameref{int:InterfacesResults}}
		\item[Deployed on:]~\faSquareO~\vpett{\nameref{node:DeploymentRiskEstimationProcessorNodePilotDeployment}}, \faSquareO~\vpett{\nameref{node:DeploymentRiskEstimationProcessorNodeDevDeployment}}, \faSquareO~\vpett{\nameref{node:DeploymenteHealthPlatformRiskEstimationProcessorNode}}
		\item[Visible on diagrams:]~\cref{diag:Component:Contextdiagramfortheclientserverview,diag:Component:DecompositionoftheRiskEstimationProcessor,diag:Component:Primarydiagramoftheclientserverview,diag:Deployment:DevelopmentTestDeployment20patients3clinmodels,diag:Deployment:PilotDeployment200patients5clinmodels,diag:Deployment:Primarydeploymentdiagram,diag:Interaction:CombiningClinicalModelJobresults}		
	\end{description}

\subsection{RiskEstimationScheduler}\label{comp:ComponentseHealthPlatformRiskEstimationScheduler}
	\begin{description}[noitemsep,nolistsep]
		\item[Responsibility:]~The \vpett{\nameref{comp:ComponentseHealthPlatformRiskEstimationScheduler}} is responsible for taking in new and scheduling requests for risk estimations.
It keeps track of the throughput of incoming jobs and changes the scheduling from first-in/first-out to dynamic priority earliest deadline first when going into overload modus.
		\item[Super-components:]~\iconcomponent{}~\vpett{\nameref{comp:ComponentseHealthPlatform}}
		\item[Sub-components:]~None
		\item[Provided interfaces:]~\iconprovided{}~\vpett{\nameref{int:InterfacesFetchJobs}}, \iconprovided{}~\vpett{\nameref{int:InterfacesLaunchRiskEstimation}}
		\item[Required interfaces:]~\iconrequired{}~\vpett{\nameref{int:InterfacesClinicalJobMgmt}}, \iconrequired{}~\vpett{\nameref{int:InterfacesFetchClinicalModels}}, \iconrequired{}~\vpett{\nameref{int:InterfacesJobMgmt}}, \iconrequired{}~\vpett{\nameref{int:InterfacesOtherDataMgmt}}, \iconrequired{}~\vpett{\nameref{int:InterfacesSensorDataMgmt}}
		\item[Deployed on:]~\faSquareO~\vpett{\nameref{node:DeploymentRiskEstimationManagementNodePilotDeployment}}, \faSquareO~\vpett{\nameref{node:DeploymentRiskEstimationManagmentNodeDevDeployment}}, \faSquareO~\vpett{\nameref{node:DeploymenteHealthPlatformRiskEstimationMgmtNode}}
		\item[Visible on diagrams:]~\cref{diag:Component:Primarydiagramoftheclientserverview,diag:Deployment:DevelopmentTestDeployment20patients3clinmodels,diag:Deployment:PilotDeployment200patients5clinmodels,diag:Deployment:Primarydeploymentdiagram,diag:Interaction:Incomingsensordata,diag:Interaction:Riskestimationprocess}		
	\end{description}

\subsection{UserAuthenticationService}\label{comp:ComponentsUserAuthenticationService}
	\begin{description}[noitemsep,nolistsep]
		\item[Responsibility:]~Authenticates human users (administrators, clinicians) accessing PMS interfaces. Verifies credentials against the \vpett{\nameref{comp:ComponentsPatientRecordRepository}} via the UserAuthentication interface. Distinct from the AuthenticationValidator inside the CommunicationGateway, which authenticates patient-gateway devices, not human users.
		\item[Super-components:]~None
		\item[Sub-components:]~None
		\item[Provided interfaces:]~None
		\item[Required interfaces:]~\iconrequired{}~\vpett{\nameref{int:InterfacesUserAuthentication}}
		\item[Deployed on:]~\faSquareO~\vpett{\nameref{node:DeploymentDevPatientRecordsNodeDevDeployment}}, \faSquareO~\vpett{\nameref{node:DeploymentPatientRecordsNodePilotDeployment}}, \faSquareO~\vpett{\nameref{node:DeploymentPatientRecordsNode}}
		\item[Visible on diagrams:]~\cref{diag:Component:Primarydiagramoftheclientserverview,diag:Deployment:DevelopmentTestDeployment20patients3clinmodels,diag:Deployment:PilotDeployment200patients5clinmodels,diag:Deployment:Primarydeploymentdiagram}		
	\end{description}


% END COMPONENTS

% MODULES
\section{Modules}\label{sec:modules}
\subsection{Av2AuthenticationValidator}\label{mod:ComponentsAv2CommunicationGatewayAv2AuthenticationValidator}
	\begin{description}[noitemsep,nolistsep]
		\item[Responsibility:]~Validates the authentication credentials presented by patient gateways before any data is processed further. It receives parsed transmissions from the RequestReceiver via InboundGatewayData, verifies that the gateway is registered and that its credentials are valid, and forwards authenticated transmissions to the RoutingManager via ValidatedTransmission. Transmissions that fail authentication are rejected, logged with the gateway identifier, IP address, and timestamp, and a \ref{data:DatatypesFailureEvent} is raised to the FailureDetector. Only transmissions from registered and authenticated gateways are allowed to proceed.
		\item[Super-components:]~\iconcomponent{}~\vpett{\nameref{comp:ComponentsAv2CommunicationGateway}}
		\item[Super-modules:]~None
		\item[Sub-modules:]~None
		\item[Provided interfaces:]~\iconprovided{}~\vpett{\nameref{int:InterfacesAv2AuthFailureSignals}}, \iconprovided{}~\vpett{\nameref{int:InterfacesAv2ValidatedTransmission}}
		\item[Required interfaces:]~\iconrequired{}~\vpett{\nameref{int:InterfacesAv2InboundGatewayData}}
		\item[Visible on diagrams:]~\cref{diag:Component:DecompositionoftheAv2CommunicationGateway,diag:Deployment:Contextdiagramforthedeploymentview}		
	\end{description}

\subsection{Av2BackupChannelManager}\label{mod:ComponentsAv2PatientGatewayAv2BackupChannelManager}
	\begin{description}[noitemsep,nolistsep]
		\item[Responsibility:]~Manages the backup communication channel, currently SMS, used for emergency notifications when the primary channel is unavailable. It monitors whether the backup channel is available and dispatches emergency messages through it when the DegradeModeManager reports that the gateway is in degraded mode.
		\item[Super-components:]~\iconcomponent{}~\vpett{\nameref{comp:ComponentsAv2PatientGateway}}
		\item[Super-modules:]~None
		\item[Sub-modules:]~None
		\item[Provided interfaces:]~\iconprovided{}~\vpett{\nameref{int:InterfacesAv2BackupChannelDispatch}}
		\item[Required interfaces:]~\iconrequired{}~\vpett{\nameref{int:InterfacesAv2ConnectivityStatus}}, \iconrequired{}~\vpett{\nameref{int:InterfacesAv2DegradeModeStatus}}
		\item[Visible on diagrams:]~\cref{diag:Component:DecompositionoftheAv2PatientGateway}		
	\end{description}

\subsection{Av2BatteryMonitor}\label{mod:ComponentsAv2PatientGatewayAv2BatteryMonitor}
	\begin{description}[noitemsep,nolistsep]
		\item[Responsibility:]~Continuously monitors the battery charge level of the patient gateway device. It publishes battery status events and triggers an alert to the PatientAlertService when the charge drops to or below 15\%, satisfying the Av2 prevention response measure.
		\item[Super-components:]~\iconcomponent{}~\vpett{\nameref{comp:ComponentsAv2PatientGateway}}
		\item[Super-modules:]~None
		\item[Sub-modules:]~None
		\item[Provided interfaces:]~\iconprovided{}~\vpett{\nameref{int:InterfacesAv2BatteryStatus}}
		\item[Required interfaces:]~None
		\item[Visible on diagrams:]~\cref{diag:Component:DecompositionoftheAv2PatientGateway}		
	\end{description}

\subsection{Av2CommunicationGapLog}\label{mod:ComponentsAv2MonitoringServiceAv2CommunicationGapLog}
	\begin{description}[noitemsep,nolistsep]
		\item[Responsibility:]~Persistently records the start, end, and duration of every
communication interruption per gateway. Provides historical query access
for administrators and compliance. Updated by RecoveryTracker when a
failure is opened/closed.
		\item[Super-components:]~\iconcomponent{}~\vpett{\nameref{comp:ComponentsAv2MonitoringService}}
		\item[Super-modules:]~None
		\item[Sub-modules:]~None
		\item[Provided interfaces:]~\iconprovided{}~\vpett{\nameref{int:InterfacesAv2CommunicationGapHistory}}
		\item[Required interfaces:]~None
		\item[Visible on diagrams:]~\cref{diag:Component:DecompositionoftheAv2MonitoringService}		
	\end{description}

\subsection{Av2DegradeModeManager}\label{mod:ComponentsAv2PatientGatewayAv2DegradeModeManager}
	\begin{description}[noitemsep,nolistsep]
		\item[Responsibility:]~Manages the degraded mode state of the patient gateway. It transitions the gateway into degraded mode when the ServerConnectivityMonitor reports the backend as unreachable, and back to normal mode upon successful reconnection. It records the timestamp at which degraded mode was entered, enabling gap duration tracking. All other modules on the gateway consult DegradeModeStatus to adapt their behavior accordingly.
		\item[Super-components:]~\iconcomponent{}~\vpett{\nameref{comp:ComponentsAv2PatientGateway}}
		\item[Super-modules:]~None
		\item[Sub-modules:]~None
		\item[Provided interfaces:]~\iconprovided{}~\vpett{\nameref{int:InterfacesAv2DegradeModeStatus}}
		\item[Required interfaces:]~\iconrequired{}~\vpett{\nameref{int:InterfacesAv2ServerConnectionStatus}}
		\item[Visible on diagrams:]~\cref{diag:Component:DecompositionoftheAv2PatientGateway}		
	\end{description}

\subsection{Av2EmergencyNotificationService}\label{mod:ComponentsAv2PatientGatewayAv2EmergencyNotificationService}
	\begin{description}[noitemsep,nolistsep]
		\item[Responsibility:]~Handles the dispatch of emergency notifications from the patient gateway. In normal operation it uses the primary communication channel. When the gateway is in degraded mode it delegates to the BackupChannelManager to send the emergency notification via the backup channel, ensuring no emergency alerts are lost due to primary channel failure.
		\item[Super-components:]~\iconcomponent{}~\vpett{\nameref{comp:ComponentsAv2PatientGateway}}
		\item[Super-modules:]~None
		\item[Sub-modules:]~None
		\item[Provided interfaces:]~\iconprovided{}~\vpett{\nameref{int:InterfacesAv2EmergencyDispatch}}
		\item[Required interfaces:]~\iconrequired{}~\vpett{\nameref{int:InterfacesAv2BackupChannelDispatch}}, \iconrequired{}~\vpett{\nameref{int:InterfacesAv2ConnectivityStatus}}, \iconrequired{}~\vpett{\nameref{int:InterfacesAv2DegradeModeStatus}}
		\item[Visible on diagrams:]~\cref{diag:Component:DecompositionoftheAv2PatientGateway}		
	\end{description}

\subsection{Av2FailureClassifier}\label{mod:ComponentsAv2MonitoringServiceAv2FailureClassifier}
	\begin{description}[noitemsep,nolistsep]
		\item[Responsibility:]~Receives \ref{data:DatatypesFailureEvent} objects from both the MissingDataDetector and the InternalFailureDetector via the DetectedFailures interface. For each failure it looks up the risk level of the affected patient via PatientStatusOverview, delegates deadline computation to the NotificationTimingPolicy, and constructs a \ref{data:DatatypesClassifiedFailure} enriched with the escalation level and notification deadline. It emits \ref{data:DatatypesClassifiedFailure} objects via EscalationEvents to the StakeholderNotifier. When a failure is resolved it signals the RecoveryTracker via RecoveryStatus.
		\item[Super-components:]~\iconcomponent{}~\vpett{\nameref{comp:ComponentsAv2MonitoringService}}
		\item[Super-modules:]~None
		\item[Sub-modules:]~None
		\item[Provided interfaces:]~\iconprovided{}~\vpett{\nameref{int:InterfacesAv2EscalationEvents}}, \iconprovided{}~\vpett{\nameref{int:InterfacesAv2RecoveryStatus}}
		\item[Required interfaces:]~\iconrequired{}~\vpett{\nameref{int:InterfacesAv2DetectedFailures}}, \iconrequired{}~\vpett{\nameref{int:InterfacesAv2PatientStatusOverview}}
		\item[Visible on diagrams:]~\cref{diag:Component:DecompositionoftheAv2MonitoringService,diag:Deployment:Contextdiagramforthedeploymentview}		
	\end{description}

\subsection{Av2FailureDetector}\label{mod:ComponentsAv2CommunicationGatewayAv2FailureDetector}
	\begin{description}[noitemsep,nolistsep]
		\item[Responsibility:]~Monitors the health of the CommunicationGateway's own subcomponents and detects connection-level failures. It responds to the InternalHealthCheck interface by reporting the current operational status of the CommunicationGateway and its subcomponents. It also receives authentication failure signals from the AuthenticationValidator and connection timeout signals, constructing \ref{data:DatatypesFailureEvent} objects and publishing them via the FailureEvents interface to the RetryCoordinator. Internal failures must be detectable within five seconds as required by Av2.
		\item[Super-components:]~\iconcomponent{}~\vpett{\nameref{comp:ComponentsAv2CommunicationGateway}}
		\item[Super-modules:]~None
		\item[Sub-modules:]~None
		\item[Provided interfaces:]~\iconprovided{}~\vpett{\nameref{int:InterfacesAv2CommunicationGatewayHealthCheck}}, \iconprovided{}~\vpett{\nameref{int:InterfacesAv2FailureEvents}}
		\item[Required interfaces:]~\iconrequired{}~\vpett{\nameref{int:InterfacesAv2AuthFailureSignals}}, \iconrequired{}~\vpett{\nameref{int:InterfacesAv2RoutFailureSignals}}
		\item[Visible on diagrams:]~\cref{diag:Component:DecompositionoftheAv2CommunicationGateway,diag:Deployment:Contextdiagramforthedeploymentview}		
	\end{description}

\subsection{Av2GatewayRegistry}\label{mod:ComponentsAv2CommunicationGatewayAv2GatewayRegistry}
	\begin{description}[noitemsep,nolistsep]
		\item[Responsibility:]~Maintains the authoritative bidirectional mapping between \ref{data:DatatypesGatewayId}, \ref{data:DatatypesPatientId}, expected transmission interval (in seconds), and current patient risk level. It is populated during patient registration and updated whenever a patient's risk level changes via OtherDataMgmt. It is consulted by HeartbeatTracker to calibrate per-gateway expected intervals whenever risk level changes, by AuthenticationValidator to verify that an incoming gateway is registered, and by FailureClassifier to resolve a \ref{data:DatatypesGatewayId} to the associated \ref{data:DatatypesPatientId} and risk level before computing escalation deadlines. It is the only component inside CommunicationGateway that holds patient-to-gateway identity mappings.
		\item[Super-components:]~\iconcomponent{}~\vpett{\nameref{comp:ComponentsAv2CommunicationGateway}}
		\item[Super-modules:]~None
		\item[Sub-modules:]~None
		\item[Provided interfaces:]~None
		\item[Required interfaces:]~\iconrequired{}~\vpett{\nameref{int:InterfacesAv2PatientStatusOverview}}
		\item[Visible on diagrams:]~\cref{diag:Component:DecompositionoftheAv2CommunicationGateway}		
	\end{description}

\subsection{Av2HeartbeatTracker}\label{mod:ComponentsAv2CommunicationGatewayAv2HeartbeatTracker}
	\begin{description}[noitemsep,nolistsep]
		\item[Responsibility:]~Maintains a per-gateway registry of the last time a transmission was received from each patient gateway and the configured expected transmission interval for that gateway. It receives heartbeat updates from the RequestReceiver via the HeartbeatUpdate interface each time a transmission arrives. It exposes the HeartbeatRegistry interface so the MissingDataDetector in the MonitoringService can poll for overdue gateways. It also supports gateway registration so that newly enrolled patients have their expected transmission interval configured from the start.
		\item[Super-components:]~\iconcomponent{}~\vpett{\nameref{comp:ComponentsAv2CommunicationGateway}}
		\item[Super-modules:]~None
		\item[Sub-modules:]~None
		\item[Provided interfaces:]~\iconprovided{}~\vpett{\nameref{int:InterfacesAv2AvailabilityMonitoring}}
		\item[Required interfaces:]~\iconrequired{}~\vpett{\nameref{int:InterfacesAv2HeartbeatUpdate}}
		\item[Visible on diagrams:]~\cref{diag:Component:DecompositionoftheAv2CommunicationGateway,diag:Deployment:Contextdiagramforthedeploymentview}		
	\end{description}

\subsection{Av2InternalFailureDetector}\label{mod:ComponentsAv2MonitoringServiceAv2InternalFailureDetector}
	\begin{description}[noitemsep,nolistsep]
		\item[Responsibility:]~Polls the health check endpoints of internal communication subsystems — specifically the CommunicationGateway and the MessageBroker — via the InternalHealthCheck interface at a frequency of at least once per second. When a subsystem returns a DEGRADED or DOWN status, or fails to respond within one second, it constructs a \ref{data:DatatypesFailureEvent} of type INTERNAL\_SUBSYSTEM\_CRASH or COMMUNICATION\_CHANNEL\_DOWN and reports it via the DetectedFailures interface to the FailureClassifier. This polling frequency guarantees detection within five seconds as required by Av2.
		\item[Super-components:]~\iconcomponent{}~\vpett{\nameref{comp:ComponentsAv2MonitoringService}}
		\item[Super-modules:]~None
		\item[Sub-modules:]~None
		\item[Provided interfaces:]~\iconprovided{}~\vpett{\nameref{int:InterfacesAv2DetectedFailures}}
		\item[Required interfaces:]~\iconrequired{}~\vpett{\nameref{int:InterfacesAv2CommunicationGatewayHealthCheck}}, \iconrequired{}~\vpett{\nameref{int:InterfacesAv2MessageBrokerHealthCheck}}
		\item[Visible on diagrams:]~\cref{diag:Component:DecompositionoftheAv2MonitoringService}		
	\end{description}

\subsection{Av2LocalBufferingRepository}\label{mod:ComponentsAv2PatientGatewayAv2LocalBufferingRepository}
	\begin{description}[noitemsep,nolistsep]
		\item[Responsibility:]~Persistently stores sensor data packages and their timestamps on the patient gateway during periods when the primary communication channel is unavailable. It guarantees storage capacity for at least 24 hours of sensor data at the configured transmission rate, satisfying the Av2 resolution response measure. It exposes operations for writing buffered readings and reading them back for delivery upon reconnection.
		\item[Super-components:]~\iconcomponent{}~\vpett{\nameref{comp:ComponentsAv2PatientGateway}}
		\item[Super-modules:]~None
		\item[Sub-modules:]~None
		\item[Provided interfaces:]~\iconprovided{}~\vpett{\nameref{int:InterfacesAv2BufferedSensorData}}
		\item[Required interfaces:]~None
		\item[Visible on diagrams:]~\cref{diag:Component:DecompositionoftheAv2PatientGateway}		
	\end{description}

\subsection{Av2MissingDataDetector}\label{mod:ComponentsAv2MonitoringServiceAv2MissingDataDetector}
	\begin{description}[noitemsep,nolistsep]
		\item[Responsibility:]~Periodically polls the HeartbeatRegistry data exposed via the AvailabilityMonitoring interface to identify patient gateways that have not sent an update within their expected transmission interval plus the five-minute tolerance defined by Av2. When an overdue gateway is detected it constructs a \ref{data:DatatypesFailureEvent} of type MISSING\_SENSOR\_DATA and reports it via the DetectedFailures interface to the FailureClassifier. The polling frequency must be high enough to guarantee detection within transmission interval plus five minutes.
		\item[Super-components:]~\iconcomponent{}~\vpett{\nameref{comp:ComponentsAv2MonitoringService}}
		\item[Super-modules:]~None
		\item[Sub-modules:]~None
		\item[Provided interfaces:]~\iconprovided{}~\vpett{\nameref{int:InterfacesAv2DetectedFailures}}
		\item[Required interfaces:]~\iconrequired{}~\vpett{\nameref{int:InterfacesAv2AvailabilityMonitoring}}
		\item[Visible on diagrams:]~\cref{diag:Component:DecompositionoftheAv2MonitoringService}		
	\end{description}

\subsection{Av2NetworkConnectivityMonitor}\label{mod:ComponentsAv2PatientGatewayAv2NetworkConnectivityMonitor}
	\begin{description}[noitemsep,nolistsep]
		\item[Responsibility:]~Continuously monitors the availability and quality of the network connection on the patient gateway. It publishes connectivity status events indicating signal strength and network type, enabling the PatientAlertService to warn the patient of poor or absent connectivity.
		\item[Super-components:]~\iconcomponent{}~\vpett{\nameref{comp:ComponentsAv2PatientGateway}}
		\item[Super-modules:]~None
		\item[Sub-modules:]~None
		\item[Provided interfaces:]~\iconprovided{}~\vpett{\nameref{int:InterfacesAv2ConnectivityStatus}}
		\item[Required interfaces:]~None
		\item[Visible on diagrams:]~\cref{diag:Component:DecompositionoftheAv2PatientGateway}		
	\end{description}

\subsection{Av2PatientAlertService}\label{mod:ComponentsAv2PatientGatewayAv2PatientAlertService}
	\begin{description}[noitemsep,nolistsep]
		\item[Responsibility:]~Presents warnings to the patient on the gateway device. It alerts the patient when the battery drops to or below 15\% and when network connectivity is absent or insufficient. These warnings fulfill the Av2 prevention response measures for battery and connectivity.
		\item[Super-components:]~\iconcomponent{}~\vpett{\nameref{comp:ComponentsAv2PatientGateway}}
		\item[Super-modules:]~None
		\item[Sub-modules:]~None
		\item[Provided interfaces:]~None
		\item[Required interfaces:]~\iconrequired{}~\vpett{\nameref{int:InterfacesAv2BatteryStatus}}, \iconrequired{}~\vpett{\nameref{int:InterfacesAv2ConnectivityStatus}}
		\item[Visible on diagrams:]~\cref{diag:Component:DecompositionoftheAv2PatientGateway}		
	\end{description}

\subsection{Av2RecoveryTracker}\label{mod:ComponentsAv2MonitoringServiceAv2RecoveryTracker}
	\begin{description}[noitemsep,nolistsep]
		\item[Responsibility:]~Tracks the lifecycle of active failures from detection through resolution. It receives \ref{data:DatatypesClassifiedFailure} signals via EscalationEvents when a failure is detected and RecoveryStatus signals from the FailureClassifier when a failure is resolved. It records the failure start time, detection time, and resolution time, computes the total communication gap duration, and signals the RecoverySyncService and RedeploymentCoordinator via the RecoveryCoordination interface to begin recovery operations. It also records how long communication was interrupted per the Av2 detection requirement.
		\item[Super-components:]~\iconcomponent{}~\vpett{\nameref{comp:ComponentsAv2MonitoringService}}
		\item[Super-modules:]~None
		\item[Sub-modules:]~None
		\item[Provided interfaces:]~\iconprovided{}~\vpett{\nameref{int:InterfacesAv2RecoveryCoordination}}
		\item[Required interfaces:]~\iconrequired{}~\vpett{\nameref{int:InterfacesAv2CommunicationGapHistory}}, \iconrequired{}~\vpett{\nameref{int:InterfacesAv2EscalationEvents}}, \iconrequired{}~\vpett{\nameref{int:InterfacesAv2RecoveryStatus}}
		\item[Visible on diagrams:]~\cref{diag:Component:DecompositionoftheAv2MonitoringService}		
	\end{description}

\subsection{Av2RequestReceiver}\label{mod:ComponentsAv2CommunicationGatewayAv2RequestReceiver}
	\begin{description}[noitemsep,nolistsep]
		\item[Responsibility:]~Receives raw inbound transmissions from patient gateways via the GatewayDataIngress interface, performs initial parsing and format validation, and forwards validated inbound data to the AuthenticationValidator via the InboundGatewayData interface. It also extracts the gateway identifier and receipt timestamp from each transmission and records them via the HeartbeatUpdate interface to the HeartbeatTracker, ensuring that every received transmission is counted as a heartbeat regardless of whether it passes authentication.
		\item[Super-components:]~\iconcomponent{}~\vpett{\nameref{comp:ComponentsAv2CommunicationGateway}}
		\item[Super-modules:]~None
		\item[Sub-modules:]~None
		\item[Provided interfaces:]~\iconprovided{}~\vpett{\nameref{int:InterfacesAv2HeartbeatUpdate}}, \iconprovided{}~\vpett{\nameref{int:InterfacesAv2InboundGatewayData}}
		\item[Required interfaces:]~\iconrequired{}~\vpett{\nameref{int:InterfacesAv2GatewayDataIngress}}
		\item[Visible on diagrams:]~\cref{diag:Component:DecompositionoftheAv2CommunicationGateway,diag:Deployment:Contextdiagramforthedeploymentview}		
	\end{description}

\subsection{Av2RetryAndSynchronizationService}\label{mod:ComponentsAv2PatientGatewayAv2RetryAndSynchronizationService}
	\begin{description}[noitemsep,nolistsep]
		\item[Responsibility:]~Implements the exponential backoff retry strategy for delivering buffered sensor data to the PMS backend after a communication failure. It reads pending data from the LocalBufferingRepository and retries delivery with increasing intervals until the backend confirms receipt, at which point it clears the delivered entries from the buffer.
		\item[Super-components:]~\iconcomponent{}~\vpett{\nameref{comp:ComponentsAv2PatientGateway}}
		\item[Super-modules:]~None
		\item[Sub-modules:]~None
		\item[Provided interfaces:]~None
		\item[Required interfaces:]~\iconrequired{}~\vpett{\nameref{int:InterfacesAv2BufferedSensorData}}, \iconrequired{}~\vpett{\nameref{int:InterfacesAv2DegradeModeStatus}}, \iconrequired{}~\vpett{\nameref{int:InterfacesAv2DeliveryAck}}, \iconrequired{}~\vpett{\nameref{int:InterfacesAv2GatewayResyncCommand}}, \iconrequired{}~\vpett{\nameref{int:InterfacesAv2ServerConnectionStatus}}
		\item[Visible on diagrams:]~\cref{diag:Component:DecompositionoftheAv2PatientGateway}		
	\end{description}

\subsection{Av2RetryCoordinator}\label{mod:ComponentsAv2CommunicationGatewayAv2RetryCoordinator}
	\begin{description}[noitemsep,nolistsep]
		\item[Responsibility:]~Coordinates retry attempts for messages that the RoutingManager failed to deliver to downstream backend services. It receives \ref{data:DatatypesFailureEvent} objects from the FailureDetector via FailureEvents, determines whether a retry is appropriate, and schedules retries using exponential backoff. It publishes retry-ready messages via RetryDispatch to the RoutingManager. It cancels retries once delivery is confirmed or the failure has persisted long enough to require escalation to the MonitoringService.
		\item[Super-components:]~\iconcomponent{}~\vpett{\nameref{comp:ComponentsAv2CommunicationGateway}}
		\item[Super-modules:]~None
		\item[Sub-modules:]~None
		\item[Provided interfaces:]~\iconprovided{}~\vpett{\nameref{int:InterfacesAv2RetryDispatch}}
		\item[Required interfaces:]~\iconrequired{}~\vpett{\nameref{int:InterfacesAv2FailureEvents}}
		\item[Visible on diagrams:]~\cref{diag:Component:DecompositionoftheAv2CommunicationGateway,diag:Deployment:Contextdiagramforthedeploymentview}		
	\end{description}

\subsection{Av2RoutingManager}\label{mod:ComponentsAv2CommunicationGatewayAv2RoutingManager}
	\begin{description}[noitemsep,nolistsep]
		\item[Responsibility:]~Routes authenticated and validated gateway transmissions to the appropriate downstream service. Sensor data packages are forwarded to the MessageBroker via BufferedDataDispatch. Emergency notifications are forwarded to the EmergencyNotificationService. It also accepts retry-ready messages from the RetryCoordinator via RetryDispatch and re-attempts delivery of previously failed transmissions. It is the only module inside the CommunicationGateway that communicates with downstream backend services.
		\item[Super-components:]~\iconcomponent{}~\vpett{\nameref{comp:ComponentsAv2CommunicationGateway}}
		\item[Super-modules:]~None
		\item[Sub-modules:]~None
		\item[Provided interfaces:]~\iconprovided{}~\vpett{\nameref{int:InterfacesAv2BufferedDataDispatch}}, \iconprovided{}~\vpett{\nameref{int:InterfacesAv2RoutFailureSignals}}
		\item[Required interfaces:]~\iconrequired{}~\vpett{\nameref{int:InterfacesAv2RetryDispatch}}, \iconrequired{}~\vpett{\nameref{int:InterfacesAv2ValidatedTransmission}}
		\item[Visible on diagrams:]~\cref{diag:Component:DecompositionoftheAv2CommunicationGateway,diag:Deployment:Contextdiagramforthedeploymentview}		
	\end{description}

\subsection{Av2SensorCommunicator}\label{mod:ComponentsAv2PatientGatewayAv2SensorCommunicator}
	\begin{description}[noitemsep,nolistsep]
		\item[Responsibility:]~Acquires raw sensor readings from the wearable sensors attached to the patient gateway via the SensorDataAcquisition interface and places them into the DataBuffer for processing and transmission.
		\item[Super-components:]~\iconcomponent{}~\vpett{\nameref{comp:ComponentsAv2PatientGateway}}
		\item[Super-modules:]~None
		\item[Sub-modules:]~None
		\item[Provided interfaces:]~\iconprovided{}~\vpett{\nameref{int:InterfacesDataBuffer}}
		\item[Required interfaces:]~\iconrequired{}~\vpett{\nameref{int:InterfacesAv2SensorDataAcquisition}}
		\item[Visible on diagrams:]~\cref{diag:Component:DecompositionoftheAv2PatientGateway}		
	\end{description}

\subsection{Av2SensorDataTransmitter}\label{mod:ComponentsAv2PatientGatewayAv2SensorDataTransmitter}
	\begin{description}[noitemsep,nolistsep]
		\item[Responsibility:]~Reads sensor data from the DataBuffer provided by the SensorCommunicator and transmits it to the CommunicationGateway. In normal mode it transmits immediately. In degraded mode it hands data to the LocalBufferingRepository instead, and upon recovery delegates flushing to the RetryAndSynchronizationService.
		\item[Super-components:]~\iconcomponent{}~\vpett{\nameref{comp:ComponentsAv2PatientGateway}}
		\item[Super-modules:]~None
		\item[Sub-modules:]~None
		\item[Provided interfaces:]~\iconprovided{}~\vpett{\nameref{int:InterfacesAv2GatewayDataIngress}}
		\item[Required interfaces:]~\iconrequired{}~\vpett{\nameref{int:InterfacesAv2BufferedSensorData}}, \iconrequired{}~\vpett{\nameref{int:InterfacesAv2DegradeModeStatus}}, \iconrequired{}~\vpett{\nameref{int:InterfacesDataBuffer}}
		\item[Visible on diagrams:]~\cref{diag:Component:DecompositionoftheAv2PatientGateway}		
	\end{description}

\subsection{Av2ServerConnectivityMonitor}\label{mod:ComponentsAv2PatientGatewayAv2ServerConnectivityMonitor}
	\begin{description}[noitemsep,nolistsep]
		\item[Responsibility:]~Monitors the reachability of the PMS backend from the patient gateway by periodically sending heartbeat probes. It publishes server connection status events that drive the DegradeModeManager's decision to enter or exit degraded mode.
		\item[Super-components:]~\iconcomponent{}~\vpett{\nameref{comp:ComponentsAv2PatientGateway}}
		\item[Super-modules:]~None
		\item[Sub-modules:]~None
		\item[Provided interfaces:]~\iconprovided{}~\vpett{\nameref{int:InterfacesAv2ServerConnectionStatus}}
		\item[Required interfaces:]~None
		\item[Visible on diagrams:]~\cref{diag:Component:DecompositionoftheAv2PatientGateway}		
	\end{description}

\subsection{Av2SLAMonitor}\label{mod:ComponentsAv2MonitoringServiceAv2SLAMonitor}
	\begin{description}[noitemsep,nolistsep]
		\item[Responsibility:]~Aggregates per-gateway connectivity events and computes
rolling availability metrics (per year), bandwidth samples, and coverage
estimates. Compares against the SLA thresholds (99.5\% availability,
80\% coverage, 64Kb/s bandwidth) and raises a \ref{data:DatatypesFailureEvent} of type
SLA\_VIOLATION when thresholds are breached over the measurement window.
		\item[Super-components:]~\iconcomponent{}~\vpett{\nameref{comp:ComponentsAv2MonitoringService}}
		\item[Super-modules:]~None
		\item[Sub-modules:]~None
		\item[Provided interfaces:]~\iconprovided{}~\vpett{\nameref{int:InterfacesAv2DetectedFailures}}, \iconprovided{}~\vpett{\nameref{int:InterfacesAv2SLAReporting}}
		\item[Required interfaces:]~\iconrequired{}~\vpett{\nameref{int:InterfacesAv2AvailabilityMonitoring}}
		\item[Visible on diagrams:]~\cref{diag:Component:DecompositionoftheAv2MonitoringService,diag:Deployment:Contextdiagramforthedeploymentview}		
	\end{description}

\subsection{Av2StakeholderNotifier}\label{mod:ComponentsAv2MonitoringServiceAv2StakeholderNotifier}
	\begin{description}[noitemsep,nolistsep]
		\item[Responsibility:]~Receives \ref{data:DatatypesClassifiedFailure} objects via EscalationEvents and dispatches notifications to the appropriate stakeholders before the deadline encoded in each \ref{data:DatatypesClassifiedFailure}. For INTERNAL\_SUBSYSTEM\_CRASH and COMMUNICATION\_CHANNEL\_DOWN failures it notifies the system administrator. For GATEWAY\_UNREACHABLE and MISSING\_SENSOR\_DATA failures it notifies both the system administrator and initiates patient contact. It uses the NotificationDispatch interface to send the actual notifications. It must meet the notification deadline in 90\% of cases as required by Av2.
		\item[Super-components:]~\iconcomponent{}~\vpett{\nameref{comp:ComponentsAv2MonitoringService}}
		\item[Super-modules:]~None
		\item[Sub-modules:]~None
		\item[Provided interfaces:]~\iconprovided{}~\vpett{\nameref{int:InterfacesAv2NotificationDispatch}}
		\item[Required interfaces:]~\iconrequired{}~\vpett{\nameref{int:InterfacesAv2EscalationEvents}}, \iconrequired{}~\vpett{\nameref{int:InterfacesAv2PatientContactLookup}}
		\item[Visible on diagrams:]~\cref{diag:Component:DecompositionoftheAv2MonitoringService,diag:Deployment:Contextdiagramforthedeploymentview}		
	\end{description}

\subsection{JobMLModelSelector}\label{mod:ComponentseHealthPlatformRiskEstimationProcessorJobMLModelSelector}
	\begin{description}[noitemsep,nolistsep]
		\item[Responsibility:]~This module is responsible for determining the appropriate \ref{data:DatatypesMLModel} to use for making the prediction on the provided sensor data, as there could be multiple \ref{data:DatatypesMLModel} candidates for making predictions.
		\item[Super-components:]~\iconcomponent{}~\vpett{\nameref{comp:ComponentseHealthPlatformRiskEstimationProcessor}} $\triangleright$ \iconcomponent{}~\vpett{\nameref{comp:ComponentseHealthPlatform}}
		\item[Super-modules:]~None
		\item[Sub-modules:]~None
		\item[Provided interfaces:]~\iconprovided{}~\vpett{\nameref{int:InterfacesApplyClinicalModel}}
		\item[Required interfaces:]~None
		\item[Visible on diagrams:]~\cref{diag:Component:DecompositionoftheRiskEstimationProcessor,diag:Interaction:Computeclinicalmodelresult}		
	\end{description}

\subsection{JobProcessor}\label{mod:ComponentseHealthPlatformRiskEstimationProcessorJobProcessor}
	\begin{description}[noitemsep,nolistsep]
		\item[Responsibility:]~The \vpett{\nameref{mod:ComponentseHealthPlatformRiskEstimationProcessorJobProcessor}} module coordinates the execution of the \ref{data:DatatypesClinicalModelJob}s to ensure that the required \ref{data:DatatypesMLModel}s are fetched, and the prediction is performed with the appropriate \ref{data:DatatypesMLModel}.
		\item[Super-components:]~\iconcomponent{}~\vpett{\nameref{comp:ComponentseHealthPlatformRiskEstimationProcessor}} $\triangleright$ \iconcomponent{}~\vpett{\nameref{comp:ComponentseHealthPlatform}}
		\item[Super-modules:]~None
		\item[Sub-modules:]~None
		\item[Provided interfaces:]~None
		\item[Required interfaces:]~\iconrequired{}~\vpett{\nameref{int:InterfacesApplyClinicalModel}}, \iconrequired{}~\vpett{\nameref{int:InterfacesFetchJobs}}, \iconrequired{}~\vpett{\nameref{int:InterfacesFetchMLModel}}, \iconrequired{}~\vpett{\nameref{int:InterfacesResults}}, \iconrequired{}~\vpett{\nameref{int:InterfacesSensorDataAnalysis}}
		\item[Visible on diagrams:]~\cref{diag:Component:DecompositionoftheRiskEstimationProcessor,diag:Interaction:Computeclinicalmodelresult}		
	\end{description}

\subsection{MLaaSLibrary}\label{mod:ModulesMLaaSLibrary}
	\begin{description}[noitemsep,nolistsep]
		\item[Responsibility:]~\vpett{\nameref{comp:ComponentsMLaaS}} development library for working with the \vpett{\nameref{comp:ComponentsMLaaS}} providers.
It supports storage \& retrieval of \ref{data:DatatypesMLModel}s, and exchanging \ref{data:DatatypesMLModelUpdate}s with \vpett{\nameref{comp:ComponentsMLaaS}} providers.
		\item[Super-components:]~None
		\item[Super-modules:]~None
		\item[Sub-modules:]~None
		\item[Provided interfaces:]~\iconprovided{}~\vpett{\nameref{int:InterfacesMLaaSAPI}}
		\item[Required interfaces:]~None
		\item[Visible on diagrams:]~\cref{diag:Component:Contextdiagramforthedecompositionview}		
	\end{description}

\subsection{MLLibrary}\label{mod:ModulesMLLibrary}
	\begin{description}[noitemsep,nolistsep]
		\item[Responsibility:]~Library for working with \ref{data:DatatypesMLModel}s locally and using them to make predictions based on the incoming sensor data.
The library can also process external updates to apply them to a specific \ref{data:DatatypesMLModel}.
		\item[Super-components:]~None
		\item[Super-modules:]~None
		\item[Sub-modules:]~None
		\item[Provided interfaces:]~\iconprovided{}~\vpett{\nameref{int:InterfacesMLAPI}}
		\item[Required interfaces:]~None
		\item[Visible on diagrams:]~\cref{diag:Component:Contextdiagramforthedecompositionview}		
	\end{description}

\subsection{MLModelRepository}\label{mod:ComponentseHealthPlatformMLModelManagerMLModelRepository}
	\begin{description}[noitemsep,nolistsep]
		\item[Responsibility:]~This module stores \ref{data:DatatypesMLModel}s locally and makes them available for use by the \vpett{\nameref{comp:ComponentseHealthPlatformRiskEstimationProcessor}}  when processing \ref{data:DatatypesClinicalModelJob}s.
		\item[Super-components:]~\iconcomponent{}~\vpett{\nameref{comp:ComponentseHealthPlatform}} $\triangleright$ \iconcomponent{}~\vpett{\nameref{comp:ComponentseHealthPlatformMLModelManager}}
		\item[Super-modules:]~None
		\item[Sub-modules:]~None
		\item[Provided interfaces:]~\iconprovided{}~\vpett{\nameref{int:InterfacesMLModelStorage}}
		\item[Required interfaces:]~None
		\item[Visible on diagrams:]~\cref{diag:Component:DecompositionoftheMLModelManager,diag:Interaction:MLModelSynchronization,diag:Interaction:ProcessingMLModelupdates}		
	\end{description}

\subsection{MLModelRetrieval}\label{mod:ComponentseHealthPlatformRiskEstimationProcessorMLModelRetrieval}
	\begin{description}[noitemsep,nolistsep]
		\item[Responsibility:]~This module is responsible for retrieving the requested \ref{data:DatatypesMLModel}s from the \vpett{\nameref{comp:ComponentseHealthPlatformMLModelManager}}.
		\item[Super-components:]~\iconcomponent{}~\vpett{\nameref{comp:ComponentseHealthPlatformRiskEstimationProcessor}} $\triangleright$ \iconcomponent{}~\vpett{\nameref{comp:ComponentseHealthPlatform}}
		\item[Super-modules:]~None
		\item[Sub-modules:]~None
		\item[Provided interfaces:]~\iconprovided{}~\vpett{\nameref{int:InterfacesFetchMLModel}}
		\item[Required interfaces:]~\iconrequired{}~\vpett{\nameref{int:InterfacesMLModelMgmt}}
		\item[Visible on diagrams:]~\cref{diag:Component:DecompositionoftheRiskEstimationProcessor,diag:Interaction:Computeclinicalmodelresult}		
	\end{description}

\subsection{MLModelStorageManager}\label{mod:ComponentseHealthPlatformMLModelManagerMLModelStorageManager}
	\begin{description}[noitemsep,nolistsep]
		\item[Responsibility:]~This module is responsible for managing the locally stored \ref{data:DatatypesMLModel}s, triggering the synchronization and retrieval of new \ref{data:DatatypesMLModel}s, and coordinating the updates to the existing \ref{data:DatatypesMLModel}s.
		\item[Super-components:]~\iconcomponent{}~\vpett{\nameref{comp:ComponentseHealthPlatform}} $\triangleright$ \iconcomponent{}~\vpett{\nameref{comp:ComponentseHealthPlatformMLModelManager}}
		\item[Super-modules:]~None
		\item[Sub-modules:]~None
		\item[Provided interfaces:]~\iconprovided{}~\vpett{\nameref{int:InterfacesMLModelMgmt}}
		\item[Required interfaces:]~\iconrequired{}~\vpett{\nameref{int:InterfacesMLModelStorage}}, \iconrequired{}~\vpett{\nameref{int:InterfacesMLModelSync}}
		\item[Visible on diagrams:]~\cref{diag:Component:DecompositionoftheMLModelManager,diag:Interaction:MLModelSynchronization,diag:Interaction:ProcessingMLModelupdates}		
	\end{description}

\subsection{MLModelUpdateProcessor}\label{mod:ComponentseHealthPlatformMLModelManagerMLModelUpdateProcessor}
	\begin{description}[noitemsep,nolistsep]
		\item[Responsibility:]~The MLModelUpdater processes and applies updates to the local \ref{data:DatatypesMLModel}s to prevent updates to these models from requiring to retrieve the full \ref{data:DatatypesMLModel}s.
		\item[Super-components:]~\iconcomponent{}~\vpett{\nameref{comp:ComponentseHealthPlatform}} $\triangleright$ \iconcomponent{}~\vpett{\nameref{comp:ComponentseHealthPlatformMLModelManager}}
		\item[Super-modules:]~None
		\item[Sub-modules:]~None
		\item[Provided interfaces:]~\iconprovided{}~\vpett{\nameref{int:InterfacesMLModelSync}}
		\item[Required interfaces:]~\iconrequired{}~\vpett{\nameref{int:InterfacesMLaaSAPI}}, \iconrequired{}~\vpett{\nameref{int:InterfacesMLaaSService}}, \iconrequired{}~\vpett{\nameref{int:InterfacesMLModelStorage}}, \iconrequired{}~\vpett{\nameref{int:InterfacesMLModelUpdateMgmt}}
		\item[Visible on diagrams:]~\cref{diag:Component:Contextdiagramforthedecompositionview,diag:Component:DecompositionoftheMLModelManager,diag:Interaction:MLModelSynchronization,diag:Interaction:ProcessingMLModelupdates}		
	\end{description}

\subsection{ModelUpdater}\label{mod:ComponentseHealthPlatformMLModelManagerModelUpdater}
	\begin{description}[noitemsep,nolistsep]
		\item[Responsibility:]~This module is responsible for processing \ref{data:DatatypesMLModelUpdate}s and applying them to the \ref{data:DatatypesMLModel}.
		\item[Super-components:]~\iconcomponent{}~\vpett{\nameref{comp:ComponentseHealthPlatform}} $\triangleright$ \iconcomponent{}~\vpett{\nameref{comp:ComponentseHealthPlatformMLModelManager}}
		\item[Super-modules:]~None
		\item[Sub-modules:]~None
		\item[Provided interfaces:]~\iconprovided{}~\vpett{\nameref{int:InterfacesMLModelUpdateMgmt}}
		\item[Required interfaces:]~\iconrequired{}~\vpett{\nameref{int:InterfacesMLAPI}}
		\item[Visible on diagrams:]~\cref{diag:Component:Contextdiagramforthedecompositionview,diag:Component:DecompositionoftheMLModelManager,diag:Interaction:ProcessingMLModelupdates}		
	\end{description}

\subsection{P1CommandRouter}\label{mod:ComponentsP1PhysicianCommandServiceP1CommandRouter}
	\begin{description}[noitemsep,nolistsep]
		\item[Responsibility:]~Provides P1IPhysicianCommand at the component boundary. Thin dispatcher that routes each operation to the matching handler.
		\item[Super-components:]~\iconcomponent{}~\vpett{\nameref{comp:ComponentsP1PhysicianCommandService}}
		\item[Super-modules:]~None
		\item[Sub-modules:]~None
		\item[Provided interfaces:]~\iconprovided{}~\vpett{\nameref{int:InterfacesP1IPhysicianCommand}}
		\item[Required interfaces:]~\iconrequired{}~\vpett{\nameref{int:InterfacesP1IConfigCommand}}, \iconrequired{}~\vpett{\nameref{int:InterfacesP1IOnDemandCommand}}, \iconrequired{}~\vpett{\nameref{int:InterfacesP1IRiskLevelCommand}}
		\item[Visible on diagrams:]~\cref{diag:Component:DecompositionoftheP1PhysicianCommandService}		
	\end{description}

\subsection{P1ConfigurationHandler}\label{mod:ComponentsP1PhysicianCommandServiceP1ConfigurationHandler}
	\begin{description}[noitemsep,nolistsep]
		\item[Responsibility:]~UC7 handler. Overwrites the per-patient \ref{data:DatatypesClinicalModelConfiguration} via P1ClinicalConfigMgmt, invalidates the clinical-model cache via ClinicalModelCacheMgmt, then launches a normal-priority risk-estimation job via LaunchRiskEstimation so Response item 2 (re-assessment after config) is met. The launch carries no fresh sensor data, so the scheduler runs the job against the patient's latest stored package.
		\item[Super-components:]~\iconcomponent{}~\vpett{\nameref{comp:ComponentsP1PhysicianCommandService}}
		\item[Super-modules:]~None
		\item[Sub-modules:]~None
		\item[Provided interfaces:]~\iconprovided{}~\vpett{\nameref{int:InterfacesP1IConfigCommand}}
		\item[Required interfaces:]~\iconrequired{}~\vpett{\nameref{int:InterfacesClinicalModelCacheMgmt}}, \iconrequired{}~\vpett{\nameref{int:InterfacesLaunchRiskEstimation}}, \iconrequired{}~\vpett{\nameref{int:InterfacesP1ClinicalConfigMgmt}}
		\item[Visible on diagrams:]~\cref{diag:Component:DecompositionoftheP1PhysicianCommandService}		
	\end{description}

\subsection{P1CorrelationTracker}\label{mod:ComponentsP1PhysicianCommandServiceP1CorrelationTracker}
	\begin{description}[noitemsep,nolistsep]
		\item[Responsibility:]~Mints CorrelationIds for UC9 requests and matches incoming RiskEvents back to the originating handler. Single locus an Av2 failover must preserve so UC9 result delivery survives.
		\item[Super-components:]~\iconcomponent{}~\vpett{\nameref{comp:ComponentsP1PhysicianCommandService}}
		\item[Super-modules:]~None
		\item[Sub-modules:]~None
		\item[Provided interfaces:]~\iconprovided{}~\vpett{\nameref{int:InterfacesP1ICorrelationTracking}}
		\item[Required interfaces:]~None
		\item[Visible on diagrams:]~\cref{diag:Component:DecompositionoftheP1PhysicianCommandService}		
	\end{description}

\subsection{P1DurableNotificationLog}\label{mod:ComponentsP1NotificationDispatcherP1DurableNotificationLog}
	\begin{description}[noitemsep,nolistsep]
		\item[Responsibility:]~Write-ahead log. Red and yellow notifications are persisted here before being acked from P1IRiskEvents so they survive a crash.
		\item[Super-components:]~\iconcomponent{}~\vpett{\nameref{comp:ComponentsP1NotificationDispatcher}}
		\item[Super-modules:]~None
		\item[Sub-modules:]~None
		\item[Provided interfaces:]~\iconprovided{}~\vpett{\nameref{int:InterfacesP1INotificationLog}}
		\item[Required interfaces:]~None
		\item[Visible on diagrams:]~\cref{diag:Component:DecompositionoftheP1NotificationDispatcher}		
	\end{description}

\subsection{P1EHRProxyModule}\label{mod:ComponentsP1PatientDataServiceP1EHRProxyModule}
	\begin{description}[noitemsep,nolistsep]
		\item[Responsibility:]~Fronts \vpett{\nameref{comp:ComponentsHIS}} for EHR reads and writes via PatientRecordMgmt. Stale-tolerant cache on reads. Cache is invalidated structurally on PMS-initiated writes only. External updates made directly on the \vpett{\nameref{comp:ComponentsHIS}} are not observed, which is acceptable in normal mode and flagged as a sensitivity point.
		\item[Super-components:]~\iconcomponent{}~\vpett{\nameref{comp:ComponentsP1PatientDataService}}
		\item[Super-modules:]~None
		\item[Sub-modules:]~None
		\item[Provided interfaces:]~\iconprovided{}~\vpett{\nameref{int:InterfacesPatientRecordMgmt}}
		\item[Required interfaces:]~\iconrequired{}~\vpett{\nameref{int:InterfacesP1IHISAccess}}
		\item[Visible on diagrams:]~\cref{diag:Component:DecompositionoftheP1PatientDataService}		
	\end{description}

\subsection{P1NotificationInboxModule}\label{mod:ComponentsP1NotificationDispatcherP1NotificationInboxModule}
	\begin{description}[noitemsep,nolistsep]
		\item[Responsibility:]~Provides P1INotificationInbox at the component boundary. Owns per-physician subscription state and the pull surface used by \vpett{\nameref{comp:ComponentsP1PhysicianGateway}} (getPendingNotifications and acknowledgeNotification).
		\item[Super-components:]~\iconcomponent{}~\vpett{\nameref{comp:ComponentsP1NotificationDispatcher}}
		\item[Super-modules:]~None
		\item[Sub-modules:]~None
		\item[Provided interfaces:]~\iconprovided{}~\vpett{\nameref{int:InterfacesP1INotificationInbox}}
		\item[Required interfaces:]~\iconrequired{}~\vpett{\nameref{int:InterfacesP1INotificationLog}}, \iconrequired{}~\vpett{\nameref{int:InterfacesP1IPriorityBuffer}}, \iconrequired{}~\vpett{\nameref{int:InterfacesP1IRecipientRegistry}}
		\item[Visible on diagrams:]~\cref{diag:Component:DecompositionoftheP1NotificationDispatcher,diag:Deployment:Contextdiagramforthedeploymentview}		
	\end{description}

\subsection{P1OnDemandConsultationHandler}\label{mod:ComponentsP1PhysicianCommandServiceP1OnDemandConsultationHandler}
	\begin{description}[noitemsep,nolistsep]
		\item[Responsibility:]~UC9 orchestrator. Mints a \ref{data:DatatypesCorrelationId} via \vpett{\nameref{mod:ComponentsP1PhysicianCommandServiceP1CorrelationTracker}}, fetches sensor data via P1IOnDemandSensorFetch, stores it via `SensorDataMgmt.addSensorData(..., triggerEstimation=false)` so the implicit scheduled job is suppressed, then launches a priority risk job via LaunchRiskEstimation carrying the \ref{data:DatatypesCorrelationId}. When the matching \ref{data:DatatypesRiskEvent} arrives on P1IRiskEventListener, the handler pushes the result back to the originating physician via P1IConsultationResultPush. No polling.
		\item[Super-components:]~\iconcomponent{}~\vpett{\nameref{comp:ComponentsP1PhysicianCommandService}}
		\item[Super-modules:]~None
		\item[Sub-modules:]~None
		\item[Provided interfaces:]~\iconprovided{}~\vpett{\nameref{int:InterfacesP1IOnDemandCommand}}
		\item[Required interfaces:]~\iconrequired{}~\vpett{\nameref{int:InterfacesLaunchRiskEstimation}}, \iconrequired{}~\vpett{\nameref{int:InterfacesP1ICorrelationTracking}}, \iconrequired{}~\vpett{\nameref{int:InterfacesP1IOnDemandSensorFetch}}, \iconrequired{}~\vpett{\nameref{int:InterfacesP1IRiskEvents}}, \iconrequired{}~\vpett{\nameref{int:InterfacesSensorDataMgmt}}
		\item[Visible on diagrams:]~\cref{diag:Component:DecompositionoftheP1PhysicianCommandService}		
	\end{description}

\subsection{P1PatientStatusModule}\label{mod:ComponentsP1PatientDataServiceP1PatientStatusModule}
	\begin{description}[noitemsep,nolistsep]
		\item[Responsibility:]~Owns patient-status records. On `setEstimatedPatientStatus` that actually changes the stored status, the module pushes a \ref{data:DatatypesRiskEvent} to every matching subscription via P1IRiskEventListener.
		\item[Super-components:]~\iconcomponent{}~\vpett{\nameref{comp:ComponentsP1PatientDataService}}
		\item[Super-modules:]~None
		\item[Sub-modules:]~None
		\item[Provided interfaces:]~\iconprovided{}~\vpett{\nameref{int:InterfacesOtherDataMgmt}}, \iconprovided{}~\vpett{\nameref{int:InterfacesP1IRiskEvents}}
		\item[Required interfaces:]~None
		\item[Visible on diagrams:]~\cref{diag:Component:DecompositionoftheP1PatientDataService}		
	\end{description}

\subsection{P1PriorityBuffer}\label{mod:ComponentsP1NotificationDispatcherP1PriorityBuffer}
	\begin{description}[noitemsep,nolistsep]
		\item[Responsibility:]~In-memory priority queue ordering red over yellow over green. Enforces the 100 per minute threshold and drops green first under overload.
		\item[Super-components:]~\iconcomponent{}~\vpett{\nameref{comp:ComponentsP1NotificationDispatcher}}
		\item[Super-modules:]~None
		\item[Sub-modules:]~None
		\item[Provided interfaces:]~\iconprovided{}~\vpett{\nameref{int:InterfacesP1IPriorityBuffer}}
		\item[Required interfaces:]~None
		\item[Visible on diagrams:]~\cref{diag:Component:DecompositionoftheP1NotificationDispatcher}		
	\end{description}

\subsection{P1RecipientRegistry}\label{mod:ComponentsP1NotificationDispatcherP1RecipientRegistry}
	\begin{description}[noitemsep,nolistsep]
		\item[Responsibility:]~Maps a patient-status event to the physicians who should be notified about it. UC5 recipient resolution lives here. Populated via subscribeToNotifications on P1INotificationInbox.
		\item[Super-components:]~\iconcomponent{}~\vpett{\nameref{comp:ComponentsP1NotificationDispatcher}}
		\item[Super-modules:]~None
		\item[Sub-modules:]~None
		\item[Provided interfaces:]~\iconprovided{}~\vpett{\nameref{int:InterfacesP1IRecipientRegistry}}
		\item[Required interfaces:]~None
		\item[Visible on diagrams:]~\cref{diag:Component:DecompositionoftheP1NotificationDispatcher}		
	\end{description}

\subsection{P1RiskEventSubscriber}\label{mod:ComponentsP1NotificationDispatcherP1RiskEventSubscriber}
	\begin{description}[noitemsep,nolistsep]
		\item[Responsibility:]~Provides P1IRiskEventListener at the component boundary. On event, queries \vpett{\nameref{mod:ComponentsP1NotificationDispatcherP1RecipientRegistry}}, persists red and yellow to \vpett{\nameref{mod:ComponentsP1NotificationDispatcherP1DurableNotificationLog}}, then enqueues to \vpett{\nameref{mod:ComponentsP1NotificationDispatcherP1PriorityBuffer}}. Owns the outbound Av2EmergencyDispatch socket for backup-channel escalation of red notifications.
		\item[Super-components:]~\iconcomponent{}~\vpett{\nameref{comp:ComponentsP1NotificationDispatcher}}
		\item[Super-modules:]~None
		\item[Sub-modules:]~None
		\item[Provided interfaces:]~None
		\item[Required interfaces:]~\iconrequired{}~\vpett{\nameref{int:InterfacesP1INotificationLog}}, \iconrequired{}~\vpett{\nameref{int:InterfacesP1IPriorityBuffer}}, \iconrequired{}~\vpett{\nameref{int:InterfacesP1IRecipientRegistry}}, \iconrequired{}~\vpett{\nameref{int:InterfacesP1IRiskEvents}}
		\item[Visible on diagrams:]~\cref{diag:Component:DecompositionoftheP1NotificationDispatcher}		
	\end{description}

\subsection{P1RiskLevelHandler}\label{mod:ComponentsP1PhysicianCommandServiceP1RiskLevelHandler}
	\begin{description}[noitemsep,nolistsep]
		\item[Responsibility:]~UC8 handler. Writes the new internal risk level via `OtherDataMgmt.setEstimatedPatientStatus`, which fires P1IRiskEvents so the tier index, recipient fan-out, and any UC9 correlations all observe the change. Then forwards the change to the EHR via PatientRecordMgmt.
		\item[Super-components:]~\iconcomponent{}~\vpett{\nameref{comp:ComponentsP1PhysicianCommandService}}
		\item[Super-modules:]~None
		\item[Sub-modules:]~None
		\item[Provided interfaces:]~\iconprovided{}~\vpett{\nameref{int:InterfacesP1IRiskLevelCommand}}
		\item[Required interfaces:]~\iconrequired{}~\vpett{\nameref{int:InterfacesPatientRecordMgmt}}
		\item[Visible on diagrams:]~\cref{diag:Component:DecompositionoftheP1PhysicianCommandService}		
	\end{description}

\subsection{P1SensorIngestModule}\label{mod:ComponentsP1PatientDataServiceP1SensorIngestModule}
	\begin{description}[noitemsep,nolistsep]
		\item[Responsibility:]~Accepts and persists raw sensor data. Launches a risk-estimation job via LaunchRiskEstimation only when the caller sets `triggerEstimation=true` (the UC4 default). UC9 callers pass `false` to suppress the implicit scheduled job in favour of the priority one \vpett{\nameref{mod:ComponentsP1PhysicianCommandServiceP1OnDemandConsultationHandler}} launches explicitly.
		\item[Super-components:]~\iconcomponent{}~\vpett{\nameref{comp:ComponentsP1PatientDataService}}
		\item[Super-modules:]~None
		\item[Sub-modules:]~None
		\item[Provided interfaces:]~\iconprovided{}~\vpett{\nameref{int:InterfacesSensorDataMgmt}}
		\item[Required interfaces:]~\iconrequired{}~\vpett{\nameref{int:InterfacesLaunchRiskEstimation}}
		\item[Visible on diagrams:]~\cref{diag:Component:DecompositionoftheP1PatientDataService}		
	\end{description}

\subsection{SensorDataRiskAssessment}\label{mod:ComponentseHealthPlatformRiskEstimationProcessorSensorDataRiskAssessment}
	\begin{description}[noitemsep,nolistsep]
		\item[Responsibility:]~This module performs the actual risk assessment by using the \ref{data:DatatypesMLModel} to make a prediction on the sensor data.
		\item[Super-components:]~\iconcomponent{}~\vpett{\nameref{comp:ComponentseHealthPlatformRiskEstimationProcessor}} $\triangleright$ \iconcomponent{}~\vpett{\nameref{comp:ComponentseHealthPlatform}}
		\item[Super-modules:]~None
		\item[Sub-modules:]~None
		\item[Provided interfaces:]~\iconprovided{}~\vpett{\nameref{int:InterfacesSensorDataAnalysis}}
		\item[Required interfaces:]~\iconrequired{}~\vpett{\nameref{int:InterfacesKVStore}}, \iconrequired{}~\vpett{\nameref{int:InterfacesMLAPI}}
		\item[Visible on diagrams:]~\cref{diag:Component:Contextdiagramforthedecompositionview,diag:Component:DecompositionoftheRiskEstimationProcessor,diag:Interaction:Computeclinicalmodelresult}		
	\end{description}

\pdfinfo{/SAPlugin (v10.0.0)}
% END MODULES

% INTERFACES
\section{Interfaces} \label{sec:interfaces}
  %%%%%%%% ApplyClinicalModel
  \subsection{ApplyClinicalModel}\label{int:InterfacesApplyClinicalModel}
    \begin{description}[noitemsep,nolistsep]
      \item[Provided by:] \iconcomponent{}~\vpett{\nameref{mod:ComponentseHealthPlatformRiskEstimationProcessorJobMLModelSelector}}
      \item[Required by:] \iconcomponent{}~\vpett{\nameref{mod:ComponentseHealthPlatformRiskEstimationProcessorJobProcessor}} 
      \item[Operations:] ~
    \begin{itemize}[noitemsep,nolistsep,leftmargin=-.25cm]
      \item \textsf{String determineMLModelForClinicalModel(\ref{data:DatatypesClinicalModel} clinicalModel, \ref{data:DatatypesClinicalModelConfiguration} clinicalModelConfiguration)}
        \begin{itemize}[noitemsep,nolistsep]
           \item Effect: Apply the \ref{data:DatatypesClinicalModel} to obtain the most appropriate \ref{data:DatatypesMLModel} to use for the calculation. 
           \item Returns: \ref{data:Datatypesorghl7fhirr5modelIdentifier} of the \ref{data:DatatypesMLModel} to use.
           \item Sequence Diagrams: \cref{diag:Interaction:Computeclinicalmodelresult}
        \end{itemize}
    \end{itemize}
      \item[Diagrams:] None
    \end{description}

  %%%%%%%% Av2AuthFailureSignals
  \subsection{Av2AuthFailureSignals}\label{int:InterfacesAv2AuthFailureSignals}
    \begin{description}[noitemsep,nolistsep]
      \item[Provided by:] \iconcomponent{}~\vpett{\nameref{mod:ComponentsAv2CommunicationGatewayAv2AuthenticationValidator}}
      \item[Required by:] \iconcomponent{}~\vpett{\nameref{mod:ComponentsAv2CommunicationGatewayAv2FailureDetector}} 
      \item[Operations:] ~
    \begin{itemize}[noitemsep,nolistsep,leftmargin=-.25cm]
      \item \textsf{void reportAuthFailure(\ref{data:DatatypesGatewayId} gatewayId, string failureReason, \ref{data:DatatypesTimestamp} attemptedAt)}
        \begin{itemize}[noitemsep,nolistsep]
           \item Effect: Reports a single authentication failure to FailureDetector. The failureReason describes why authentication failed, for example unregistered gateway or invalid credentials. FailureDetector aggregates repeated failures from the same gateway to detect systematic attack or misconfiguration. 
           \item Sequence Diagrams: None
        \end{itemize}
    \end{itemize}
      \item[Diagrams:] None
    \end{description}

  %%%%%%%% Av2AvailabilityMonitoring
  \subsection{Av2AvailabilityMonitoring}\label{int:InterfacesAv2AvailabilityMonitoring}
    \begin{description}[noitemsep,nolistsep]
      \item[Provided by:] \iconcomponent{}~\vpett{\nameref{comp:ComponentsAv2CommunicationGateway}}, \iconcomponent{}~\vpett{\nameref{mod:ComponentsAv2CommunicationGatewayAv2HeartbeatTracker}}
      \item[Required by:] \iconcomponent{}~\vpett{\nameref{mod:ComponentsAv2MonitoringServiceAv2MissingDataDetector}}, \iconcomponent{}~\vpett{\nameref{comp:ComponentsAv2MonitoringService}}, \iconcomponent{}~\vpett{\nameref{mod:ComponentsAv2MonitoringServiceAv2SLAMonitor}} 
           \item[Description]: Exposes per-gateway heartbeat data so the MissingDataDetector can identify overdue gateways.
      \item[Operations:] ~
    \begin{itemize}[noitemsep,nolistsep,leftmargin=-.25cm]
      \item \textsf{void getExpectedInterval(\ref{data:DatatypesGatewayId} gatewayId)}
        \begin{itemize}[noitemsep,nolistsep]
           \item Effect: Returns the configured expected transmission interval in seconds for the specified gateway. Used by MissingDataDetector to compute the detection deadline: last heartbeat plus expected interval plus five minute tolerance.
Returns: The expected transmission interval in seconds as an integer. 
           \item Sequence Diagrams: None
        \end{itemize}
      \item \textsf{void getLastHeartbeat(\ref{data:DatatypesGatewayId} gatewayId)}
        \begin{itemize}[noitemsep,nolistsep]
           \item Effect: Returns the timestamp of the most recently received transmission from the specified gateway. Used by MissingDataDetector to determine how long a gateway has been silent and whether the silence exceeds the configured tolerance window.
Returns: The \ref{data:DatatypesTimestamp} of the last received transmission from the specified gateway, or null if no transmission has ever been received. 
           \item Sequence Diagrams: None
        \end{itemize}
      \item \textsf{void recordHeartBeat(\ref{data:DatatypesGatewayId} gatewayId, \ref{data:DatatypesTimestamp} timestamp)}
        \begin{itemize}[noitemsep,nolistsep]
           \item Effect: Records that a transmission was received from the specified gateway at the given timestamp. Updates the last-seen time for that gateway in the HeartbeatTracker. Called by RequestReceiver every time any transmission arrives, regardless of whether it passes authentication, so that connectivity is tracked independently of authentication status. 
           \item Sequence Diagrams: None
        \end{itemize}
      \item \textsf{void registerGateway(\ref{data:DatatypesGatewayId} gatewayId, int expectedIntervalSeconds)}
        \begin{itemize}[noitemsep,nolistsep]
           \item Effect: Registers a new patient gateway in the HeartbeatTracker with its configured expected transmission interval in seconds. Called during patient enrolment so that the HeartbeatTracker immediately begins monitoring for overdue transmissions from that gateway. If a gateway with the given identifier is already registered its interval is updated. 
           \item Sequence Diagrams: None
        \end{itemize}
    \end{itemize}
      \item[Diagrams:] None
    \end{description}

  %%%%%%%% Av2BackupChannelDispatch
  \subsection{Av2BackupChannelDispatch}\label{int:InterfacesAv2BackupChannelDispatch}
    \begin{description}[noitemsep,nolistsep]
      \item[Provided by:] \iconcomponent{}~\vpett{\nameref{mod:ComponentsAv2PatientGatewayAv2BackupChannelManager}}
      \item[Required by:] \iconcomponent{}~\vpett{\nameref{mod:ComponentsAv2PatientGatewayAv2EmergencyNotificationService}} 
           \item[Description]: Published by the BackupChannelManager. Provides operations for sending emergency messages via the backup channel and for querying whether the backup channel is currently available. Consumed by the EmergencyNotificationService module inside the PatientGateway.
      \item[Operations:] ~
    \begin{itemize}[noitemsep,nolistsep,leftmargin=-.25cm]
      \item \textsf{\ref{data:DatatypesbackupChannelType} getBackupChannelType()}
        \begin{itemize}[noitemsep,nolistsep]
           \item Effect: Returns the type of backup communication channel currently configured for this patient gateway. Allows EmergencyNotificationService to log which backup channel was used when dispatching an emergency notification.
 
           \item Returns: A \ref{data:DatatypesbackupChannelType} enum value indicating the backup channel in use, currently always SMS.
           \item Sequence Diagrams: None
        \end{itemize}
      \item \textsf{boolean isBackupChannelAvailable()}
        \begin{itemize}[noitemsep,nolistsep]
           \item Effect: Checks whether the backup SMS channel is currently available and capable of transmitting an emergency notification. Called by EmergencyNotificationService before attempting backup channel dispatch. 
           \item Returns: True if the backup SMS channel is operational and can accept outbound messages. False if the backup channel is also unavailable.
           \item Sequence Diagrams: None
        \end{itemize}
      \item \textsf{\ref{data:DatatypesBackupDeliveryResult} sendViaBackupChannel(\ref{data:DatatypesEmergencyMessage} message, \ref{data:DatatypesGatewayId} gatewayId)}
        \begin{itemize}[noitemsep,nolistsep]
           \item Effect: Transmits the given emergency notification via the backup SMS channel through the telecom network to the SMSGatewayProvider, which delivers it to the EmergencyBackupReceiver on the PMS backend. Called by EmergencyNotificationService when the gateway is in degraded mode and the primary channel is unavailable. 
           \item Returns: A \ref{data:DatatypesBackupDeliveryResult} indicating whether the SMS was successfully submitted to the telecom network, the timestamp of submission, and the channel used.
           \item Sequence Diagrams: None
        \end{itemize}
    \end{itemize}
      \item[Diagrams:] None
    \end{description}

  %%%%%%%% Av2BackupEmergencyIngress
  \subsection{Av2BackupEmergencyIngress}\label{int:InterfacesAv2BackupEmergencyIngress}
    \begin{description}[noitemsep,nolistsep]
      \item[Provided by:] \iconcomponent{}~\vpett{\nameref{comp:ComponentsAv2EmergencyBackupReceiver}}
      \item[Required by:] \iconcomponent{}~\vpett{\nameref{comp:ComponentsAv2PatientGateway}} 
      \item[Operations:] ~
    \begin{itemize}[noitemsep,nolistsep,leftmargin=-.25cm]
      \item \textsf{void receiveEmergencyNotification(\ref{data:DatatypesGatewayId} gatewayId, \ref{data:DatatypesEmergencyMessage} message, string deviceToken)}
        \begin{itemize}[noitemsep,nolistsep]
           \item Effect: Receives an emergency notification transmitted by the patient gateway via the SMS backup channel. Authenticates the sender using the pre-registered device token embedded in the SMS payload. If authentication succeeds, forwards the validated emergency to the MessageBroker. Returns: void. 
           \item Sequence Diagrams: None
        \end{itemize}
    \end{itemize}
      \item[Diagrams:] None
    \end{description}

  %%%%%%%% Av2BatteryStatus
  \subsection{Av2BatteryStatus}\label{int:InterfacesAv2BatteryStatus}
    \begin{description}[noitemsep,nolistsep]
      \item[Provided by:] \iconcomponent{}~\vpett{\nameref{mod:ComponentsAv2PatientGatewayAv2BatteryMonitor}}
      \item[Required by:] \iconcomponent{}~\vpett{\nameref{mod:ComponentsAv2PatientGatewayAv2PatientAlertService}} 
           \item[Description]: Published by the BatteryMonitor. Provides the current battery charge level and supports subscription to threshold-crossing alerts. Consumed by the PatientAlertService and the EmergencyNotificationService module to adapt behavior based on device power state.
      \item[Operations:] ~
    \begin{itemize}[noitemsep,nolistsep,leftmargin=-.25cm]
      \item \textsf{int getBatteryLevel()}
        \begin{itemize}[noitemsep,nolistsep]
           \item Effect: Returns the current battery charge level of the patient gateway device as a percentage between 0 and 100. Called periodically by PatientAlertService to determine whether the 15\% warning threshold has been crossed per Av2 prevention response measures. 
           \item Returns: An integer between 0 and 100 representing the current battery charge percentage.
           \item Sequence Diagrams: None
        \end{itemize}
      \item \textsf{boolean isBelowThreshhold()}
        \begin{itemize}[noitemsep,nolistsep]
           \item Effect: Returns whether the battery charge level has dropped to or below the 15\% warning threshold defined by Av2 prevention response measures. PatientAlertService calls this to decide whether to display a low battery warning to the patient. 
           \item Returns: True if the current battery level is at or below 15\%. False if the battery level is above 15\%.
           \item Sequence Diagrams: None
        \end{itemize}
    \end{itemize}
      \item[Diagrams:] None
    \end{description}

  %%%%%%%% Av2BrokerCommitNotification
  \subsection{Av2BrokerCommitNotification}\label{int:InterfacesAv2BrokerCommitNotification}
    \begin{description}[noitemsep,nolistsep]
      \item[Provided by:] \iconcomponent{}~\vpett{\nameref{comp:ComponentsAv2MessageBroker}}
      \item[Required by:] \iconcomponent{}~\vpett{\nameref{comp:ComponentsAv2AcknowledgementService}} 
      \item[Operations:] ~
    \begin{itemize}[noitemsep,nolistsep,leftmargin=-.25cm]
      \item \textsf{void onMessageCommitted(\ref{data:DatatypesGatewayId} gatewayId, string messageId, \ref{data:DatatypesTimestamp} commitedAt)}
        \begin{itemize}[noitemsep,nolistsep]
           \item Effect: Notifies the AcknowledgementService that the MessageBroker has durably committed a message. This triggers the AcknowledgementService to emit a DeliveryAck back to the originating PatientGateway so it can clear the entry from its LocalBufferingRepository. 
           \item Sequence Diagrams: None
        \end{itemize}
    \end{itemize}
      \item[Diagrams:] None
    \end{description}

  %%%%%%%% Av2BufferedDataDispatch
  \subsection{Av2BufferedDataDispatch}\label{int:InterfacesAv2BufferedDataDispatch}
    \begin{description}[noitemsep,nolistsep]
      \item[Provided by:] \iconcomponent{}~\vpett{\nameref{comp:ComponentsAv2CommunicationGateway}}, \iconcomponent{}~\vpett{\nameref{mod:ComponentsAv2CommunicationGatewayAv2RoutingManager}}
      \item[Required by:] \iconcomponent{}~\vpett{\nameref{comp:ComponentsAv2MessageBroker}} 
           \item[Description]: Published by the RoutingManager at the CommunicationGateway boundary. Consumed by the MessageBroker. Carries validated and routed sensor data packages from the CommunicationGateway to the MessageBroker queue.
      \item[Operations:] ~
    \begin{itemize}[noitemsep,nolistsep,leftmargin=-.25cm]
      \item \textsf{void dispatchEmergencyNotification(\ref{data:DatatypesGatewayId} gatewayId, \ref{data:DatatypesEmergencyMessage} alert, \ref{data:DatatypesTimestamp} recievedAt)}
        \begin{itemize}[noitemsep,nolistsep]
           \item Effect: Forwards a validated emergency notification from the RoutingManager to the MessageBroker's durable queue. Emergency notifications are prioritised over regular sensor data in the queue. The MessageBroker durably commits the notification before returning, ensuring it is not lost even if the DataIngestionService is momentarily unavailable. 
           \item Sequence Diagrams: None
        \end{itemize}
      \item \textsf{void dispatchSensorData(\ref{data:DatatypesGatewayId} gatewayId, \ref{data:DatatypesSensorDataPackage} data, \ref{data:DatatypesTimestamp} recievedAt)}
        \begin{itemize}[noitemsep,nolistsep]
           \item Effect: Forwards a validated and authenticated sensor data package from the RoutingManager to the MessageBroker's durable queue. The receivedAt timestamp is the server-side time at which the CommunicationGateway received the transmission. The MessageBroker durably commits the package before returning, ensuring no data loss if downstream services are temporarily unavailable. 
           \item Sequence Diagrams: None
        \end{itemize}
    \end{itemize}
      \item[Diagrams:] None
    \end{description}

  %%%%%%%% Av2BufferedSensorData
  \subsection{Av2BufferedSensorData}\label{int:InterfacesAv2BufferedSensorData}
    \begin{description}[noitemsep,nolistsep]
      \item[Provided by:] \iconcomponent{}~\vpett{\nameref{mod:ComponentsAv2PatientGatewayAv2LocalBufferingRepository}}
      \item[Required by:] \iconcomponent{}~\vpett{\nameref{mod:ComponentsAv2PatientGatewayAv2RetryAndSynchronizationService}}, \iconcomponent{}~\vpett{\nameref{mod:ComponentsAv2PatientGatewayAv2SensorDataTransmitter}} 
           \item[Description]: Published by the LocalBufferingRepository. Provides a to store BufferedData when in DegradeMode and a way for to read the pending data when synchronizing again with the backend. Consumed by the SensorDataTransmitter and the RetryAndSynchronizationService.
      \item[Operations:] ~
    \begin{itemize}[noitemsep,nolistsep,leftmargin=-.25cm]
      \item \textsf{void clearDeliveredData()}
        \begin{itemize}[noitemsep,nolistsep]
           \item Effect: Removes all entries from the LocalBufferingRepository that have been confirmed as delivered by the backend via DeliveryAck. Called by RetryAndSynchronizationService after receiving acknowledgement to prevent redelivery of already-processed data. 
           \item Sequence Diagrams: None
        \end{itemize}
      \item \textsf{List\textless{}Tuple\textless{}\ref{data:DatatypesSensorDataPackage}, \ref{data:DatatypesTimestamp}\textgreater{}\textgreater{} getPendingData()}
        \begin{itemize}[noitemsep,nolistsep]
           \item Effect: Returns all sensor data packages currently stored in the LocalBufferingRepository awaiting retransmission to the PMS backend. Called by RetryAndSynchronizationService when the gateway reconnects and begins flushing buffered data. 
           \item Returns: An ordered list of tuples each containing a \ref{data:DatatypesSensorDataPackage} and its acquisition timestamp, ordered chronologically from oldest to newest to ensure the backend receives data in the correct order.
           \item Sequence Diagrams: None
        \end{itemize}
      \item \textsf{boolean isRepositoryFull()}
        \begin{itemize}[noitemsep,nolistsep]
           \item Effect: Checks whether the LocalBufferingRepository has reached its maximum storage capacity. 
           \item Returns: True if the repository is full and cannot store additional sensor data packages without overwriting the oldest entries. False if storage capacity remains available.
           \item Sequence Diagrams: None
        \end{itemize}
      \item \textsf{void storeBufferedData()}
        \begin{itemize}[noitemsep,nolistsep]
           \item Effect: Persists a sensor data package and its acquisition timestamp to the LocalBufferingRepository during degraded mode when the primary channel is unavailable. The repository guarantees storage capacity for at least 24 hours of sensor data at the configured transmission rate per Av2 resolution response measures. 
           \item Sequence Diagrams: None
        \end{itemize}
    \end{itemize}
      \item[Diagrams:] None
    \end{description}

  %%%%%%%% Av2CommunicationGapHistory
  \subsection{Av2CommunicationGapHistory}\label{int:InterfacesAv2CommunicationGapHistory}
    \begin{description}[noitemsep,nolistsep]
      \item[Provided by:] \iconcomponent{}~\vpett{\nameref{mod:ComponentsAv2MonitoringServiceAv2CommunicationGapLog}}
      \item[Required by:] \iconcomponent{}~\vpett{\nameref{mod:ComponentsAv2MonitoringServiceAv2RecoveryTracker}} 
      \item[Operations:] ~
    \begin{itemize}[noitemsep,nolistsep,leftmargin=-.25cm]
      \item \textsf{List\textless{}{\colorbox{red!30}{\underline{CommunicationGap}}}\textgreater{} getGapHistory(\ref{data:DatatypesGatewayId} gatewayId)}
        \begin{itemize}[noitemsep,nolistsep]
           \item Effect: Returns all recorded communication gaps for the specified gateway within the given time period. Used by administrators and compliance processes to review the history of connectivity interruptions for a specific patient gateway. 
           \item Returns: A list of CommunicationGap records each containing the gateway identifier, the start and end timestamps of the interruption, and the computed duration. Ordered chronologically from oldest to newest.
           \item Sequence Diagrams: None
        \end{itemize}
      \item \textsf{void recordGap(\ref{data:DatatypesGatewayId} gatewayId, \ref{data:DatatypesTimestamp} start, \ref{data:DatatypesTimestamp} end)}
        \begin{itemize}[noitemsep,nolistsep]
           \item Effect: Persists a completed communication gap record for the specified gateway. Called by RecoveryTracker when a failure is resolved. The gap duration is computed from start and end and stored for compliance reporting and administrator review. 
           \item Sequence Diagrams: None
        \end{itemize}
    \end{itemize}
      \item[Diagrams:] None
    \end{description}

  %%%%%%%% Av2CommunicationGatewayHealthCheck
  \subsection{Av2CommunicationGatewayHealthCheck}\label{int:InterfacesAv2CommunicationGatewayHealthCheck}
    \begin{description}[noitemsep,nolistsep]
      \item[Provided by:] \iconcomponent{}~\vpett{\nameref{comp:ComponentsAv2CommunicationGateway}}, \iconcomponent{}~\vpett{\nameref{mod:ComponentsAv2CommunicationGatewayAv2FailureDetector}}
      \item[Required by:] \iconcomponent{}~\vpett{\nameref{mod:ComponentsAv2MonitoringServiceAv2InternalFailureDetector}}, \iconcomponent{}~\vpett{\nameref{comp:ComponentsAv2MonitoringService}} 
           \item[Description]: Exposes a health check endpoint on the CommunicationGateway so that the InternalFailureDetector in MonitoringService can poll its operational status at high frequency. The CommunicationGateway aggregates the health of its internal subcomponents — RequestReceiver, AuthenticationValidator, RoutingManager, FailureDetector, RetryCoordinator, HeartbeatTracker — and returns a combined status. A DEGRADED or DOWN status or a failure to respond within one second triggers a \ref{data:DatatypesFailureEvent} of type INTERNAL\_SUBSYSTEM\_CRASH in the InternalFailureDetector.
      \item[Operations:] ~
    \begin{itemize}[noitemsep,nolistsep,leftmargin=-.25cm]
      \item \textsf{\ref{data:DatatypesHealthStatus} ping()}
        \begin{itemize}[noitemsep,nolistsep]
           \item Effect: Returns the current top-level operational status of the CommunicationGateway. Returns OK if all subcomponents are functioning normally, DEGRADED if one or more subcomponents are impaired but the gateway is still processing transmissions, DOWN if the gateway cannot process transmissions. 
           \item Returns: the current top-level operational status of the CommunicationGateway
           \item Sequence Diagrams: None
        \end{itemize}
    \end{itemize}
      \item[Diagrams:] None
    \end{description}

  %%%%%%%% Av2ConnectivityStatus
  \subsection{Av2ConnectivityStatus}\label{int:InterfacesAv2ConnectivityStatus}
    \begin{description}[noitemsep,nolistsep]
      \item[Provided by:] \iconcomponent{}~\vpett{\nameref{mod:ComponentsAv2PatientGatewayAv2NetworkConnectivityMonitor}}
      \item[Required by:] \iconcomponent{}~\vpett{\nameref{mod:ComponentsAv2PatientGatewayAv2BackupChannelManager}}, \iconcomponent{}~\vpett{\nameref{mod:ComponentsAv2PatientGatewayAv2EmergencyNotificationService}}, \iconcomponent{}~\vpett{\nameref{mod:ComponentsAv2PatientGatewayAv2PatientAlertService}} 
           \item[Description]: Published by the NetworkConnectivityMonitor. Provides the current network connection state, signal strength, and network type. Consumed by the PatientAlertService, the BackupChannelManagerand EmergencyNotificationService.
      \item[Operations:] ~
    \begin{itemize}[noitemsep,nolistsep,leftmargin=-.25cm]
      \item \textsf{\ref{data:DatatypesSignalStrength} getSignalStrength()}
        \begin{itemize}[noitemsep,nolistsep]
           \item Effect: Returns the current signal strength of the network connection as measured by the NetworkConnectivityMonitor. Called by PatientAlertService to determine whether to warn the patient about insufficient connectivity. 
           \item Returns: A \ref{data:DatatypesSignalStrength} enum value: NONE if no signal, WEAK if signal is insufficient for reliable transmission, MODERATE if functional but potentially lossy, STRONG if reliable connectivity is available.
           \item Sequence Diagrams: None
        \end{itemize}
      \item \textsf{\ref{data:DatatypesNetworkType} getnetworkType()}
        \begin{itemize}[noitemsep,nolistsep]
           \item Effect: Returns the type of network connection currently active on the patient gateway device. Used by SLAMonitor to determine whether the active connection is the cellular network governed by the SLA with the telecom operator. 
           \item Returns: A \ref{data:DatatypesNetworkType} enum value: NONE if no connection, CELLULAR if connected via mobile network, WIFI if connected via WiFi.
           \item Sequence Diagrams: None
        \end{itemize}
      \item \textsf{boolean isConnected()}
        \begin{itemize}[noitemsep,nolistsep]
           \item Effect: Returns whether the patient gateway currently has any active network connection. Used by PatientAlertService to trigger a connectivity warning when the patient is completely offline, and by BackupChannelManager to decide if the backup SMS channel should be activated. 
           \item Returns: True if a network connection of any strength is currently active. False if the device has no network connectivity.
           \item Sequence Diagrams: None
        \end{itemize}
    \end{itemize}
      \item[Diagrams:] None
    \end{description}

  %%%%%%%% Av2DegradeModeStatus
  \subsection{Av2DegradeModeStatus}\label{int:InterfacesAv2DegradeModeStatus}
    \begin{description}[noitemsep,nolistsep]
      \item[Provided by:] \iconcomponent{}~\vpett{\nameref{mod:ComponentsAv2PatientGatewayAv2DegradeModeManager}}
      \item[Required by:] \iconcomponent{}~\vpett{\nameref{mod:ComponentsAv2PatientGatewayAv2BackupChannelManager}}, \iconcomponent{}~\vpett{\nameref{mod:ComponentsAv2PatientGatewayAv2EmergencyNotificationService}}, \iconcomponent{}~\vpett{\nameref{mod:ComponentsAv2PatientGatewayAv2RetryAndSynchronizationService}}, \iconcomponent{}~\vpett{\nameref{mod:ComponentsAv2PatientGatewayAv2SensorDataTransmitter}} 
           \item[Description]: Published by the DegradeModeManager. Reports whether the gateway is currently in degraded mode, the reason for entering degraded mode, and the timestamp at which degraded mode began. Consumed by the SensorDataTransmitter, RetryAndSynchronizationService, BackupChannelManager, and EmergencyNotificationService module.
      \item[Operations:] ~
    \begin{itemize}[noitemsep,nolistsep,leftmargin=-.25cm]
      \item \textsf{void enterDegradeMode(\ref{data:DatatypesDegradedModeReason} reason)}
        \begin{itemize}[noitemsep,nolistsep]
           \item Effect: Transitions the gateway into degraded mode, recording the reason and the current timestamp as the degraded-since time. Called by DegradeModeManager when ServerConnectivityMonitor reports the backend is unreachable. All components consulting DegradeModeStatus will subsequently receive true from isInDegradeMode. 
           \item Sequence Diagrams: None
        \end{itemize}
      \item \textsf{void exitDegradeMode()}
        \begin{itemize}[noitemsep,nolistsep]
           \item Effect: Transitions the gateway back to normal operating mode from degraded mode. Called by DegradeModeManager when ServerConnectivityMonitor reports the backend is reachable again. Clears the degraded-since timestamp. 
           \item Sequence Diagrams: None
        \end{itemize}
      \item \textsf{\ref{data:DatatypesTimestamp} getdegradedSince()}
        \begin{itemize}[noitemsep,nolistsep]
           \item Effect: Returns the timestamp at which the gateway entered its current degraded mode period. Used by RetryAndSynchronizationService to calculate the total duration of the outage and to determine which buffered readings need to be retransmitted.
 
           \item Returns: The \ref{data:DatatypesTimestamp} at which degraded mode was entered. Returns null if the gateway is not currently in degraded mode.
           \item Sequence Diagrams: None
        \end{itemize}
      \item \textsf{boolean isInDegradeMode()}
        \begin{itemize}[noitemsep,nolistsep]
           \item Effect: Returns whether the patient gateway is currently operating in degraded mode. Consulted by SensorDataTransmitter, RetryAndSynchronizationService, BackupChannelManager, and EmergencyNotificationService to adapt their behaviour accordingly. 
           \item Returns: True if the gateway is currently in degraded mode. False if operating normally.
           \item Sequence Diagrams: None
        \end{itemize}
    \end{itemize}
      \item[Diagrams:] None
    \end{description}

  %%%%%%%% Av2DeliveryAck
  \subsection{Av2DeliveryAck}\label{int:InterfacesAv2DeliveryAck}
    \begin{description}[noitemsep,nolistsep]
      \item[Provided by:] \iconcomponent{}~\vpett{\nameref{comp:ComponentsAv2AcknowledgementService}}
      \item[Required by:] \iconcomponent{}~\vpett{\nameref{comp:ComponentsAv2PatientGateway}}, \iconcomponent{}~\vpett{\nameref{mod:ComponentsAv2PatientGatewayAv2RetryAndSynchronizationService}} 
      \item[Operations:] ~
    \begin{itemize}[noitemsep,nolistsep,leftmargin=-.25cm]
      \item \textsf{void acknowledgeDelivery(\ref{data:DatatypesGatewayId} gatewayId, string messageId, \ref{data:DatatypesTimestamp} acknowledgedAt)}
        \begin{itemize}[noitemsep,nolistsep]
           \item Effect: Signals to the PatientGateway (and its RetryAndSynchronizationService) that a previously transmitted message has been durably committed by the backend. Upon receipt, the gateway clears the corresponding entry from its LocalBufferingRepository. 
           \item Sequence Diagrams: None
        \end{itemize}
    \end{itemize}
      \item[Diagrams:] None
    \end{description}

  %%%%%%%% Av2DetectedFailures
  \subsection{Av2DetectedFailures}\label{int:InterfacesAv2DetectedFailures}
    \begin{description}[noitemsep,nolistsep]
      \item[Provided by:] \iconcomponent{}~\vpett{\nameref{mod:ComponentsAv2MonitoringServiceAv2InternalFailureDetector}}, \iconcomponent{}~\vpett{\nameref{mod:ComponentsAv2MonitoringServiceAv2MissingDataDetector}}, \iconcomponent{}~\vpett{\nameref{mod:ComponentsAv2MonitoringServiceAv2SLAMonitor}}
      \item[Required by:] \iconcomponent{}~\vpett{\nameref{mod:ComponentsAv2MonitoringServiceAv2FailureClassifier}} 
           \item[Description]: Carries \ref{data:DatatypesFailureEvent} objects from the MissingDataDetector and InternalFailureDetector to the FailureClassifier. Both detectors provide this interface; the FailureClassifier requires it from both.
      \item[Operations:] ~
    \begin{itemize}[noitemsep,nolistsep,leftmargin=-.25cm]
      \item \textsf{List\textless{}\ref{data:DatatypesFailureEvent}\textgreater{} getActiveFailures()}
        \begin{itemize}[noitemsep,nolistsep]
           \item Effect: Returns all currently active failures that have been reported but not yet resolved. Used by FailureClassifier to avoid processing duplicate failure reports for the same ongoing failure event. 
           \item Returns: A list of all active \ref{data:DatatypesFailureEvent} objects currently being tracked, each containing the failure type, affected gateway or subsystem, and detection timestamp.
           \item Sequence Diagrams: None
        \end{itemize}
      \item \textsf{void reportFailure(\ref{data:DatatypesFailureEvent} failure)}
        \begin{itemize}[noitemsep,nolistsep]
           \item Effect: Reports a newly detected failure to the FailureClassifier. Called by MissingDataDetector when a gateway exceeds its expected transmission interval plus the five minute tolerance, by InternalFailureDetector when a subsystem health check fails, and by SLAMonitor when an SLA threshold is breached. The FailureClassifier enriches the event with escalation level and notification deadline. 
           \item Sequence Diagrams: None
        \end{itemize}
    \end{itemize}
      \item[Diagrams:] None
    \end{description}

  %%%%%%%% Av2EmergencyDispatch
  \subsection{Av2EmergencyDispatch}\label{int:InterfacesAv2EmergencyDispatch}
    \begin{description}[noitemsep,nolistsep]
      \item[Provided by:] \iconcomponent{}~\vpett{\nameref{mod:ComponentsAv2PatientGatewayAv2EmergencyNotificationService}}, \iconcomponent{}~\vpett{\nameref{comp:ComponentsAv2PatientGateway}}
      \item[Required by:] \iconcomponent{}~\vpett{\nameref{comp:ComponentsP1NotificationDispatcher}} 
           \item[Description]: Published by the EmergencyNotificationService module inside the PatientGateway. Provides a single operation for sending an emergency notification, routing it via the primary channel in normal mode or the backup channel in degraded mode. Consumed by sensor-level risk detection logic on the gateway.
      \item[Operations:] ~
    \begin{itemize}[noitemsep,nolistsep,leftmargin=-.25cm]
      \item \textsf{void sendEmergencyNotification(\ref{data:DatatypesEmergencyMessage} alert)}
        \begin{itemize}[noitemsep,nolistsep]
           \item Effect: Sends an emergency 
           \item Sequence Diagrams: None
        \end{itemize}
    \end{itemize}
      \item[Diagrams:] None
    \end{description}

  %%%%%%%% Av2EscalationEvents
  \subsection{Av2EscalationEvents}\label{int:InterfacesAv2EscalationEvents}
    \begin{description}[noitemsep,nolistsep]
      \item[Provided by:] \iconcomponent{}~\vpett{\nameref{mod:ComponentsAv2MonitoringServiceAv2FailureClassifier}}
      \item[Required by:] \iconcomponent{}~\vpett{\nameref{mod:ComponentsAv2MonitoringServiceAv2RecoveryTracker}}, \iconcomponent{}~\vpett{\nameref{mod:ComponentsAv2MonitoringServiceAv2StakeholderNotifier}} 
           \item[Description]: Carries \ref{data:DatatypesClassifiedFailure} objects from the FailureClassifier to both the StakeholderNotifier and the RecoveryTracker.
      \item[Operations:] ~
    \begin{itemize}[noitemsep,nolistsep,leftmargin=-.25cm]
      \item \textsf{\ref{data:DatatypesClassifiedFailure} classifyFailure(\ref{data:DatatypesFailureEvent} failure)}
        \begin{itemize}[noitemsep,nolistsep]
           \item Effect: Receives a raw \ref{data:DatatypesFailureEvent} from FailureClassifier and returns a \ref{data:DatatypesClassifiedFailure} enriched with the escalation level and notification deadline derived from the affected patient's risk level. High-risk patients produce a deadline of 2 minutes, medium-risk 5 minutes, low-risk 10 minutes, per Av2 resolution response measures. 
           \item Returns: A \ref{data:DatatypesClassifiedFailure} containing the original \ref{data:DatatypesFailureEvent}, the computed \ref{data:DatatypesEscalationLevel}, and the \ref{data:DatatypesTimestamp} by which the relevant stakeholder must be notified.
           \item Sequence Diagrams: None
        \end{itemize}
      \item \textsf{\ref{data:DatatypesEscalationLevel} getEscalationLevel(\ref{data:DatatypesFailureEvent} failure)}
        \begin{itemize}[noitemsep,nolistsep]
           \item Effect: Returns the escalation level that would be assigned to the given failure based on the affected patient's current risk level, without creating a full \ref{data:DatatypesClassifiedFailure}. Used by StakeholderNotifier to prioritise notification dispatch when multiple failures are active simultaneously. 
           \item Returns: The \ref{data:DatatypesEscalationLevel} for the failure: HIGH for high-risk patients, MEDIUM for medium-risk, LOW for low-risk.
           \item Sequence Diagrams: None
        \end{itemize}
    \end{itemize}
      \item[Diagrams:] None
    \end{description}

  %%%%%%%% Av2FailureEvents
  \subsection{Av2FailureEvents}\label{int:InterfacesAv2FailureEvents}
    \begin{description}[noitemsep,nolistsep]
      \item[Provided by:] \iconcomponent{}~\vpett{\nameref{mod:ComponentsAv2CommunicationGatewayAv2FailureDetector}}
      \item[Required by:] \iconcomponent{}~\vpett{\nameref{mod:ComponentsAv2CommunicationGatewayAv2RetryCoordinator}} 
           \item[Description]: Published by the FailureDetector. Carries connection-level and internal failure notifications to the RetryCoordinator. Only used inside the CommunicationGateway.
      \item[Operations:] ~
    \begin{itemize}[noitemsep,nolistsep,leftmargin=-.25cm]
      \item \textsf{List\textless{}{\colorbox{red!30}{\underline{ConnectionFailure}}}\textgreater{} getActiveConnectionFailures()}
        \begin{itemize}[noitemsep,nolistsep]
           \item Effect: Returns the list of currently active connection failures being tracked by FailureDetector for which retry attempts are ongoing. Used by RetryCoordinator to avoid scheduling duplicate retries for the same gateway. 
           \item Returns: A list of the current connection failures
           \item Sequence Diagrams: None
        \end{itemize}
      \item \textsf{void reportConnectionFailure(\ref{data:DatatypesGatewayId} gatewayId, \ref{data:DatatypesTimestamp} attempt)}
        \begin{itemize}[noitemsep,nolistsep]
           \item Effect: Reports a connection-level or delivery failure detected within the CommunicationGateway to the RetryCoordinator. Called by FailureDetector when it receives a RoutingFailureSignal or AuthFailureSignal that warrants a retry attempt. The attempt timestamp is used by RetryCoordinator to compute the next retry time using exponential backoff. 
           \item Sequence Diagrams: None
        \end{itemize}
    \end{itemize}
      \item[Diagrams:] None
    \end{description}

  %%%%%%%% Av2GatewayDataIngress
  \subsection{Av2GatewayDataIngress}\label{int:InterfacesAv2GatewayDataIngress}
    \begin{description}[noitemsep,nolistsep]
      \item[Provided by:] \iconcomponent{}~\vpett{\nameref{comp:ComponentsAv2PatientGateway}}, \iconcomponent{}~\vpett{\nameref{mod:ComponentsAv2PatientGatewayAv2SensorDataTransmitter}}
      \item[Required by:] \iconcomponent{}~\vpett{\nameref{comp:ComponentsAv2CommunicationGateway}}, \iconcomponent{}~\vpett{\nameref{mod:ComponentsAv2CommunicationGatewayAv2RequestReceiver}} 
           \item[Description]: Published by the PatientGateway at its external boundary. Provides the entry point through which the gateway transmits sensor data and emergency notifications to the CommunicationGateway. In degraded mode outbound transmissions are buffered locally rather than sent immediately.
      \item[Operations:] ~
    \begin{itemize}[noitemsep,nolistsep,leftmargin=-.25cm]
      \item \textsf{{\colorbox{red!30}{\underline{TransmissionAck}}} transmitEmergencyNotification(\ref{data:DatatypesGatewayId} gatewayId, \ref{data:DatatypesEmergencyMessage} alert, \ref{data:DatatypesTimestamp} sendAt)}
        \begin{itemize}[noitemsep,nolistsep]
           \item Effect: Transmits an emergency notification from the PatientGateway to the backend CommunicationGateway over the primary telecom channel. Used when the primary channel is available. If the primary channel is unavailable the EmergencyNotificationService delegates to BackupChannelDispatch instead and this operation is not called.
Returns: A TransmissionAck confirming the backend has accepted and queued the emergency notification for downstream processing. 
           \item Returns: A TransmissionAck confirming the backend has accepted and queued the emergency notification for downstream processing.
           \item Sequence Diagrams: None
        \end{itemize}
      \item \textsf{{\colorbox{red!30}{\underline{TransmissionAck}}} transmitSensorData(\ref{data:DatatypesGatewayId} gatewayId, \ref{data:DatatypesSensorDataPackage} data, \ref{data:DatatypesTimestamp} sendAt)}
        \begin{itemize}[noitemsep,nolistsep]
           \item Effect: Transmits a sensor data package from the PatientGateway to the backend CommunicationGateway over the primary channel. The sendAt timestamp records the moment of transmission at the gateway. The backend responds with a TransmissionAck confirming successful receipt. 
           \item Returns: A TransmissionAck confirming the backend has accepted the transmission. Contains the message identifier assigned by the backend and the server-side receipt timestamp.
           \item Sequence Diagrams: None
        \end{itemize}
    \end{itemize}
      \item[Diagrams:] None
    \end{description}

  %%%%%%%% Av2GatewayResyncCommand
  \subsection{Av2GatewayResyncCommand}\label{int:InterfacesAv2GatewayResyncCommand}
    \begin{description}[noitemsep,nolistsep]
      \item[Provided by:] \iconcomponent{}~\vpett{\nameref{comp:ComponentsAv2RecoverySyncService}}
      \item[Required by:] \iconcomponent{}~\vpett{\nameref{comp:ComponentsAv2PatientGateway}}, \iconcomponent{}~\vpett{\nameref{mod:ComponentsAv2PatientGatewayAv2RetryAndSynchronizationService}} 
      \item[Operations:] ~
    \begin{itemize}[noitemsep,nolistsep,leftmargin=-.25cm]
      \item \textsf{void triggerResync(\ref{data:DatatypesGatewayId} gatewayId, \ref{data:DatatypesTimestamp} fromTimestamp)}
        \begin{itemize}[noitemsep,nolistsep]
           \item Effect: Instructs the specified PatientGateway to flush its LocalBufferingRepository and retransmit all buffered sensor data recorded after fromTimestamp. Called by RecoverySyncService after a communication failure is resolved 
           \item Sequence Diagrams: None
        \end{itemize}
    \end{itemize}
      \item[Diagrams:] None
    \end{description}

  %%%%%%%% Av2HeartbeatUpdate
  \subsection{Av2HeartbeatUpdate}\label{int:InterfacesAv2HeartbeatUpdate}
    \begin{description}[noitemsep,nolistsep]
      \item[Provided by:] \iconcomponent{}~\vpett{\nameref{mod:ComponentsAv2CommunicationGatewayAv2RequestReceiver}}
      \item[Required by:] \iconcomponent{}~\vpett{\nameref{mod:ComponentsAv2CommunicationGatewayAv2HeartbeatTracker}} 
           \item[Description]: Used by the RequestReceiver to notify the HeartbeatTracker that a transmission has been received from a gateway. Only used inside the CommunicationGateway.
      \item[Operations:] ~
    \begin{itemize}[noitemsep,nolistsep,leftmargin=-.25cm]
      \item \textsf{void recordTransmission(\ref{data:DatatypesGatewayId} gatewayId, \ref{data:DatatypesTimestamp} recievedAt)}
        \begin{itemize}[noitemsep,nolistsep]
           \item Effect: Notifies HeartbeatTracker that a transmission was received from the specified gateway at the given server-side timestamp. Called by RequestReceiver for every inbound transmission regardless of whether it passes authentication. Updates the last-seen time for that gateway so MissingDataDetector can accurately compute how long the gateway has been silent. 
           \item Sequence Diagrams: None
        \end{itemize}
    \end{itemize}
      \item[Diagrams:] None
    \end{description}

  %%%%%%%% Av2InboundGatewayData
  \subsection{Av2InboundGatewayData}\label{int:InterfacesAv2InboundGatewayData}
    \begin{description}[noitemsep,nolistsep]
      \item[Provided by:] \iconcomponent{}~\vpett{\nameref{mod:ComponentsAv2CommunicationGatewayAv2RequestReceiver}}
      \item[Required by:] \iconcomponent{}~\vpett{\nameref{mod:ComponentsAv2CommunicationGatewayAv2AuthenticationValidator}} 
           \item[Description]: Carries parsed and format-validated inbound transmissions from the RequestReceiver to the AuthenticationValidator. Only used inside the CommunicationGateway.
      \item[Operations:] ~
    \begin{itemize}[noitemsep,nolistsep,leftmargin=-.25cm]
      \item \textsf{void forwardParsedTransmission(\ref{data:DatatypesGatewayId} gatewayId, {\colorbox{red!30}{\underline{ParsedTransmission}}} transmission, \ref{data:DatatypesTimestamp} recievedAt)}
        \begin{itemize}[noitemsep,nolistsep]
           \item Effect: Forwards a format-validated and parsed inbound transmission from RequestReceiver to AuthenticationValidator for credential verification. The receivedAt timestamp is the server-side time at which the CommunicationGateway received the raw transmission. Only transmissions that pass initial format validation are forwarded, malformed transmissions are dropped by RequestReceiver before this interface is called. 
           \item Sequence Diagrams: None
        \end{itemize}
    \end{itemize}
      \item[Diagrams:] None
    \end{description}

  %%%%%%%% Av2MessageBrokerHealthCheck
  \subsection{Av2MessageBrokerHealthCheck}\label{int:InterfacesAv2MessageBrokerHealthCheck}
    \begin{description}[noitemsep,nolistsep]
      \item[Provided by:] \iconcomponent{}~\vpett{\nameref{comp:ComponentsAv2MessageBroker}}
      \item[Required by:] \iconcomponent{}~\vpett{\nameref{mod:ComponentsAv2MonitoringServiceAv2InternalFailureDetector}}, \iconcomponent{}~\vpett{\nameref{comp:ComponentsAv2MonitoringService}} 
           \item[Description]: Exposes a health check endpoint on the MessageBroker so that the InternalFailureDetector in MonitoringService can independently poll its operational status. The MessageBroker reports its own queue depth, throughput, and availability. A failure to respond within one second or a DOWN status triggers a \ref{data:DatatypesFailureEvent} of type COMMUNICATION\_CHANNEL\_DOWN in the InternalFailureDetector, since MessageBroker failure directly breaks the data path between the CommunicationGateway and the DataIngestionService.
      \item[Operations:] ~
    \begin{itemize}[noitemsep,nolistsep,leftmargin=-.25cm]
      \item \textsf{\ref{data:DatatypesHealthStatus} ping()}
        \begin{itemize}[noitemsep,nolistsep]
           \item Effect: Returns the current operational status of the MessageBroker. Returns OK if the broker is accepting and delivering messages normally, DEGRADED if queue depth is abnormally high or throughput is reduced, DOWN if the broker is not accepting new messages. 
           \item Returns: the current operational status of the MessageBroker
           \item Sequence Diagrams: None
        \end{itemize}
    \end{itemize}
      \item[Diagrams:] None
    \end{description}

  %%%%%%%% Av2NotificationDispatch
  \subsection{Av2NotificationDispatch}\label{int:InterfacesAv2NotificationDispatch}
    \begin{description}[noitemsep,nolistsep]
      \item[Provided by:] \iconcomponent{}~\vpett{\nameref{comp:ComponentsAv2MonitoringService}}, \iconcomponent{}~\vpett{\nameref{mod:ComponentsAv2MonitoringServiceAv2StakeholderNotifier}}
      \item[Required by:] \iconcomponent{}~\vpett{\nameref{comp:ComponentsAv2AdminPortal}} 
           \item[Description]: The outbound notification channel through which the StakeholderNotifier sends alerts to stakeholders.
      \item[Operations:] ~
    \begin{itemize}[noitemsep,nolistsep,leftmargin=-.25cm]
      \item \textsf{void notifyAdministrator(\ref{data:DatatypesClassifiedFailure} failure, string message)}
        \begin{itemize}[noitemsep,nolistsep]
           \item Effect: Dispatches a notification to the system administrator describing the classified failure. The message is delivered via the NotificationDispatcher within the deadline encoded in the \ref{data:DatatypesClassifiedFailure} — 2 minutes for HIGH escalation level, 5 minutes for MEDIUM, 10 minutes for LOW — per Av2 resolution response measures. Must succeed in 90\% of cases within the deadline. 
           \item Sequence Diagrams: None
        \end{itemize}
      \item \textsf{void notifyPatientContact(\ref{data:DatatypesPatientId} patientId, \ref{data:DatatypesClassifiedFailure} failure, string message)}
        \begin{itemize}[noitemsep,nolistsep]
           \item Effect: Dispatches a notification to the registered contact person of the patient identified by patientId. Used for GATEWAY\_UNREACHABLE and MISSING\_SENSOR\_DATA failures where the patient or their contact must be informed. The patient's contact information is retrieved from \vpett{\nameref{comp:ComponentsPatientRecordRepository}} via PatientContactLookup before dispatch. 
           \item Sequence Diagrams: None
        \end{itemize}
    \end{itemize}
      \item[Diagrams:] None
    \end{description}

  %%%%%%%% Av2PatientContactLookup
  \subsection{Av2PatientContactLookup}\label{int:InterfacesAv2PatientContactLookup}
    \begin{description}[noitemsep,nolistsep]
      \item[Provided by:] \iconcomponent{}~\vpett{\nameref{comp:ComponentsPatientRecordRepository}}
      \item[Required by:] \iconcomponent{}~\vpett{\nameref{comp:ComponentsAv2MonitoringService}}, \iconcomponent{}~\vpett{\nameref{mod:ComponentsAv2MonitoringServiceAv2StakeholderNotifier}} 
      \item[Operations:] ~
    \begin{itemize}[noitemsep,nolistsep,leftmargin=-.25cm]
      \item \textsf{\ref{data:Datatypesorghl7fhirr5modelPatientContactComponent} getContactInfo(\ref{data:DatatypesPatientId} id)}
        \begin{itemize}[noitemsep,nolistsep]
           \item Effect: Retrieves the contact information of the registered contact person for the specified patient from the \vpett{\nameref{comp:ComponentsPatientRecordRepository}}. Used by StakeholderNotifier to obtain the phone number or email address needed to notify the patient's contact when a GATEWAY\_UNREACHABLE or MISSING\_SENSOR\_DATA failure is detected. 
           \item Returns: A ContactInfo object containing the name, phone number, and email address of the patient's registered contact person.
           \item Sequence Diagrams: None
        \end{itemize}
    \end{itemize}
      \item[Diagrams:] None
    \end{description}

  %%%%%%%% Av2PatientStatusOverview
  \subsection{Av2PatientStatusOverview}\label{int:InterfacesAv2PatientStatusOverview}
    \begin{description}[noitemsep,nolistsep]
      \item[Provided by:] \iconcomponent{}~\vpett{\nameref{comp:ComponentsPatientRecordRepository}}
      \item[Required by:] \iconcomponent{}~\vpett{\nameref{comp:ComponentsAv2CommunicationGateway}}, \iconcomponent{}~\vpett{\nameref{mod:ComponentsAv2MonitoringServiceAv2FailureClassifier}}, \iconcomponent{}~\vpett{\nameref{mod:ComponentsAv2CommunicationGatewayAv2GatewayRegistry}}, \iconcomponent{}~\vpett{\nameref{comp:ComponentsPatientStatusService}} 
      \item[Operations:] ~
    \begin{itemize}[noitemsep,nolistsep,leftmargin=-.25cm]
      \item \textsf{\ref{data:DatatypesPatientStatus} getPatientStatus(\ref{data:DatatypesPatientId} patientId)}
        \begin{itemize}[noitemsep,nolistsep]
           \item Effect: Returns the current risk level of the specified patient as stored in the \vpett{\nameref{comp:ComponentsPatientRecordRepository}}. Used by FailureClassifier to determine escalation urgency, and by GatewayRegistery to calibrate heartbeat intervals 
           \item Returns: The current \ref{data:DatatypesPatientStatus} of the specified patient including their risk level classification.
           \item Sequence Diagrams: None
        \end{itemize}
    \end{itemize}
      \item[Diagrams:] None
    \end{description}

  %%%%%%%% Av2QueuedSensorDataMgmt
  \subsection{Av2QueuedSensorDataMgmt}\label{int:InterfacesAv2QueuedSensorDataMgmt}
    \begin{description}[noitemsep,nolistsep]
      \item[Provided by:] \iconcomponent{}~\vpett{\nameref{comp:ComponentsAv2MessageBroker}}
      \item[Required by:] \iconcomponent{}~\vpett{\nameref{comp:ComponentsAv2DataIngestionService}}, \iconcomponent{}~\vpett{\nameref{comp:ComponentsAv2EmergencyBackupReceiver}}, \iconcomponent{}~\vpett{\nameref{comp:ComponentsAv2RecoverySyncService}} 
      \item[Operations:] ~
    \begin{itemize}[noitemsep,nolistsep,leftmargin=-.25cm]
      \item \textsf{\ref{data:DatatypesSensorDataPackage} consume()}
        \begin{itemize}[noitemsep,nolistsep]
           \item Effect: Blocks until a sensor data package is available and returns it for processing. Called by DataIngestionService and RecoverySyncService to drain the queue. 
           \item Returns: The next \ref{data:DatatypesSensorDataPackage} dequeued from the MessageBroker in delivery order.
           \item Sequence Diagrams: None
        \end{itemize}
      \item \textsf{void enqueue(\ref{data:DatatypesGatewayId} gatewayId, \ref{data:DatatypesSensorDataPackage} data, \ref{data:DatatypesTimestamp} receivedAt)}
        \begin{itemize}[noitemsep,nolistsep]
           \item Effect: Adds a sensor data package to the MessageBroker's durable queue for downstream processing by the DataIngestionService. Called by the RoutingManager after a validated transmission is received. 
           \item Sequence Diagrams: None
        \end{itemize}
      \item \textsf{int getQueueDepth()}
        \begin{itemize}[noitemsep,nolistsep]
           \item Effect: Returns the current number of messages waiting in the queue. Used by RecoverySyncService to determine when the backlog from a degraded-mode recovery has been fully drained. 
           \item Returns: The number of messages currently waiting in the queue as an integer.
           \item Sequence Diagrams: None
        \end{itemize}
    \end{itemize}
      \item[Diagrams:] None
    \end{description}

  %%%%%%%% Av2RecoveryCoordination
  \subsection{Av2RecoveryCoordination}\label{int:InterfacesAv2RecoveryCoordination}
    \begin{description}[noitemsep,nolistsep]
      \item[Provided by:] \iconcomponent{}~\vpett{\nameref{comp:ComponentsAv2MonitoringService}}, \iconcomponent{}~\vpett{\nameref{mod:ComponentsAv2MonitoringServiceAv2RecoveryTracker}}
      \item[Required by:] \iconcomponent{}~\vpett{\nameref{comp:ComponentsAv2RecoverySyncService}}, \iconcomponent{}~\vpett{\nameref{comp:ComponentsAv2RedeploymentCoordinator}} 
      \item[Operations:] ~
    \begin{itemize}[noitemsep,nolistsep,leftmargin=-.25cm]
      \item \textsf{void initiateRecovery(\ref{data:DatatypesFailureEventId} id)}
        \begin{itemize}[noitemsep,nolistsep]
           \item Effect: Signals RecoverySyncService and RedeploymentCoordinator to begin post-failure recovery actions for the specified failure. RecoverySyncService drains any queued sensor data that accumulated in the MessageBroker during the outage and instructs the affected PatientGateway to flush its LocalBufferingRepository. Called by RecoveryTracker when a failure is resolved. 
           \item Sequence Diagrams: None
        \end{itemize}
      \item \textsf{void onRecoveryComplete(\ref{data:DatatypesFailureEventId} id, \ref{data:DatatypesTimestamp} resolvedAt)}
        \begin{itemize}[noitemsep,nolistsep]
           \item Effect: Notifies the MonitoringService that all recovery actions for the specified failure have been completed. Used to close the failure record, update SLA metrics with the total outage duration, and clear the failure from the active failure list. 
           \item Sequence Diagrams: None
        \end{itemize}
    \end{itemize}
      \item[Diagrams:] None
    \end{description}

  %%%%%%%% Av2RecoveryStatus
  \subsection{Av2RecoveryStatus}\label{int:InterfacesAv2RecoveryStatus}
    \begin{description}[noitemsep,nolistsep]
      \item[Provided by:] \iconcomponent{}~\vpett{\nameref{mod:ComponentsAv2MonitoringServiceAv2FailureClassifier}}
      \item[Required by:] \iconcomponent{}~\vpett{\nameref{mod:ComponentsAv2MonitoringServiceAv2RecoveryTracker}} 
           \item[Description]: Signals the RecoveryTracker when a previously detected failure has been resolved, allowing it to close the failure record and trigger recovery coordination.
      \item[Operations:] ~
    \begin{itemize}[noitemsep,nolistsep,leftmargin=-.25cm]
      \item \textsf{void onFailureResolved(\ref{data:DatatypesFailureEventId} failureId, \ref{data:DatatypesTimestamp} resolvedAt)}
        \begin{itemize}[noitemsep,nolistsep]
           \item Effect: Signals RecoveryTracker that the failure identified by failureId has been resolved at the given timestamp. RecoveryTracker uses resolvedAt together with the failure's detectedAt timestamp to compute the total communication gap duration, closes the failure record, and triggers RecoveryCoordination to begin post-failure synchronisation. 
           \item Sequence Diagrams: None
        \end{itemize}
    \end{itemize}
      \item[Diagrams:] None
    \end{description}

  %%%%%%%% Av2RedeploymentControl
  \subsection{Av2RedeploymentControl}\label{int:InterfacesAv2RedeploymentControl}
    \begin{description}[noitemsep,nolistsep]
      \item[Provided by:] \iconcomponent{}~\vpett{\nameref{comp:ComponentsAv2RedeploymentCoordinator}}
      \item[Required by:] \iconcomponent{}~\vpett{\nameref{comp:ComponentsAv2AdminPortal}} 
      \item[Operations:] ~
    \begin{itemize}[noitemsep,nolistsep,leftmargin=-.25cm]
      \item \textsf{\ref{data:DatatypesRedeploymentStatus} getRedeploymentStatus(\ref{data:DatatypesRedeploymentJobId} id)}
        \begin{itemize}[noitemsep,nolistsep]
           \item Effect: Returns the current status of the redeployment or rollback job identified by the given job identifier. Called periodically by AdminPortal to track job progress. If the job identifier does not exist a NoSuchJobException is thrown. 
           \item Returns: The current \ref{data:DatatypesRedeploymentStatus} of the job: IN\_PROGRESS if still running, COMPLETED if successful, FAILED if the deployment failed, or ROLLED\_BACK if a rollback completed successfully.
           \item Sequence Diagrams: None
        \end{itemize}
      \item \textsf{\ref{data:DatatypesRedeploymentJob} triggerRedeployment(string subsystemId)}
        \begin{itemize}[noitemsep,nolistsep]
           \item Effect: Initiates a redeployment of the specified subsystem to its latest available version. Called by AdminPortal when the system administrator decides to redeploy a failing component. The RedeploymentCoordinator begins the redeployment immediately and the operation returns as soon as the job is created. The administrator must poll getRedeploymentStatus to track progress. The entire redeployment must complete within 5 minutes per Av2 response measures. 
           \item Returns: A \ref{data:DatatypesRedeploymentJob} object containing the assigned job identifier, the target subsystem, the start timestamp, and the initial status of IN\_PROGRESS.
           \item Sequence Diagrams: None
        \end{itemize}
      \item \textsf{\ref{data:DatatypesRedeploymentJob} triggerRollback(string subsystemId, string targetVersion)}
        \begin{itemize}[noitemsep,nolistsep]
           \item Effect: Initiates a rollback of the specified subsystem to the given previously working version. Called by AdminPortal when redeployment fails or when the administrator decides to revert to a stable state. Like triggerRedeployment, the operation returns immediately with a job handle and the administrator polls for progress. Must complete within 5 minutes per Av2 response measures. 
           \item Returns: A \ref{data:DatatypesRedeploymentJob} object containing the assigned job identifier, the target subsystem, the target version, the start timestamp, and the initial status of IN\_PROGRESS.
           \item Sequence Diagrams: None
        \end{itemize}
    \end{itemize}
      \item[Diagrams:] None
    \end{description}

  %%%%%%%% Av2RetryDispatch
  \subsection{Av2RetryDispatch}\label{int:InterfacesAv2RetryDispatch}
    \begin{description}[noitemsep,nolistsep]
      \item[Provided by:] \iconcomponent{}~\vpett{\nameref{mod:ComponentsAv2CommunicationGatewayAv2RetryCoordinator}}
      \item[Required by:] \iconcomponent{}~\vpett{\nameref{mod:ComponentsAv2CommunicationGatewayAv2RoutingManager}} 
           \item[Description]: Published by the RetryCoordinator. Provides retry-ready messages to the RoutingManager for re-delivery. Only used inside the CommunicationGateway.
      \item[Operations:] ~
    \begin{itemize}[noitemsep,nolistsep,leftmargin=-.25cm]
      \item \textsf{void cancelRetry(\ref{data:DatatypesGatewayId} gatewayId)}
        \begin{itemize}[noitemsep,nolistsep]
           \item Effect: Cancels all pending retry attempts for the specified gateway. Called by RetryCoordinator when delivery is confirmed via BrokerCommitNotification, or when the failure has persisted long enough to require escalation to MonitoringService rather than further retrying. 
           \item Sequence Diagrams: None
        \end{itemize}
      \item \textsf{\ref{data:DatatypesRetryState} getRetryState(\ref{data:DatatypesGatewayId} gatewayId)}
        \begin{itemize}[noitemsep,nolistsep]
           \item Effect: Returns the current retry state for the specified gateway including the number of attempts made, the next scheduled retry time, and the current backoff interval. Used by RetryCoordinator to determine whether to continue retrying or escalate. 
           \item Returns: A \ref{data:DatatypesRetryState} object containing the gateway identifier, current attempt number, timestamp of the next scheduled retry, and the current backoff interval duration.
           \item Sequence Diagrams: None
        \end{itemize}
      \item \textsf{void scheduleRetry(\ref{data:DatatypesGatewayId} gatewayId, int attemt)}
        \begin{itemize}[noitemsep,nolistsep]
           \item Effect: Schedules a retry delivery attempt for the specified gateway. The attemptNumber determines the backoff interval — each successive attempt doubles the wait time before the next retry. Called by RetryCoordinator after receiving a connection failure from FailureDetector. RoutingManager picks up the retry signal and reattempts delivery to the MessageBroker. 
           \item Sequence Diagrams: None
        \end{itemize}
    \end{itemize}
      \item[Diagrams:] None
    \end{description}

  %%%%%%%% Av2RoutFailureSignals
  \subsection{Av2RoutFailureSignals}\label{int:InterfacesAv2RoutFailureSignals}
    \begin{description}[noitemsep,nolistsep]
      \item[Provided by:] \iconcomponent{}~\vpett{\nameref{mod:ComponentsAv2CommunicationGatewayAv2RoutingManager}}
      \item[Required by:] \iconcomponent{}~\vpett{\nameref{mod:ComponentsAv2CommunicationGatewayAv2FailureDetector}} 
      \item[Operations:] ~
% no operations
      \item[Diagrams:] None
    \end{description}

  %%%%%%%% Av2SensorDataAcquisition
  \subsection{Av2SensorDataAcquisition}\label{int:InterfacesAv2SensorDataAcquisition}
    \begin{description}[noitemsep,nolistsep]
      \item[Provided by:] None
      \item[Required by:] \iconcomponent{}~\vpett{\nameref{mod:ComponentsAv2PatientGatewayAv2SensorCommunicator}} 
           \item[Description]: Published by the external monitoring device. Provides the sensor data to the PatientGatewayApplication. Consumed by the SensorCommunicator
      \item[Operations:] ~
    \begin{itemize}[noitemsep,nolistsep,leftmargin=-.25cm]
      \item \textsf{boolean isSensorAvailable()}
        \begin{itemize}[noitemsep,nolistsep]
           \item Effect: Returns true if the monitoring device is connected and producing readings. Used by SensorCommunicator to detect sensor disconnection and alert PatientAlertService accordingly. 
           \item Returns: True if the monitoring device is connected and producing readings. False if the device is disconnected or not responding.
           \item Sequence Diagrams: None
        \end{itemize}
      \item \textsf{\ref{data:DatatypesSensorDataPackage} readLatestReading()}
        \begin{itemize}[noitemsep,nolistsep]
           \item Effect: Returns the most recent sensor data package produced by the monitoring device, containing one or more SensorDataElements covering all active sensor types (e.g. heart rate, blood pressure, ECG). Called periodically by SensorCommunicator at the configured transmission rate. 
           \item Returns: sensor data package
           \item Sequence Diagrams: None
        \end{itemize}
    \end{itemize}
      \item[Diagrams:] None
    \end{description}

  %%%%%%%% Av2ServerConnectionStatus
  \subsection{Av2ServerConnectionStatus}\label{int:InterfacesAv2ServerConnectionStatus}
    \begin{description}[noitemsep,nolistsep]
      \item[Provided by:] \iconcomponent{}~\vpett{\nameref{mod:ComponentsAv2PatientGatewayAv2ServerConnectivityMonitor}}
      \item[Required by:] \iconcomponent{}~\vpett{\nameref{mod:ComponentsAv2PatientGatewayAv2DegradeModeManager}}, \iconcomponent{}~\vpett{\nameref{mod:ComponentsAv2PatientGatewayAv2RetryAndSynchronizationService}} 
           \item[Description]: Published by the ServerConnectivityMonitor. Reports whether the PMS backend is currently reachable from the patient gateway and the timestamp of the last successful contact. Consumed by the DegradeModeManager to decide when to enter or exit degraded mode.
      \item[Operations:] ~
    \begin{itemize}[noitemsep,nolistsep,leftmargin=-.25cm]
      \item \textsf{\ref{data:DatatypesTimestamp} getLastSuccesfulContactTime()}
        \begin{itemize}[noitemsep,nolistsep]
           \item Effect: Returns the timestamp of the most recent successful heartbeat probe response from the PMS backend. Used by DegradeModeManager and RetryAndSynchronizationService to calculate how long the gateway has been operating without backend connectivity. 
           \item Returns: The \ref{data:DatatypesTimestamp} of the last successful contact with the PMS backend. If no successful contact has ever been made, returns null.
           \item Sequence Diagrams: None
        \end{itemize}
      \item \textsf{boolean isServerReachable()}
        \begin{itemize}[noitemsep,nolistsep]
           \item Effect: Returns whether the PMS backend is currently reachable from the patient gateway based on the most recent heartbeat probe sent by ServerConnectivityMonitor. Used by DegradeModeManager to decide whether to enter or exit degraded mode. 
           \item Returns: True if the most recent heartbeat probe received a response from the backend within the expected timeout. False if the backend did not respond, indicating the primary channel or backend is unavailable.
           \item Sequence Diagrams: None
        \end{itemize}
    \end{itemize}
      \item[Diagrams:] None
    \end{description}

  %%%%%%%% Av2SLAReporting
  \subsection{Av2SLAReporting}\label{int:InterfacesAv2SLAReporting}
    \begin{description}[noitemsep,nolistsep]
      \item[Provided by:] \iconcomponent{}~\vpett{\nameref{comp:ComponentsAv2MonitoringService}}, \iconcomponent{}~\vpett{\nameref{mod:ComponentsAv2MonitoringServiceAv2SLAMonitor}}
      \item[Required by:] \iconcomponent{}~\vpett{\nameref{comp:ComponentsAv2AdminPortal}} 
      \item[Operations:] ~
    \begin{itemize}[noitemsep,nolistsep,leftmargin=-.25cm]
      \item \textsf{{\colorbox{red!30}{\underline{SLAReport}}} getSLAReport(\ref{data:Datatypesorghl7fhirr5modelPeriod} period)}
        \begin{itemize}[noitemsep,nolistsep]
           \item Effect: Returns a compiled SLA compliance report for the given time period covering all monitored gateways. The report includes rolling availability percentage, mobile network coverage ratio, and average bandwidth measurements, compared against the SLA thresholds of 99.5\% availability, 80\% coverage, and 64Kb/s bandwidth. Used by AdminPortal to display compliance status to the system administrator.
 
           \item Returns: An SLAReport containing per-gateway and aggregate SLA metrics for the requested period, including any threshold violations detected during that period.
           \item Sequence Diagrams: None
        \end{itemize}
    \end{itemize}
      \item[Diagrams:] None
    \end{description}

  %%%%%%%% Av2ValidatedTransmission
  \subsection{Av2ValidatedTransmission}\label{int:InterfacesAv2ValidatedTransmission}
    \begin{description}[noitemsep,nolistsep]
      \item[Provided by:] \iconcomponent{}~\vpett{\nameref{mod:ComponentsAv2CommunicationGatewayAv2AuthenticationValidator}}
      \item[Required by:] \iconcomponent{}~\vpett{\nameref{mod:ComponentsAv2CommunicationGatewayAv2RoutingManager}} 
           \item[Description]: Carries transmissions that have passed authentication validation from the AuthenticationValidator to the RoutingManager. Only used inside the CommunicationGateway.
      \item[Operations:] ~
    \begin{itemize}[noitemsep,nolistsep,leftmargin=-.25cm]
      \item \textsf{void forwardValidatedTransmission(\ref{data:DatatypesGatewayId} gatewayId, {\colorbox{red!30}{\underline{ParsedTransmission}}} transmission, \ref{data:DatatypesTimestamp} recievedAt)}
        \begin{itemize}[noitemsep,nolistsep]
           \item Effect: Forwards a transmission that has passed authentication validation from AuthenticationValidator to RoutingManager for downstream dispatch. Only transmissions from registered gateways with valid credentials reach this interface. The RoutingManager uses the transmission type to decide whether to route it as sensor data or as an emergency notification. 
           \item Sequence Diagrams: None
        \end{itemize}
    \end{itemize}
      \item[Diagrams:] None
    \end{description}

  %%%%%%%% ClinicalJobMgmt
  \subsection{ClinicalJobMgmt}\label{int:InterfacesClinicalJobMgmt}
    \begin{description}[noitemsep,nolistsep]
      \item[Provided by:] \iconcomponent{}~\vpett{\nameref{comp:ComponentseHealthPlatformClinicalJobCreator}}
      \item[Required by:] \iconcomponent{}~\vpett{\nameref{comp:ComponentseHealthPlatformRiskEstimationScheduler}} 
      \item[Operations:] ~
    \begin{itemize}[noitemsep,nolistsep,leftmargin=-.25cm]
      \item \textsf{List\textless{}Tuple\textless{}\ref{data:DatatypesClinicalModelJob}, Double\textgreater{}\textgreater{} createClinicalModelJobsForPatient(\ref{data:DatatypesPatientId} patientId, Tuple\textless{}\ref{data:DatatypesSensorDataPackage}, \ref{data:DatatypesTimestamp}\textgreater{} sensorData)}
        \begin{itemize}[noitemsep,nolistsep]
           \item Effect: Construct the list of \ref{data:DatatypesClinicalModelJob}s for a specified patient. These jobs are populated with the necessary data for their execution by the \vpett{\nameref{comp:ComponentseHealthPlatformRiskEstimationProcessor}}. 
           \item Returns: A list of tuples containing the \ref{data:DatatypesClinicalModelJob} and its weight.
           \item Sequence Diagrams: \cref{diag:Interaction:Riskestimationprocess}
        \end{itemize}
    \end{itemize}
      \item[Diagrams:] None
    \end{description}

  %%%%%%%% ClinicalModelCacheMgmt
  \subsection{ClinicalModelCacheMgmt}\label{int:InterfacesClinicalModelCacheMgmt}
    \begin{description}[noitemsep,nolistsep]
      \item[Provided by:] \iconcomponent{}~\vpett{\nameref{comp:ComponentseHealthPlatformClinicalModelCache}}
      \item[Required by:] \iconcomponent{}~\vpett{\nameref{mod:ComponentsP1PhysicianCommandServiceP1ConfigurationHandler}}, \iconcomponent{}~\vpett{\nameref{comp:ComponentsP1PhysicianCommandService}} 
      \item[Operations:] ~
    \begin{itemize}[noitemsep,nolistsep,leftmargin=-.25cm]
      \item \textsf{void invalidateCacheEntries(\ref{data:DatatypesPatientId} patientId)}
        \begin{itemize}[noitemsep,nolistsep]
           \item Effect: Invalidate all cached items for the patient identified by patientId. If the cache holds nothing for this patient, nothing changes. After invalidation, the next request leads to a fresh fetch from the database and a fresh cache population. 
           \item Sequence Diagrams: None
        \end{itemize}
    \end{itemize}
      \item[Diagrams:] None
    \end{description}

  %%%%%%%% ClinicalModelStorage
  \subsection{ClinicalModelStorage}\label{int:InterfacesClinicalModelStorage}
    \begin{description}[noitemsep,nolistsep]
      \item[Provided by:] \iconcomponent{}~\vpett{\nameref{comp:ComponentseHealthPlatformClinicalModelDB}}
      \item[Required by:] \iconcomponent{}~\vpett{\nameref{comp:ComponentseHealthPlatformClinicalModelCache}} 
      \item[Operations:] ~
    \begin{itemize}[noitemsep,nolistsep,leftmargin=-.25cm]
      \item \textsf{Map\textless{}\ref{data:DatatypesClinicalModel}, Tuple\textless{}\ref{data:DatatypesClinicalModelConfiguration}, Double\textgreater{}\textgreater{} getAssignedClinicalModelsAndConfigForPatient(\ref{data:DatatypesPatientId} patientId) throws \ref{ex:ExceptionsNoSuchPatientException}}
        \begin{itemize}[noitemsep,nolistsep]
           \item Effect: Fetch and return the ids of the \ref{data:DatatypesClinicalModel}s assigned to the patient identified by patientId and their configurations for this patient. The configurations are returned as a \ref{data:DatatypesClinicalModelConfiguration} containing a map of configuration parameters and their value. 
           \item Returns: A map with for every applicable \ref{data:DatatypesClinicalModel}, a tuple containing the corresponding \ref{data:DatatypesClinicalModelConfiguration} for the specified patient and the corresponding weight factors used for combining the results.
           \item Sequence Diagrams: \cref{diag:Interaction:Riskestimationprocess}
        \end{itemize}
    \end{itemize}
      \item[Diagrams:] None
    \end{description}

  %%%%%%%% DataBuffer
  \subsection{DataBuffer}\label{int:InterfacesDataBuffer}
    \begin{description}[noitemsep,nolistsep]
      \item[Provided by:] \iconcomponent{}~\vpett{\nameref{mod:ComponentsAv2PatientGatewayAv2SensorCommunicator}}
      \item[Required by:] \iconcomponent{}~\vpett{\nameref{mod:ComponentsAv2PatientGatewayAv2SensorDataTransmitter}} 
           \item[Description]: Published by the SensorCommunicator. Provides a buffer of acquired sensor readings awaiting transmission. Consumed by the SensorDataTransmitter
      \item[Operations:] ~
    \begin{itemize}[noitemsep,nolistsep,leftmargin=-.25cm]
      \item \textsf{void bufferReading(\ref{data:DatatypesSensorDataPackage} reading, \ref{data:DatatypesTimestamp} time)}
        \begin{itemize}[noitemsep,nolistsep]
           \item Effect: Adds the given sensor data package and its acquisition timestamp to the in-memory buffer maintained by SensorCommunicator. Called each time a new reading is acquired from the wearable sensor hardware. If the buffer is full, the oldest entry is overwritten to ensure the most recent data is always preserved. 
           \item Sequence Diagrams: None
        \end{itemize}
      \item \textsf{void clearDeliveredReadings(\ref{data:DatatypesTimestamp} upTo)}
        \begin{itemize}[noitemsep,nolistsep]
           \item Effect: Removes all buffered readings with an acquisition timestamp at or before the given timestamp from the buffer. Called by SensorDataTransmitter after the CommunicationGateway confirms receipt of a batch of readings via TransmissionAck. 
           \item Sequence Diagrams: None
        \end{itemize}
      \item \textsf{List\textless{}Tuple\textless{}\ref{data:DatatypesSensorDataPackage}, \ref{data:DatatypesTimestamp}\textgreater{}\textgreater{} getBufferedReading()}
        \begin{itemize}[noitemsep,nolistsep]
           \item Effect: Returns all sensor data packages currently held in the buffer awaiting transmission. Called by SensorDataTransmitter to retrieve readings for delivery to the CommunicationGateway. 
           \item Returns: An ordered list of tuples each containing a \ref{data:DatatypesSensorDataPackage} and the timestamp at which it was acquired. Ordered from oldest to newest so that SensorDataTransmitter delivers data in chronological order.
           \item Sequence Diagrams: None
        \end{itemize}
      \item \textsf{boolean isBufferFull()}
        \begin{itemize}[noitemsep,nolistsep]
           \item Effect: Checks whether the buffer has reached its maximum capacity.
 
           \item Returns: True if the buffer is at maximum capacity and cannot accept new readings without overwriting existing ones. False if space remains available.
           \item Sequence Diagrams: None
        \end{itemize}
    \end{itemize}
      \item[Diagrams:] None
    \end{description}

  %%%%%%%% FetchClinicalModels
  \subsection{FetchClinicalModels}\label{int:InterfacesFetchClinicalModels}
    \begin{description}[noitemsep,nolistsep]
      \item[Provided by:] \iconcomponent{}~\vpett{\nameref{comp:ComponentseHealthPlatformClinicalModelCache}}
      \item[Required by:] \iconcomponent{}~\vpett{\nameref{comp:ComponentseHealthPlatformClinicalJobCreator}}, \iconcomponent{}~\vpett{\nameref{comp:ComponentseHealthPlatformRiskEstimationScheduler}} 
      \item[Operations:] ~
    \begin{itemize}[noitemsep,nolistsep,leftmargin=-.25cm]
      \item \textsf{Map\textless{}\ref{data:DatatypesClinicalModel}, Tuple\textless{}\ref{data:DatatypesClinicalModelConfiguration}, Double\textgreater{}\textgreater{} getAssignedClinicalModelsAndConfigForPatient(\ref{data:DatatypesPatientId} patientId)}
        \begin{itemize}[noitemsep,nolistsep]
           \item Effect: The \vpett{\nameref{comp:ComponentseHealthPlatformClinicalModelCache}} will return the clinical models assigned to the patient identified by patientId and their configurations for this patient. The configurations are
a map of configuration parameters and their value. If the cache contains data for the patient with given id, the data from the cache is returned. If not, the \vpett{\nameref{comp:ComponentseHealthPlatformClinicalModelCache}} will
fetch all clinical models assigned to the patient with given patientId and their configurations for this patient, store these in the cache and return them. 
           \item Returns: A map with for every applicable \ref{data:DatatypesClinicalModel}, a tuple containing the corresponding \ref{data:DatatypesClinicalModelConfiguration} for the specified patient and the corresponding weight factors used for combining the results.
           \item Sequence Diagrams: \cref{diag:Interaction:Riskestimationprocess}
        \end{itemize}
    \end{itemize}
      \item[Diagrams:] None
    \end{description}

  %%%%%%%% FetchJobs
  \subsection{FetchJobs}\label{int:InterfacesFetchJobs}
    \begin{description}[noitemsep,nolistsep]
      \item[Provided by:] \iconcomponent{}~\vpett{\nameref{comp:ComponentseHealthPlatformRiskEstimationScheduler}}
      \item[Required by:] \iconcomponent{}~\vpett{\nameref{mod:ComponentseHealthPlatformRiskEstimationProcessorJobProcessor}}, \iconcomponent{}~\vpett{\nameref{comp:ComponentseHealthPlatformRiskEstimationProcessor}} 
      \item[Operations:] ~
    \begin{itemize}[noitemsep,nolistsep,leftmargin=-.25cm]
      \item \textsf{\ref{data:DatatypesClinicalModelJob} getNextJob()}
        \begin{itemize}[noitemsep,nolistsep]
           \item Effect: Returns the next clinical model computation job that must be performed (i.e., the first job in the queue). 
           \item Returns: The next clinical model computation job that must be performed.
           \item Sequence Diagrams: \cref{diag:Interaction:Computeclinicalmodelresult}
        \end{itemize}
    \end{itemize}
      \item[Diagrams:] None
    \end{description}

  %%%%%%%% FetchMLModel
  \subsection{FetchMLModel}\label{int:InterfacesFetchMLModel}
    \begin{description}[noitemsep,nolistsep]
      \item[Provided by:] \iconcomponent{}~\vpett{\nameref{mod:ComponentseHealthPlatformRiskEstimationProcessorMLModelRetrieval}}
      \item[Required by:] \iconcomponent{}~\vpett{\nameref{mod:ComponentseHealthPlatformRiskEstimationProcessorJobProcessor}} 
      \item[Operations:] ~
    \begin{itemize}[noitemsep,nolistsep,leftmargin=-.25cm]
      \item \textsf{\ref{data:DatatypesMLModel} fetchModel(String modelId) throws \ref{ex:ExceptionsNoSuchMLModelException}}
        \begin{itemize}[noitemsep,nolistsep]
           \item Effect: Fetch the \ref{data:DatatypesMLModel} corresponding with the provided modelId. 
           \item Returns: The requested \ref{data:DatatypesMLModel}.
           \item Sequence Diagrams: \cref{diag:Interaction:Computeclinicalmodelresult}
        \end{itemize}
    \end{itemize}
      \item[Diagrams:] None
    \end{description}

  %%%%%%%% healthAPI
  \subsection{healthAPI}\label{int:InterfaceshealthAPI}
    \begin{description}[noitemsep,nolistsep]
      \item[Provided by:] \iconcomponent{}~\vpett{\nameref{comp:ComponentsHIS}}
      \item[Required by:] \iconcomponent{}~\vpett{\nameref{comp:ComponentsP1HISAdapter}} 
           \item[Description]: This interface provides a reduced set of simplified operations that are sufficient in the context of the patient monitoring service that is being developed.
It is derived from the HL7 FHIR standard for healthcare interoperability, based on the HAPI FHIR implementation of the FHIR specification in Java. The relevant OpenAPIs are described at: \url{https://hapi.fhir.org/baseR4/swagger-ui/}, while the FHIR standard is available at: \url{https://hl7.org/fhir/}.
The descriptions from the data types used in this interface follow directly from the API documentation (\url{https://hapifhir.io/hapi-fhir/apidocs/hapi-fhir-structures-r5/org/hl7/fhir/r5/model/package-summary.html}).
      \item[Operations:] ~
    \begin{itemize}[noitemsep,nolistsep,leftmargin=-.25cm]
      \item \textsf{void deleteObservation(string id) throws \ref{ex:ExceptionsAuthorizationException2}, \ref{ex:ExceptionsAuthenticationException2}}
        \begin{itemize}[noitemsep,nolistsep]
           \item Effect: Delete an \ref{data:Datatypesorghl7fhirr5modelObservation} instance with the specified id. 
           \item Sequence Diagrams: None
        \end{itemize}
      \item \textsf{void deletePatient(string id) throws \ref{ex:ExceptionsAuthorizationException2}, \ref{ex:ExceptionsAuthenticationException2}}
        \begin{itemize}[noitemsep,nolistsep]
           \item Effect: Delete a \ref{data:Datatypesorghl7fhirr5modelPatient} instance with the specfied id. 
           \item Sequence Diagrams: None
        \end{itemize}
      \item \textsf{void deleteRiskAssessment(string id) throws \ref{ex:ExceptionsAuthorizationException2}, \ref{ex:ExceptionsAuthenticationException2}}
        \begin{itemize}[noitemsep,nolistsep]
           \item Effect: Delete a \ref{data:Datatypesorghl7fhirr5modelRiskAssessment} instance with the specified id. 
           \item Sequence Diagrams: None
        \end{itemize}
      \item \textsf{\ref{data:Datatypesorghl7fhirr5modelObservation} getObservation(string id) throws \ref{ex:ExceptionsAuthorizationException2}, \ref{ex:ExceptionsAuthenticationException2}}
        \begin{itemize}[noitemsep,nolistsep]
           \item Effect: Retrieve an observation instance 
           \item Returns: The \ref{data:Datatypesorghl7fhirr5modelObservation} with the specified identifier.
           \item Sequence Diagrams: None
        \end{itemize}
      \item \textsf{\ref{data:Datatypesorghl7fhirr5modelPatient} getPatient(string id) throws \ref{ex:ExceptionsAuthorizationException2}, \ref{ex:ExceptionsAuthenticationException2}}
        \begin{itemize}[noitemsep,nolistsep]
           \item Effect: Retrieve a patient record with the specified identifier. 
           \item Returns: The \ref{data:Datatypesorghl7fhirr5modelPatient} with the specified identifier.
           \item Sequence Diagrams: None
        \end{itemize}
      \item \textsf{\ref{data:Datatypesorghl7fhirr5modelRiskAssessment} getRiskAssessment(string id) throws \ref{ex:ExceptionsAuthorizationException2}}
        \begin{itemize}[noitemsep,nolistsep]
           \item Effect: Retrieve a \ref{data:Datatypesorghl7fhirr5modelRiskAssessment} instance with the specified id. 
           \item Returns: The \ref{data:Datatypesorghl7fhirr5modelRiskAssessment} with the specified identifier.
           \item Sequence Diagrams: None
        \end{itemize}
      \item \textsf{\ref{data:Datatypesorghl7fhirr5modelObservation} saveObservation(\ref{data:Datatypesorghl7fhirr5modelObservation} observation) throws \ref{ex:ExceptionsAuthorizationException2}, \ref{ex:ExceptionsAuthenticationException2}}
        \begin{itemize}[noitemsep,nolistsep]
           \item Effect: Create a new \ref{data:Datatypesorghl7fhirr5modelObservation} instance or in case an existing \ref{data:Datatypesorghl7fhirr5modelObservation} instance exists with the same identifier, update this \ref{data:Datatypesorghl7fhirr5modelObservation} instance. 
           \item Returns: The \ref{data:Datatypesorghl7fhirr5modelObservation} object.
           \item Sequence Diagrams: None
        \end{itemize}
      \item \textsf{\ref{data:Datatypesorghl7fhirr5modelPatient} savePatient(\ref{data:Datatypesorghl7fhirr5modelPatient} patient) throws \ref{ex:ExceptionsAuthorizationException2}}
        \begin{itemize}[noitemsep,nolistsep]
           \item Effect: Create a new \ref{data:Datatypesorghl7fhirr5modelPatient} instance or update an existing \ref{data:Datatypesorghl7fhirr5modelPatient} instance. 
           \item Returns: The updated \ref{data:Datatypesorghl7fhirr5modelPatient} object.
           \item Sequence Diagrams: None
        \end{itemize}
      \item \textsf{\ref{data:Datatypesorghl7fhirr5modelRiskAssessment} saveRiskAssessment(\ref{data:Datatypesorghl7fhirr5modelRiskAssessment} riskAssessment) throws \ref{ex:ExceptionsAuthorizationException2}, \ref{ex:ExceptionsAuthenticationException2}}
        \begin{itemize}[noitemsep,nolistsep]
           \item Effect: Create a new \ref{data:Datatypesorghl7fhirr5modelRiskAssessment} instance or update an existing \ref{data:Datatypesorghl7fhirr5modelRiskAssessment} instance. 
           \item Returns: The updated \ref{data:Datatypesorghl7fhirr5modelRiskAssessment} object.
           \item Sequence Diagrams: None
        \end{itemize}
      \item \textsf{List\textless{}\ref{data:Datatypesorghl7fhirr5modelObservation}\textgreater{} searchObservation(\ref{data:DatatypesQuery}\textless{}\ref{data:DatatypesDate}\textgreater{} date, \ref{data:DatatypesQuery}\textless{}String\textgreater{} dataAbsentReason, \ref{data:DatatypesQuery}\textless{}\ref{data:Datatypesorghl7fhirr5modelPatient}\textgreater{} subject, \ref{data:DatatypesQuery}\textless{}String\textgreater{} valueConcept, \ref{data:DatatypesQuery}\textless{}\ref{data:DatatypesDate}\textgreater{} valueDate, \ref{data:DatatypesQuery}\textless{}\ref{data:Datatypesorghl7fhirr5modelObservation}\textgreater{} derivedFrom, \ref{data:DatatypesQuery}\textless{}\ref{data:Datatypesorghl7fhirr5modelPatient}\textgreater{} patient, \ref{data:DatatypesQuery}\textless{}\ref{data:Datatypesorghl7fhirr5modelQuantity}\textgreater{} valueQuantity, \ref{data:DatatypesQuery}\textless{}String\textgreater{} identifier, \ref{data:DatatypesQuery}\textless{}String\textgreater{} performer, \ref{data:DatatypesQuery}\textless{}String\textgreater{} method, \ref{data:DatatypesQuery}\textless{}String\textgreater{} category, \ref{data:DatatypesQuery}\textless{}String\textgreater{} device, \ref{data:DatatypesQuery}\textless{}\ref{data:Datatypesorghl7fhirr5modelObservationStatus}\textgreater{} status) throws \ref{ex:ExceptionsAuthorizationException2}}
        \begin{itemize}[noitemsep,nolistsep]
           \item Effect: Search for \ref{data:Datatypesorghl7fhirr5modelObservation} instances matching the specified query criteria. 
           \item Returns: The list of \ref{data:Datatypesorghl7fhirr5modelObservation}s matching the specified search criteria.
           \item Sequence Diagrams: None
        \end{itemize}
      \item \textsf{List\textless{}\ref{data:Datatypesorghl7fhirr5modelPatient}\textgreater{} searchPatient(\ref{data:DatatypesQuery}\textless{}\ref{data:DatatypesDate}\textgreater{} birthDate, \ref{data:DatatypesQuery}\textless{}Boolean\textgreater{} deceased, \ref{data:DatatypesQuery}\textless{}String\textgreater{} addressState, \ref{data:DatatypesQuery}\textless{}String\textgreater{} administrativeGender, \ref{data:DatatypesQuery}\textless{}String\textgreater{} link, \ref{data:DatatypesQuery}\textless{}String\textgreater{} language, \ref{data:DatatypesQuery}\textless{}String\textgreater{} addressCountry, \ref{data:DatatypesQuery}\textless{}\ref{data:DatatypesDate}\textgreater{} deathDate, \ref{data:DatatypesQuery}\textless{}String\textgreater{} phonetic, \ref{data:DatatypesQuery}\textless{}String\textgreater{} telecom, \ref{data:DatatypesQuery}\textless{}String\textgreater{} addressCity, \ref{data:DatatypesQuery}\textless{}String\textgreater{} email, \ref{data:DatatypesQuery}\textless{}String\textgreater{} given, \ref{data:DatatypesQuery}\textless{}String\textgreater{} identifier, \ref{data:DatatypesQuery}\textless{}String\textgreater{} address, \ref{data:DatatypesQuery}\textless{}String\textgreater{} generalPractitioner, \ref{data:DatatypesQuery}\textless{}Boolean\textgreater{} active, \ref{data:DatatypesQuery}\textless{}String\textgreater{} addressPostalCode, \ref{data:DatatypesQuery}\textless{}String\textgreater{} phone, \ref{data:DatatypesQuery}\textless{}String\textgreater{} organizationCustodian, \ref{data:DatatypesQuery}\textless{}String\textgreater{} name, \ref{data:DatatypesQuery}\textless{}String\textgreater{} family) throws \ref{ex:ExceptionsAuthorizationException2}, \ref{ex:ExceptionsAuthenticationException2}}
        \begin{itemize}[noitemsep,nolistsep]
           \item Effect: Search for \ref{data:Datatypesorghl7fhirr5modelPatient} instances matching the specified query criteria. 
           \item Returns: The list of \ref{data:Datatypesorghl7fhirr5modelPatient}s matching the specified search criteria.
           \item Sequence Diagrams: None
        \end{itemize}
      \item \textsf{List\textless{}\ref{data:Datatypesorghl7fhirr5modelRiskAssessment}\textgreater{} searchRiskAssessment(\ref{data:DatatypesQuery}\textless{}\ref{data:DatatypesDate}\textgreater{} date, \ref{data:DatatypesQuery}\textless{}String\textgreater{} identifier, \ref{data:DatatypesQuery}\textless{}String\textgreater{} performer, \ref{data:DatatypesQuery}\textless{}String\textgreater{} method, \ref{data:DatatypesQuery}\textless{}\ref{data:Datatypesorghl7fhirr5modelRange}\textgreater{} probability, \ref{data:DatatypesQuery}\textless{}\ref{data:Datatypesorghl7fhirr5modelPatient}\textgreater{} subject, \ref{data:DatatypesQuery}\textless{}String\textgreater{} condition, \ref{data:DatatypesQuery}\textless{}\ref{data:Datatypesorghl7fhirr5modelPatient}\textgreater{} patient, \ref{data:DatatypesQuery}\textless{}\ref{data:DatatypesBigDecimal}\textgreater{} risk) throws \ref{ex:ExceptionsAuthenticationException2}, \ref{ex:ExceptionsAuthorizationException2}}
        \begin{itemize}[noitemsep,nolistsep]
           \item Effect: Search for a \ref{data:Datatypesorghl7fhirr5modelRiskAssessment} instance matching the provided query criteria. 
           \item Returns: The list of \ref{data:Datatypesorghl7fhirr5modelRiskAssessment}s matching the specified search criteria.
           \item Sequence Diagrams: None
        \end{itemize}
    \end{itemize}
      \item[Diagrams:] None
    \end{description}

  %%%%%%%% JobMgmt
  \subsection{JobMgmt}\label{int:InterfacesJobMgmt}
    \begin{description}[noitemsep,nolistsep]
      \item[Provided by:] \iconcomponent{}~\vpett{\nameref{comp:ComponentseHealthPlatformRiskEstimationCombiner}}
      \item[Required by:] \iconcomponent{}~\vpett{\nameref{comp:ComponentseHealthPlatformRiskEstimationScheduler}} 
      \item[Operations:] ~
    \begin{itemize}[noitemsep,nolistsep,leftmargin=-.25cm]
      \item \textsf{void addJobSet(\ref{data:DatatypesRiskEstimationId} id, List\textless{}Tuple\textless{}\ref{data:DatatypesClinicalModelJobId}, Double\textgreater{}\textgreater{} jobIds)}
        \begin{itemize}[noitemsep,nolistsep]
           \item Effect: Adds a set of identifiers for jobs belonging to a
single risk estimation identified by id. The partial
results of each clinical model computation job identified by an
element in jobIds have to be combined in order to find
the final result of the risk estimation as a whole. 
           \item Sequence Diagrams: \cref{diag:Interaction:Riskestimationprocess}
        \end{itemize}
    \end{itemize}
      \item[Diagrams:] None
    \end{description}

  %%%%%%%% KVStore
  \subsection{KVStore}\label{int:InterfacesKVStore}
    \begin{description}[noitemsep,nolistsep]
      \item[Provided by:] \iconcomponent{}~\vpett{\nameref{comp:ComponentsKVStorageService}}
      \item[Required by:] \iconcomponent{}~\vpett{\nameref{comp:ComponentseHealthPlatformRiskEstimationProcessor}}, \iconcomponent{}~\vpett{\nameref{mod:ComponentseHealthPlatformRiskEstimationProcessorSensorDataRiskAssessment}}, \iconcomponent{}~\vpett{\nameref{comp:ComponentseHealthPlatform}} 
      \item[Operations:] ~
    \begin{itemize}[noitemsep,nolistsep,leftmargin=-.25cm]
      \item \textsf{\ref{data:DatatypesObject} retrieve(String key) throws \ref{ex:ExceptionsIOException}}
        \begin{itemize}[noitemsep,nolistsep]
           \item Effect: Retrieve the object with the provided key from the Key-Value storage. 
           \item Returns: The object retrieved from storage.
           \item Sequence Diagrams: \cref{diag:Interaction:Computeclinicalmodelresult}
        \end{itemize}
      \item \textsf{void store(String key, \ref{data:DatatypesObject} value) throws \ref{ex:ExceptionsIOException}}
        \begin{itemize}[noitemsep,nolistsep]
           \item Effect: Store the provided object in the Key-Value storage under the provided key. 
           \item Sequence Diagrams: \cref{diag:Interaction:Computeclinicalmodelresult}
        \end{itemize}
    \end{itemize}
      \item[Diagrams:] None
    \end{description}

  %%%%%%%% LaunchRiskEstimation
  \subsection{LaunchRiskEstimation}\label{int:InterfacesLaunchRiskEstimation}
    \begin{description}[noitemsep,nolistsep]
      \item[Provided by:] \iconcomponent{}~\vpett{\nameref{comp:ComponentseHealthPlatformRiskEstimationScheduler}}
      \item[Required by:] \iconcomponent{}~\vpett{\nameref{mod:ComponentsP1PhysicianCommandServiceP1ConfigurationHandler}}, \iconcomponent{}~\vpett{\nameref{mod:ComponentsP1PhysicianCommandServiceP1OnDemandConsultationHandler}}, \iconcomponent{}~\vpett{\nameref{comp:ComponentsP1PatientDataService}}, \iconcomponent{}~\vpett{\nameref{comp:ComponentsP1PhysicianCommandService}}, \iconcomponent{}~\vpett{\nameref{mod:ComponentsP1PatientDataServiceP1SensorIngestModule}} 
      \item[Operations:] ~
    \begin{itemize}[noitemsep,nolistsep,leftmargin=-.25cm]
      \item \textsf{void launchRiskEstimation(\ref{data:DatatypesPatientId} patientId, \ref{data:DatatypesSensorDataPackage} newSensorData, \ref{data:DatatypesTimestamp} receivedAt, \ref{data:DatatypesPriority} priority, \ref{data:DatatypesCorrelationId} correlationId)}
        \begin{itemize}[noitemsep,nolistsep]
           \item Effect: Submit a risk-estimation job for the patient. When `newSensorData` is supplied (UC4, UC9), the job runs against that package. When omitted (UC7 reconfigure), the scheduler runs the job against the patient's most recent stored sensor data. The scheduler orders the job by `priority` (UC9 HIGH jumps ahead of scheduled jobs) and carries `correlationId` through to the resulting P1IRiskEvents event so subscribers can match it. Scheduling is otherwise unchanged: FIFO in normal mode, dynamic priority (earliest-deadline-first with tighter deadlines for high-risk patients) in overload mode. 
           \item Sequence Diagrams: \cref{diag:Interaction:Riskestimationprocess,diag:Interaction:Incomingsensordata}
        \end{itemize}
    \end{itemize}
      \item[Diagrams:] None
    \end{description}

  %%%%%%%% MLaaSAPI
  \subsection{MLaaSAPI}\label{int:InterfacesMLaaSAPI}
    \begin{description}[noitemsep,nolistsep]
      \item[Provided by:] \iconcomponent{}~\vpett{\nameref{mod:ModulesMLaaSLibrary}}
      \item[Required by:] \iconcomponent{}~\vpett{\nameref{mod:ComponentseHealthPlatformMLModelManagerMLModelUpdateProcessor}} 
           \item[Description]: This API provides a number of operations for managing the models deployed on the \vpett{\nameref{comp:ComponentsMLaaS}} cloud.
      \item[Operations:] ~
    \begin{itemize}[noitemsep,nolistsep,leftmargin=-.25cm]
      \item \textsf{\ref{data:DatatypesCredential} authenticate() throws \ref{ex:ExceptionsInvalidAPIKeyException}}
        \begin{itemize}[noitemsep,nolistsep]
           \item Effect: Construct a \ref{data:DatatypesCredential} object for authenticating requests with the \vpett{\nameref{comp:ComponentsMLaaS}} service provider. This operation uses the previously configured API key to construct the \ref{data:DatatypesCredential} object. If such an API key is not configured, an exception will be thrown. 
           \item Returns: Authentication token for subsequent requests.
           \item Sequence Diagrams: \cref{diag:Interaction:MLModelSynchronization,diag:Interaction:ProcessingMLModelupdates}
        \end{itemize}
      \item \textsf{void delete(\ref{data:DatatypesMLModel} model, \ref{data:DatatypesCredential} token) throws \ref{ex:ExceptionsAuthorizationException}}
        \begin{itemize}[noitemsep,nolistsep]
           \item Effect: Delete a model from the online \vpett{\nameref{comp:ComponentsMLaaS}} service. 
           \item Sequence Diagrams: None
        \end{itemize}
      \item \textsf{\ref{data:DatatypesMLModelDescription} getModelDescription(List\textless{}\ref{data:DatatypesMLModelDescription}\textgreater{} descriptions, String modelId)}
        \begin{itemize}[noitemsep,nolistsep]
           \item Effect: This method retrieves the correct \ref{data:DatatypesMLModelDescription} that corresponds with the \ref{data:DatatypesMLModel} identified by the provided modelId. 
           \item Returns: The \ref{data:DatatypesMLModelDescription} corresponding with the \ref{data:DatatypesMLModel} identified by the modelId.
           \item Sequence Diagrams: None
        \end{itemize}
      \item \textsf{String getModelId(\ref{data:DatatypesMLModelDescription} description)}
        \begin{itemize}[noitemsep,nolistsep]
           \item Effect: Extract the modelId from the \ref{data:DatatypesMLModelDescription}. 
           \item Returns: String with the modelId of the \ref{data:DatatypesMLModel} specified in the \ref{data:DatatypesMLModelDescription}.
           \item Sequence Diagrams: \cref{diag:Interaction:MLModelSynchronization}
        \end{itemize}
      \item \textsf{void register(\ref{data:DatatypesMLModel} model, \ref{data:DatatypesCredential} token) throws \ref{ex:ExceptionsAuthorizationException}}
        \begin{itemize}[noitemsep,nolistsep]
           \item Effect: Register a new model with the \vpett{\nameref{comp:ComponentsMLaaS}} service. 
           \item Sequence Diagrams: None
        \end{itemize}
      \item \textsf{void setAPIKey(String apiKey) throws \ref{ex:ExceptionsAuthenticationException}}
        \begin{itemize}[noitemsep,nolistsep]
           \item Effect: Set the credentials for authenticating with the ML-as-a-Service provider. 
           \item Sequence Diagrams: None
        \end{itemize}
    \end{itemize}
      \item[Diagrams:] None
    \end{description}

  %%%%%%%% MLaaSService
  \subsection{MLaaSService}\label{int:InterfacesMLaaSService}
    \begin{description}[noitemsep,nolistsep]
      \item[Provided by:] \iconcomponent{}~\vpett{\nameref{comp:ComponentsMLaaS}}
      \item[Required by:] \iconcomponent{}~\vpett{\nameref{comp:ComponentseHealthPlatformMLModelManager}}, \iconcomponent{}~\vpett{\nameref{mod:ComponentseHealthPlatformMLModelManagerMLModelUpdateProcessor}}, \iconcomponent{}~\vpett{\nameref{comp:ComponentseHealthPlatform}} 
           \item[Description]: This interface offiers operations for interacting with \vpett{\nameref{comp:ComponentsMLaaS}} service providers.
      \item[Operations:] ~
    \begin{itemize}[noitemsep,nolistsep,leftmargin=-.25cm]
      \item \textsf{List\textless{}\ref{data:DatatypesMLModelUpdate}\textgreater{} fetchModelUpdates(\ref{data:DatatypesCredential} credential, \ref{data:DatatypesMLModelDescription} model)}
        \begin{itemize}[noitemsep,nolistsep]
           \item Effect: Fetch a list of model updates to apply to the local \ref{data:DatatypesMLModel}. 
           \item Returns: The list of MLModelUpdates to apply locally.
           \item Sequence Diagrams: \cref{diag:Interaction:ProcessingMLModelupdates}
        \end{itemize}
      \item \textsf{List\textless{}\ref{data:DatatypesMLModelDescription}\textgreater{} getAvailableModels(\ref{data:DatatypesCredential} credential)}
        \begin{itemize}[noitemsep,nolistsep]
           \item Effect: Get the list of available ML models to use for making predictions on incoming sensor data. 
           \item Returns: The list of available MLModels.
           \item Sequence Diagrams: \cref{diag:Interaction:MLModelSynchronization}
        \end{itemize}
      \item \textsf{void pushModelUpdates(\ref{data:DatatypesCredential} credential, \ref{data:DatatypesMLModelDescription} model, List\textless{}\ref{data:DatatypesMLModelUpdate}\textgreater{} updates)}
        \begin{itemize}[noitemsep,nolistsep]
           \item Effect: Send a list of MLModelUpdates to the online service. 
           \item Sequence Diagrams: None
        \end{itemize}
      \item \textsf{\ref{data:DatatypesMLModel} retrieveModel(\ref{data:DatatypesCredential} credential, \ref{data:DatatypesMLModelDescription} model)}
        \begin{itemize}[noitemsep,nolistsep]
           \item Effect: Retrieve a specific \ref{data:DatatypesMLModel} corresponding with the provided \ref{data:DatatypesMLModelDescription}. 
           \item Returns: The requested \ref{data:DatatypesMLModel}.
           \item Sequence Diagrams: \cref{diag:Interaction:MLModelSynchronization}
        \end{itemize}
    \end{itemize}
      \item[Diagrams:] None
    \end{description}

  %%%%%%%% MLAPI
  \subsection{MLAPI}\label{int:InterfacesMLAPI}
    \begin{description}[noitemsep,nolistsep]
      \item[Provided by:] \iconcomponent{}~\vpett{\nameref{mod:ModulesMLLibrary}}
      \item[Required by:] \iconcomponent{}~\vpett{\nameref{mod:ComponentseHealthPlatformMLModelManagerModelUpdater}}, \iconcomponent{}~\vpett{\nameref{mod:ComponentseHealthPlatformRiskEstimationProcessorSensorDataRiskAssessment}} 
           \item[Description]: This is the API provided by the \vpett{\nameref{mod:ModulesMLLibrary}} to use \ref{data:DatatypesMLModel}s locally for making predictions and processing updates to these models.
      \item[Operations:] ~
    \begin{itemize}[noitemsep,nolistsep,leftmargin=-.25cm]
      \item \textsf{\ref{data:DatatypesMLModel} applyUpdate(\ref{data:DatatypesMLModel} model, \ref{data:DatatypesMLModelUpdate} update)}
        \begin{itemize}[noitemsep,nolistsep]
           \item Effect: Applies the provided \ref{data:DatatypesMLModelUpdate} on the provided \ref{data:DatatypesMLModel} and returns the resulting \ref{data:DatatypesMLModel}. 
           \item Returns: The resulting \ref{data:DatatypesMLModel} after applying the requested \ref{data:DatatypesMLModelUpdate}.
           \item Sequence Diagrams: \cref{diag:Interaction:ProcessingMLModelupdates}
        \end{itemize}
      \item \textsf{\ref{data:DatatypesPrediction} predict(\ref{data:DatatypesSensorDataPackage} sensorData)}
        \begin{itemize}[noitemsep,nolistsep]
           \item Effect: Make a prediction with the active \ref{data:DatatypesMLModel} on the provided sensordata. 
           \item Returns: The prediction based on the provided sensor data.
           \item Sequence Diagrams: \cref{diag:Interaction:Computeclinicalmodelresult}
        \end{itemize}
      \item \textsf{void verify(\ref{data:DatatypesMLModel} model) throws \ref{ex:ExceptionsInvalidMLModelException}}
        \begin{itemize}[noitemsep,nolistsep]
           \item Effect: Verifies that a provided \ref{data:DatatypesMLModel} is working correctly after applying \ref{data:DatatypesMLModelUpdate}s. 
           \item Sequence Diagrams: \cref{diag:Interaction:ProcessingMLModelupdates}
        \end{itemize}
    \end{itemize}
      \item[Diagrams:] None
    \end{description}

  %%%%%%%% MLModelMgmt
  \subsection{MLModelMgmt}\label{int:InterfacesMLModelMgmt}
    \begin{description}[noitemsep,nolistsep]
      \item[Provided by:] \iconcomponent{}~\vpett{\nameref{comp:ComponentseHealthPlatformMLModelManager}}, \iconcomponent{}~\vpett{\nameref{mod:ComponentseHealthPlatformMLModelManagerMLModelStorageManager}}
      \item[Required by:] \iconcomponent{}~\vpett{\nameref{mod:ComponentseHealthPlatformRiskEstimationProcessorMLModelRetrieval}}, \iconcomponent{}~\vpett{\nameref{comp:ComponentseHealthPlatformRiskEstimationProcessor}} 
           \item[Description]: The interface provides operations for managing different \ref{data:DatatypesMLModel}s locally, querying available models for making predictions on newly incoming sensor data.
      \item[Operations:] ~
    \begin{itemize}[noitemsep,nolistsep,leftmargin=-.25cm]
      \item \textsf{\ref{data:DatatypesMLModel} fetchModel(String modelId) throws \ref{ex:ExceptionsNoSuchMLModelException}}
        \begin{itemize}[noitemsep,nolistsep]
           \item Effect: Fetch the \ref{data:DatatypesMLModel} corresponding with the provided modelId. 
           \item Returns: The requested \ref{data:DatatypesMLModel}.
           \item Sequence Diagrams: \cref{diag:Interaction:Computeclinicalmodelresult}
        \end{itemize}
      \item \textsf{List\textless{}String\textgreater{} getAvailableModels()}
        \begin{itemize}[noitemsep,nolistsep]
           \item Effect: Get the list of available \ref{data:DatatypesMLModel}s that can be used for making \ref{data:DatatypesPrediction}s locally. 
           \item Returns: The list of identifiers of the available \ref{data:DatatypesMLModel}s that can be used.
           \item Sequence Diagrams: None
        \end{itemize}
      \item \textsf{void synchronizeModels()}
        \begin{itemize}[noitemsep,nolistsep]
           \item Effect: Synchronize the list of locally-stored \ref{data:DatatypesMLModel}s by checking for newly-available \ref{data:DatatypesMLModel}s that are not stored in the local \vpett{\nameref{mod:ComponentseHealthPlatformMLModelManagerMLModelRepository}} and retrieving these \ref{data:DatatypesMLModel}s from the \vpett{\nameref{comp:ComponentsMLaaS}} provider. 
           \item Sequence Diagrams: \cref{diag:Interaction:MLModelSynchronization}
        \end{itemize}
      \item \textsf{void updateModel(String modelId) throws \ref{ex:ExceptionsNoSuchMLModelException}}
        \begin{itemize}[noitemsep,nolistsep]
           \item Effect: Update the model with the specified id by fetching \ref{data:DatatypesMLModelUpdate}s and applying them to the model. 
           \item Sequence Diagrams: \cref{diag:Interaction:ProcessingMLModelupdates}
        \end{itemize}
    \end{itemize}
      \item[Diagrams:] None
    \end{description}

  %%%%%%%% MLModelStorage
  \subsection{MLModelStorage}\label{int:InterfacesMLModelStorage}
    \begin{description}[noitemsep,nolistsep]
      \item[Provided by:] \iconcomponent{}~\vpett{\nameref{mod:ComponentseHealthPlatformMLModelManagerMLModelRepository}}
      \item[Required by:] \iconcomponent{}~\vpett{\nameref{mod:ComponentseHealthPlatformMLModelManagerMLModelStorageManager}}, \iconcomponent{}~\vpett{\nameref{mod:ComponentseHealthPlatformMLModelManagerMLModelUpdateProcessor}} 
      \item[Operations:] ~
    \begin{itemize}[noitemsep,nolistsep,leftmargin=-.25cm]
      \item \textsf{boolean isAvailable(String modelId)}
        \begin{itemize}[noitemsep,nolistsep]
           \item Effect: Check whether an \ref{data:DatatypesMLModel} (identified by the modelId) is available in the local \vpett{\nameref{mod:ComponentseHealthPlatformMLModelManagerMLModelRepository}}. 
           \item Returns: Boolean indicating if the specified model is present in the repository.
           \item Sequence Diagrams: \cref{diag:Interaction:MLModelSynchronization,diag:Interaction:ProcessingMLModelupdates}
        \end{itemize}
      \item \textsf{\ref{data:DatatypesMLModel} retrieveModel(String modelId)}
        \begin{itemize}[noitemsep,nolistsep]
           \item Effect: Retrieve the \ref{data:DatatypesMLModel} with the specified id. 
           \item Returns: The requested \ref{data:DatatypesMLModel}.
           \item Sequence Diagrams: \cref{diag:Interaction:ProcessingMLModelupdates}
        \end{itemize}
      \item \textsf{void storeModel(\ref{data:DatatypesMLModel} model)}
        \begin{itemize}[noitemsep,nolistsep]
           \item Effect: Store an \ref{data:DatatypesMLModel} in the local \vpett{\nameref{mod:ComponentseHealthPlatformMLModelManagerMLModelRepository}} 
           \item Sequence Diagrams: \cref{diag:Interaction:MLModelSynchronization,diag:Interaction:ProcessingMLModelupdates}
        \end{itemize}
    \end{itemize}
      \item[Diagrams:] None
    \end{description}

  %%%%%%%% MLModelSync
  \subsection{MLModelSync}\label{int:InterfacesMLModelSync}
    \begin{description}[noitemsep,nolistsep]
      \item[Provided by:] \iconcomponent{}~\vpett{\nameref{mod:ComponentseHealthPlatformMLModelManagerMLModelUpdateProcessor}}
      \item[Required by:] \iconcomponent{}~\vpett{\nameref{mod:ComponentseHealthPlatformMLModelManagerMLModelStorageManager}} 
      \item[Operations:] ~
    \begin{itemize}[noitemsep,nolistsep,leftmargin=-.25cm]
      \item \textsf{void fetchNewMLModels()}
        \begin{itemize}[noitemsep,nolistsep]
           \item Effect: Fetches newly available \ref{data:DatatypesMLModel}s and store them in the \vpett{\nameref{mod:ComponentseHealthPlatformMLModelManagerMLModelRepository}}. 
           \item Sequence Diagrams: \cref{diag:Interaction:MLModelSynchronization}
        \end{itemize}
      \item \textsf{void processMLModelUpdates(String modelId) throws \ref{ex:ExceptionsNoSuchMLModelException}}
        \begin{itemize}[noitemsep,nolistsep]
           \item Effect: Retrieve and process the updates for the specified \ref{data:DatatypesMLModel} and store the updated \ref{data:DatatypesMLModel} in the \vpett{\nameref{mod:ComponentseHealthPlatformMLModelManagerMLModelRepository}}. 
           \item Sequence Diagrams: \cref{diag:Interaction:ProcessingMLModelupdates}
        \end{itemize}
    \end{itemize}
      \item[Diagrams:] None
    \end{description}

  %%%%%%%% MLModelUpdateMgmt
  \subsection{MLModelUpdateMgmt}\label{int:InterfacesMLModelUpdateMgmt}
    \begin{description}[noitemsep,nolistsep]
      \item[Provided by:] \iconcomponent{}~\vpett{\nameref{mod:ComponentseHealthPlatformMLModelManagerModelUpdater}}
      \item[Required by:] \iconcomponent{}~\vpett{\nameref{mod:ComponentseHealthPlatformMLModelManagerMLModelUpdateProcessor}} 
      \item[Operations:] ~
    \begin{itemize}[noitemsep,nolistsep,leftmargin=-.25cm]
      \item \textsf{\ref{data:DatatypesMLModel} processMLModelUpdates(\ref{data:DatatypesMLModel} model, List\textless{}\ref{data:DatatypesMLModelUpdate}\textgreater{} updates)}
        \begin{itemize}[noitemsep,nolistsep]
           \item Effect: Requests the \vpett{\nameref{mod:ComponentseHealthPlatformMLModelManagerModelUpdater}} to apply the list of  \ref{data:DatatypesMLModelUpdate}s on the provided \ref{data:DatatypesMLModel}. 
           \item Returns: The \ref{data:DatatypesMLModel} resulting from the application of the list of \ref{data:DatatypesMLModelUpdate}s on the initial \ref{data:DatatypesMLModel}.
           \item Sequence Diagrams: \cref{diag:Interaction:ProcessingMLModelupdates}
        \end{itemize}
    \end{itemize}
      \item[Diagrams:] None
    \end{description}

  %%%%%%%% OtherDataMgmt
  \subsection{OtherDataMgmt}\label{int:InterfacesOtherDataMgmt}
    \begin{description}[noitemsep,nolistsep]
      \item[Provided by:] \iconcomponent{}~\vpett{\nameref{comp:ComponentsP1PatientDataService}}, \iconcomponent{}~\vpett{\nameref{mod:ComponentsP1PatientDataServiceP1PatientStatusModule}}
      \item[Required by:] \iconcomponent{}~\vpett{\nameref{comp:ComponentseHealthPlatformClinicalJobCreator}}, \iconcomponent{}~\vpett{\nameref{comp:ComponentseHealthPlatformClinicalModelCache}}, \iconcomponent{}~\vpett{\nameref{comp:ComponentsP1PatientStatusCache}}, \iconcomponent{}~\vpett{\nameref{comp:ComponentseHealthPlatformRiskEstimationCombiner}}, \iconcomponent{}~\vpett{\nameref{comp:ComponentseHealthPlatformRiskEstimationScheduler}} 
      \item[Operations:] ~
    \begin{itemize}[noitemsep,nolistsep,leftmargin=-.25cm]
      \item \textsf{\ref{data:DatatypesPatientStatus} getPatientStatus(\ref{data:DatatypesPatientId} patientId) throws \ref{ex:ExceptionsNoSuchPatientException}}
        \begin{itemize}[noitemsep,nolistsep]
           \item Effect: Fetch the patient's most recent persisted status. 
           \item Returns: The patient's most recent persisted status.
           \item Sequence Diagrams: None
        \end{itemize}
      \item \textsf{void setEstimatedPatientStatus(\ref{data:DatatypesPatientId} patientId, \ref{data:DatatypesPatientStatus} estimatedStatus, \ref{data:DatatypesTimestamp} estimationTime) throws \ref{ex:ExceptionsNoSuchPatientException}}
        \begin{itemize}[noitemsep,nolistsep]
           \item Effect: Update the patient status and the time of estimation. If the patient's estimated risk changed, the appropriate parties are notified via P1IRiskEvents (this is the emission point). 
           \item Sequence Diagrams: \cref{diag:Interaction:Incomingsensordata,diag:Interaction:CombiningClinicalModelJobresults}
        \end{itemize}
    \end{itemize}
      \item[Diagrams:] None
    \end{description}

  %%%%%%%% P1ClinicalConfigMgmt
  \subsection{P1ClinicalConfigMgmt}\label{int:InterfacesP1ClinicalConfigMgmt}
    \begin{description}[noitemsep,nolistsep]
      \item[Provided by:] \iconcomponent{}~\vpett{\nameref{comp:ComponentseHealthPlatformClinicalModelDB}}
      \item[Required by:] \iconcomponent{}~\vpett{\nameref{mod:ComponentsP1PhysicianCommandServiceP1ConfigurationHandler}}, \iconcomponent{}~\vpett{\nameref{comp:ComponentsP1PhysicianCommandService}} 
      \item[Operations:] ~
    \begin{itemize}[noitemsep,nolistsep,leftmargin=-.25cm]
      \item \textsf{void setClinicalModelConfigForPatient(\ref{data:DatatypesPatientId} patientId, \ref{data:DatatypesClinicalModelConfiguration} config)}
        \begin{itemize}[noitemsep,nolistsep]
           \item Effect: Overwrite the per-patient \ref{data:DatatypesClinicalModelConfiguration} in \vpett{\nameref{comp:ComponentseHealthPlatformClinicalModelDB}}. The caller is expected to invalidate \vpett{\nameref{comp:ComponentseHealthPlatformClinicalModelCache}} after a successful write. 
           \item Sequence Diagrams: None
        \end{itemize}
    \end{itemize}
      \item[Diagrams:] None
    \end{description}

  %%%%%%%% P1IConfigCommand
  \subsection{P1IConfigCommand}\label{int:InterfacesP1IConfigCommand}
    \begin{description}[noitemsep,nolistsep]
      \item[Provided by:] \iconcomponent{}~\vpett{\nameref{mod:ComponentsP1PhysicianCommandServiceP1ConfigurationHandler}}
      \item[Required by:] \iconcomponent{}~\vpett{\nameref{mod:ComponentsP1PhysicianCommandServiceP1CommandRouter}} 
      \item[Operations:] ~
    \begin{itemize}[noitemsep,nolistsep,leftmargin=-.25cm]
      \item \textsf{void configurePatientRiskAssessment(\ref{data:DatatypesPatientId} patientId, \ref{data:DatatypesClinicalModelConfiguration} config)}
        \begin{itemize}[noitemsep,nolistsep]
           \item Effect: Internal delegate. \vpett{\nameref{mod:ComponentsP1PhysicianCommandServiceP1CommandRouter}} forwards the UC7 request to \vpett{\nameref{mod:ComponentsP1PhysicianCommandServiceP1ConfigurationHandler}}, which writes the config, invalidates the cache, and launches a normal-priority risk-estimation job so the new configuration takes effect immediately. 
           \item Sequence Diagrams: None
        \end{itemize}
    \end{itemize}
      \item[Diagrams:] None
    \end{description}

  %%%%%%%% P1ICorrelationTracking
  \subsection{P1ICorrelationTracking}\label{int:InterfacesP1ICorrelationTracking}
    \begin{description}[noitemsep,nolistsep]
      \item[Provided by:] \iconcomponent{}~\vpett{\nameref{mod:ComponentsP1PhysicianCommandServiceP1CorrelationTracker}}
      \item[Required by:] \iconcomponent{}~\vpett{\nameref{mod:ComponentsP1PhysicianCommandServiceP1OnDemandConsultationHandler}} 
      \item[Operations:] ~
    \begin{itemize}[noitemsep,nolistsep,leftmargin=-.25cm]
      \item \textsf{boolean consumeCorrelation(\ref{data:DatatypesCorrelationId} correlationId)}
        \begin{itemize}[noitemsep,nolistsep]
           \item Effect: Resolve a \ref{data:DatatypesCorrelationId} by removing it from the pending set if present. Report whether the event matches a tracked UC9 request. 
           \item Returns: True if the \ref{data:DatatypesCorrelationId} was pending and has now been consumed. False otherwise.
           \item Sequence Diagrams: None
        \end{itemize}
      \item \textsf{\ref{data:DatatypesCorrelationId} newCorrelationId()}
        \begin{itemize}[noitemsep,nolistsep]
           \item Effect: Mint a fresh \ref{data:DatatypesCorrelationId} and record it in the pending set so a future P1IRiskEvents push can be matched back to this UC9 request. 
           \item Returns: The freshly minted \ref{data:DatatypesCorrelationId}.
           \item Sequence Diagrams: None
        \end{itemize}
    \end{itemize}
      \item[Diagrams:] None
    \end{description}

  %%%%%%%% P1IHISAccess
  \subsection{P1IHISAccess}\label{int:InterfacesP1IHISAccess}
    \begin{description}[noitemsep,nolistsep]
      \item[Provided by:] \iconcomponent{}~\vpett{\nameref{comp:ComponentsP1HISAdapter}}
      \item[Required by:] \iconcomponent{}~\vpett{\nameref{mod:ComponentsP1PatientDataServiceP1EHRProxyModule}}, \iconcomponent{}~\vpett{\nameref{comp:ComponentsP1PatientDataService}} 
      \item[Operations:] ~
    \begin{itemize}[noitemsep,nolistsep,leftmargin=-.25cm]
      \item \textsf{\ref{data:DatatypesPatientRecord} getPatientRecord(\ref{data:DatatypesPatientId} patientId)}
        \begin{itemize}[noitemsep,nolistsep]
           \item Effect: Read the patient's record from the external \vpett{\nameref{comp:ComponentsHIS}} by combining the relevant healthAPI reads into one record-level view. 
           \item Returns: The aggregated \ref{data:DatatypesPatientRecord} (demographics, observations, and risk assessments per FHIR) retrieved from the \vpett{\nameref{comp:ComponentsHIS}}.
           \item Sequence Diagrams: None
        \end{itemize}
      \item \textsf{void updatePatientRecord(\ref{data:DatatypesPatientId} patientId, \ref{data:DatatypesPatientRecord} patientRecord)}
        \begin{itemize}[noitemsep,nolistsep]
           \item Effect: Write the patient record back to the \vpett{\nameref{comp:ComponentsHIS}} by dispatching the appropriate healthAPI save operations. 
           \item Sequence Diagrams: None
        \end{itemize}
    \end{itemize}
      \item[Diagrams:] None
    \end{description}

  %%%%%%%% P1INotificationInbox
  \subsection{P1INotificationInbox}\label{int:InterfacesP1INotificationInbox}
    \begin{description}[noitemsep,nolistsep]
      \item[Provided by:] \iconcomponent{}~\vpett{\nameref{comp:ComponentsP1NotificationDispatcher}}, \iconcomponent{}~\vpett{\nameref{mod:ComponentsP1NotificationDispatcherP1NotificationInboxModule}}
      \item[Required by:] \iconcomponent{}~\vpett{\nameref{comp:ComponentsP1PhysicianGateway}} 
      \item[Operations:] ~
    \begin{itemize}[noitemsep,nolistsep,leftmargin=-.25cm]
      \item \textsf{void acknowledgeNotification(\ref{data:DatatypesNotificationId} notificationId)}
        \begin{itemize}[noitemsep,nolistsep]
           \item Effect: Confirm delivery of the notification so the dispatcher can remove it from the priority buffer and the durable log. 
           \item Sequence Diagrams: None
        \end{itemize}
      \item \textsf{List\textless{}\ref{data:DatatypesNotification}\textgreater{} getPendingNotifications(\ref{data:DatatypesPhysicianId} physicianId)}
        \begin{itemize}[noitemsep,nolistsep]
           \item Effect: Return the physician's queued notifications in priority order (red, yellow, green) without removing them from the buffer. 
           \item Returns: The list of pending Notifications for the physician, ordered by severity then arrival.
           \item Sequence Diagrams: None
        \end{itemize}
      \item \textsf{void subscribeToNotifications(\ref{data:DatatypesPhysicianId} physicianId)}
        \begin{itemize}[noitemsep,nolistsep]
           \item Effect: Add the physician to the recipient registry so future patient-status events fan out to them. 
           \item Sequence Diagrams: None
        \end{itemize}
      \item \textsf{void unsubscribeFromNotifications(\ref{data:DatatypesPhysicianId} physicianId)}
        \begin{itemize}[noitemsep,nolistsep]
           \item Effect: Remove the physician from the recipient registry. Notifications already queued are still delivered. 
           \item Sequence Diagrams: None
        \end{itemize}
    \end{itemize}
      \item[Diagrams:] None
    \end{description}

  %%%%%%%% P1INotificationLog
  \subsection{P1INotificationLog}\label{int:InterfacesP1INotificationLog}
    \begin{description}[noitemsep,nolistsep]
      \item[Provided by:] \iconcomponent{}~\vpett{\nameref{mod:ComponentsP1NotificationDispatcherP1DurableNotificationLog}}
      \item[Required by:] \iconcomponent{}~\vpett{\nameref{mod:ComponentsP1NotificationDispatcherP1NotificationInboxModule}}, \iconcomponent{}~\vpett{\nameref{mod:ComponentsP1NotificationDispatcherP1RiskEventSubscriber}} 
      \item[Operations:] ~
    \begin{itemize}[noitemsep,nolistsep,leftmargin=-.25cm]
      \item \textsf{void markDelivered(\ref{data:DatatypesNotificationId} notificationId)}
        \begin{itemize}[noitemsep,nolistsep]
           \item Effect: Remove the notification from the durable log once the physician has acked delivery, bounding storage growth. 
           \item Sequence Diagrams: None
        \end{itemize}
      \item \textsf{void persist(\ref{data:DatatypesNotification} notification)}
        \begin{itemize}[noitemsep,nolistsep]
           \item Effect: Write the notification to the durable log before it enters the in-memory buffer so a crash cannot lose a red or yellow notification. 
           \item Sequence Diagrams: None
        \end{itemize}
      \item \textsf{List\textless{}\ref{data:DatatypesNotification}\textgreater{} recoverPending()}
        \begin{itemize}[noitemsep,nolistsep]
           \item Effect: On startup, return notifications that were persisted but never acked so \vpett{\nameref{mod:ComponentsP1NotificationDispatcherP1NotificationInboxModule}} can replay them into the priority buffer. 
           \item Returns: The list of unacked Notifications recovered from the durable log.
           \item Sequence Diagrams: None
        \end{itemize}
    \end{itemize}
      \item[Diagrams:] None
    \end{description}

  %%%%%%%% P1IOnDemandCommand
  \subsection{P1IOnDemandCommand}\label{int:InterfacesP1IOnDemandCommand}
    \begin{description}[noitemsep,nolistsep]
      \item[Provided by:] \iconcomponent{}~\vpett{\nameref{mod:ComponentsP1PhysicianCommandServiceP1OnDemandConsultationHandler}}
      \item[Required by:] \iconcomponent{}~\vpett{\nameref{mod:ComponentsP1PhysicianCommandServiceP1CommandRouter}} 
      \item[Operations:] ~
    \begin{itemize}[noitemsep,nolistsep,leftmargin=-.25cm]
      \item \textsf{\ref{data:DatatypesPatientStatus} pollConsultationResult(\ref{data:DatatypesConsultationId} consultationId)}
        \begin{itemize}[noitemsep,nolistsep]
           \item Effect: Internal delegate, \vpett{\nameref{mod:ComponentsP1PhysicianCommandServiceP1CommandRouter}} forwards the poll to \vpett{\nameref{mod:ComponentsP1PhysicianCommandServiceP1OnDemandConsultationHandler}}, which looks up the stashed result for this consultationId. The result becomes available once the matching P1IRiskEvents event (\ref{data:DatatypesCorrelationId}-filtered) has arrived; callers re-poll until non-null. The handler clears the entry after the result has been returned once. 
           \item Returns: The \ref{data:DatatypesPatientStatus} produced by the UC9 risk estimation, or null if the matching event has not yet arrived.
           \item Sequence Diagrams: None
        \end{itemize}
      \item \textsf{\ref{data:DatatypesCorrelationId} requestOnDemandConsultation(\ref{data:DatatypesPatientId} patientId)}
        \begin{itemize}[noitemsep,nolistsep]
           \item Effect: Internal delegate. \vpett{\nameref{mod:ComponentsP1PhysicianCommandServiceP1CommandRouter}} forwards the UC9 request to \vpett{\nameref{mod:ComponentsP1PhysicianCommandServiceP1OnDemandConsultationHandler}}, which mints a \ref{data:DatatypesCorrelationId}, fetches sensor data, launches a priority risk job, and (when the matching \ref{data:DatatypesRiskEvent} arrives) pushes the result back to physicianId via P1IConsultationResultPush. 
           \item Returns: The \ref{data:DatatypesCorrelationId} minted for this UC9 request.
           \item Sequence Diagrams: None
        \end{itemize}
    \end{itemize}
      \item[Diagrams:] None
    \end{description}

  %%%%%%%% P1IOnDemandSensorFetch
  \subsection{P1IOnDemandSensorFetch}\label{int:InterfacesP1IOnDemandSensorFetch}
    \begin{description}[noitemsep,nolistsep]
      \item[Provided by:] \iconcomponent{}~\vpett{\nameref{comp:ComponentsP1PatientGatewayCommander}}
      \item[Required by:] \iconcomponent{}~\vpett{\nameref{mod:ComponentsP1PhysicianCommandServiceP1OnDemandConsultationHandler}}, \iconcomponent{}~\vpett{\nameref{comp:ComponentsP1PhysicianCommandService}} 
      \item[Operations:] ~
    \begin{itemize}[noitemsep,nolistsep,leftmargin=-.25cm]
      \item \textsf{\ref{data:DatatypesSensorDataPackage} requestCurrentSensorData(\ref{data:DatatypesPatientId} patientId, \ref{data:DatatypesCorrelationId} correlationId)}
        \begin{itemize}[noitemsep,nolistsep]
           \item Effect: Synchronously ask the patient gateway to push the current sensor reading. Bounded by the commander timeout sized inside the 3-minute UC9 initiation budget. 
           \item Returns: The fresh \ref{data:DatatypesSensorDataPackage} supplied by the patient gateway for this UC9 request.
           \item Sequence Diagrams: None
        \end{itemize}
    \end{itemize}
      \item[Diagrams:] None
    \end{description}

  %%%%%%%% P1IPatientQuery
  \subsection{P1IPatientQuery}\label{int:InterfacesP1IPatientQuery}
    \begin{description}[noitemsep,nolistsep]
      \item[Provided by:] \iconcomponent{}~\vpett{\nameref{comp:ComponentsP1PatientQueryService}}
      \item[Required by:] \iconcomponent{}~\vpett{\nameref{comp:ComponentsP1PhysicianGateway}} 
      \item[Operations:] ~
    \begin{itemize}[noitemsep,nolistsep,leftmargin=-.25cm]
      \item \textsf{\ref{data:DatatypesPatientStatus} consultPatientStatus(\ref{data:DatatypesPatientId} patientId)}
        \begin{itemize}[noitemsep,nolistsep]
           \item Effect: Read the patient's current status. The request is dispatched to the priority tier (high, medium, low) matching the patient's last-known risk level from the in-memory tier index inside \vpett{\nameref{comp:ComponentsP1PatientQueryService}}, so the tier decision is O(1) on the hot path. Unknown patients fall back to the high tier so the SLA fails safe. 
           \item Returns: The patient's most recent status.
           \item Sequence Diagrams: None
        \end{itemize}
    \end{itemize}
      \item[Diagrams:] None
    \end{description}

  %%%%%%%% P1IPatientStatusRead
  \subsection{P1IPatientStatusRead}\label{int:InterfacesP1IPatientStatusRead}
    \begin{description}[noitemsep,nolistsep]
      \item[Provided by:] \iconcomponent{}~\vpett{\nameref{comp:ComponentsP1PatientStatusCache}}
      \item[Required by:] \iconcomponent{}~\vpett{\nameref{comp:ComponentsP1PatientQueryService}} 
      \item[Operations:] ~
    \begin{itemize}[noitemsep,nolistsep,leftmargin=-.25cm]
      \item \textsf{\ref{data:DatatypesPatientStatus} getPatientStatus(\ref{data:DatatypesPatientId} patientId)}
        \begin{itemize}[noitemsep,nolistsep]
           \item Effect: Cache-first read of the patient's status. On miss, the cache fetches via OtherDataMgmt and populates itself for subsequent reads. 
           \item Returns: The patient's most recent status.
           \item Sequence Diagrams: None
        \end{itemize}
    \end{itemize}
      \item[Diagrams:] None
    \end{description}

  %%%%%%%% P1IPhysicianAPI
  \subsection{P1IPhysicianAPI}\label{int:InterfacesP1IPhysicianAPI}
    \begin{description}[noitemsep,nolistsep]
      \item[Provided by:] \iconcomponent{}~\vpett{\nameref{comp:ComponentsP1PhysicianGateway}}
      \item[Required by:] None 
      \item[Operations:] ~
    \begin{itemize}[noitemsep,nolistsep,leftmargin=-.25cm]
      \item \textsf{void acknowledgeNotification(\ref{data:DatatypesNotificationId} notificationId)}
        \begin{itemize}[noitemsep,nolistsep]
           \item Effect: Confirm delivery of the notification so the dispatcher can remove it from the priority buffer and the durable log. 
           \item Sequence Diagrams: None
        \end{itemize}
      \item \textsf{void configurePatientRiskAssessment(\ref{data:DatatypesPatientId} patientId, \ref{data:DatatypesClinicalModelConfiguration} clinicalModelConfiguration)}
        \begin{itemize}[noitemsep,nolistsep]
           \item Effect: Overwrite the per-patient \ref{data:DatatypesClinicalModelConfiguration}, invalidate the clinical-model cache, and launch a normal-priority risk-estimation job so the new configuration takes effect immediately. 
           \item Sequence Diagrams: None
        \end{itemize}
      \item \textsf{\ref{data:DatatypesPatientStatus} consultPatientStatus(\ref{data:DatatypesPatientId} patientId)}
        \begin{itemize}[noitemsep,nolistsep]
           \item Effect: Return the named patient's current status. SLA tier (2 / 5 / 10 s) is picked from the patient's last-known risk level via the tier index inside \vpett{\nameref{comp:ComponentsP1PatientQueryService}}. 
           \item Returns: The patient's most recent status: risk level, supporting data, and timestamp.
           \item Sequence Diagrams: None
        \end{itemize}
      \item \textsf{List\textless{}{\colorbox{red!30}{\underline{Dataypes.Notification}}}\textgreater{} getPendingNotifications()}
        \begin{itemize}[noitemsep,nolistsep]
           \item Effect: Return the physician's queued notifications in priority order (red, yellow, green) without removing them from the buffer. 
           \item Returns: The list of pending Notifications for the physician, ordered by severity then arrival.
           \item Sequence Diagrams: None
        \end{itemize}
      \item \textsf{\ref{data:DatatypesCorrelationId} requestOnDemandConsultation(\ref{data:DatatypesPatientId} patientId, \ref{data:DatatypesPhysicianId} physicianId)}
        \begin{itemize}[noitemsep,nolistsep]
           \item Effect: Initiate an on-demand consultation. Fetches fresh sensor data, launches a priority risk estimation, and (when the result is ready) pushes it back to physicianId via P1IConsultationResultPush. No polling. 
           \item Returns: The \ref{data:DatatypesCorrelationId} minted for this request, used by the client to match the eventual push.
           \item Sequence Diagrams: None
        \end{itemize}
      \item \textsf{void subscribeToNotifications(\ref{data:DatatypesPhysicianId} physicianId)}
        \begin{itemize}[noitemsep,nolistsep]
           \item Effect: Register a physician to receive notifications about subsequent patient-status changes. 
           \item Sequence Diagrams: None
        \end{itemize}
      \item \textsf{void unsubscribeFromNotifications(\ref{data:DatatypesPhysicianId} physicianId)}
        \begin{itemize}[noitemsep,nolistsep]
           \item Effect: Remove the physician from the recipient registry. Notifications already queued for them are still delivered. 
           \item Sequence Diagrams: None
        \end{itemize}
      \item \textsf{void updatePatientRiskLevel(\ref{data:DatatypesPatientId} patientId, \ref{data:DatatypesPatientStatus} patientStatus)}
        \begin{itemize}[noitemsep,nolistsep]
           \item Effect: Update the patient's internal risk level (which fires P1IRiskEvents) and forward the change to the EHR via the \vpett{\nameref{comp:ComponentsHIS}} proxy. 
           \item Sequence Diagrams: None
        \end{itemize}
    \end{itemize}
      \item[Diagrams:] None
    \end{description}

  %%%%%%%% P1IPhysicianCommand
  \subsection{P1IPhysicianCommand}\label{int:InterfacesP1IPhysicianCommand}
    \begin{description}[noitemsep,nolistsep]
      \item[Provided by:] \iconcomponent{}~\vpett{\nameref{mod:ComponentsP1PhysicianCommandServiceP1CommandRouter}}, \iconcomponent{}~\vpett{\nameref{comp:ComponentsP1PhysicianCommandService}}
      \item[Required by:] \iconcomponent{}~\vpett{\nameref{comp:ComponentsP1PhysicianGateway}} 
      \item[Operations:] ~
    \begin{itemize}[noitemsep,nolistsep,leftmargin=-.25cm]
      \item \textsf{void configurePatientRiskAssessment(\ref{data:DatatypesPatientId} patientId, \ref{data:DatatypesClinicalModelConfiguration} clinicalModelConfiguration)}
        \begin{itemize}[noitemsep,nolistsep]
           \item Effect: Overwrite the per-patient \ref{data:DatatypesClinicalModelConfiguration}, invalidate the clinical-model cache, and launch a normal-priority risk-estimation job so the new configuration takes effect immediately. 
           \item Sequence Diagrams: None
        \end{itemize}
      \item \textsf{\ref{data:DatatypesCorrelationId} requestOnDemandConsultation(\ref{data:DatatypesPatientId} patientId, \ref{data:DatatypesPhysicianId} physicianId)}
        \begin{itemize}[noitemsep,nolistsep]
           \item Effect: Initiate an on-demand consultation. Mints a \ref{data:DatatypesCorrelationId}, fetches fresh sensor data, launches a priority risk job, and pushes the result back to physicianId via P1IConsultationResultPush when ready. 
           \item Returns: The \ref{data:DatatypesCorrelationId} minted for this UC9 request.
           \item Sequence Diagrams: None
        \end{itemize}
      \item \textsf{void updatePatientRiskLevel(\ref{data:DatatypesPatientId} patientId, \ref{data:DatatypesPatientStatus} patientStatus)}
        \begin{itemize}[noitemsep,nolistsep]
           \item Effect: Update the patient's internal risk level (which fires P1IRiskEvents) and forward the change to the EHR via the \vpett{\nameref{comp:ComponentsHIS}} proxy. 
           \item Sequence Diagrams: None
        \end{itemize}
    \end{itemize}
      \item[Diagrams:] None
    \end{description}

  %%%%%%%% P1IPriorityBuffer
  \subsection{P1IPriorityBuffer}\label{int:InterfacesP1IPriorityBuffer}
    \begin{description}[noitemsep,nolistsep]
      \item[Provided by:] \iconcomponent{}~\vpett{\nameref{mod:ComponentsP1NotificationDispatcherP1PriorityBuffer}}
      \item[Required by:] \iconcomponent{}~\vpett{\nameref{mod:ComponentsP1NotificationDispatcherP1NotificationInboxModule}}, \iconcomponent{}~\vpett{\nameref{mod:ComponentsP1NotificationDispatcherP1RiskEventSubscriber}} 
      \item[Operations:] ~
    \begin{itemize}[noitemsep,nolistsep,leftmargin=-.25cm]
      \item \textsf{void enqueue(\ref{data:DatatypesPhysicianId} physicianId, \ref{data:DatatypesNotification} notification)}
        \begin{itemize}[noitemsep,nolistsep]
           \item Effect: Add the notification to the physician's pending set. Severity drives priority ordering and the drop-green-first policy when the 100 per minute threshold is exceeded. 
           \item Sequence Diagrams: None
        \end{itemize}
      \item \textsf{List\textless{}\ref{data:DatatypesNotification}\textgreater{} peekPending(\ref{data:DatatypesPhysicianId} physicianId)}
        \begin{itemize}[noitemsep,nolistsep]
           \item Effect: Return the physician's pending notifications in priority order (red, yellow, green) without removing them from the buffer. 
           \item Returns: The list of pending Notifications for the physician, ordered by severity then arrival.
           \item Sequence Diagrams: None
        \end{itemize}
      \item \textsf{void removePending(\ref{data:DatatypesNotificationId} notificationId)}
        \begin{itemize}[noitemsep,nolistsep]
           \item Effect: Remove the notification from the buffer after the physician acks it through P1INotificationInbox. 
           \item Sequence Diagrams: None
        \end{itemize}
    \end{itemize}
      \item[Diagrams:] None
    \end{description}

  %%%%%%%% P1IRecipientRegistry
  \subsection{P1IRecipientRegistry}\label{int:InterfacesP1IRecipientRegistry}
    \begin{description}[noitemsep,nolistsep]
      \item[Provided by:] \iconcomponent{}~\vpett{\nameref{mod:ComponentsP1NotificationDispatcherP1RecipientRegistry}}
      \item[Required by:] \iconcomponent{}~\vpett{\nameref{mod:ComponentsP1NotificationDispatcherP1NotificationInboxModule}}, \iconcomponent{}~\vpett{\nameref{mod:ComponentsP1NotificationDispatcherP1RiskEventSubscriber}} 
      \item[Operations:] ~
    \begin{itemize}[noitemsep,nolistsep,leftmargin=-.25cm]
      \item \textsf{List\textless{}\ref{data:DatatypesPhysicianId}\textgreater{} findRecipients(\ref{data:DatatypesPatientId} patientId)}
        \begin{itemize}[noitemsep,nolistsep]
           \item Effect: Look up the physicians who should be notified about events for the patient. Backs UC5 recipient resolution. 
           \item Returns: The list of PhysicianIds subscribed to events for this patient.
           \item Sequence Diagrams: None
        \end{itemize}
      \item \textsf{void subscribe(\ref{data:DatatypesPhysicianId} physicianId)}
        \begin{itemize}[noitemsep,nolistsep]
           \item Effect: Add the physician to the notification recipient set. 
           \item Sequence Diagrams: None
        \end{itemize}
      \item \textsf{void unsubscribe(\ref{data:DatatypesPhysicianId} physicianId)}
        \begin{itemize}[noitemsep,nolistsep]
           \item Effect: Remove the physician from the notification recipient set. 
           \item Sequence Diagrams: None
        \end{itemize}
    \end{itemize}
      \item[Diagrams:] None
    \end{description}

  %%%%%%%% P1IRiskEventListener
  \subsection{P1IRiskEventListener}\label{int:InterfacesP1IRiskEventListener}
    \begin{description}[noitemsep,nolistsep]
      \item[Provided by:] \iconcomponent{}~\vpett{\nameref{comp:ComponentsP1NotificationDispatcher}}
      \item[Required by:] \iconcomponent{}~\vpett{\nameref{comp:ComponentsP1PatientDataService}} 
      \item[Operations:] ~
    \begin{itemize}[noitemsep,nolistsep,leftmargin=-.25cm]
      \item \textsf{void onRiskEvent(\ref{data:DatatypesRiskEvent} event)}
        \begin{itemize}[noitemsep,nolistsep]
           \item Effect: Push callback invoked by \vpett{\nameref{mod:ComponentsP1PatientDataServiceP1PatientStatusModule}} for every subscription whose \ref{data:DatatypesFilterCriteria} matches the event. Receiver applies its own handling: \vpett{\nameref{mod:ComponentsP1NotificationDispatcherP1RiskEventSubscriber}} fans out to physicians, \vpett{\nameref{mod:ComponentsP1PhysicianCommandServiceP1OnDemandConsultationHandler}} delivers UC9 results, \vpett{\nameref{comp:ComponentsP1PatientQueryService}} refreshes its tier index. 
           \item Sequence Diagrams: None
        \end{itemize}
    \end{itemize}
      \item[Diagrams:] None
    \end{description}

  %%%%%%%% P1IRiskEvents
  \subsection{P1IRiskEvents}\label{int:InterfacesP1IRiskEvents}
    \begin{description}[noitemsep,nolistsep]
      \item[Provided by:] \iconcomponent{}~\vpett{\nameref{comp:ComponentsP1PatientDataService}}, \iconcomponent{}~\vpett{\nameref{mod:ComponentsP1PatientDataServiceP1PatientStatusModule}}
      \item[Required by:] \iconcomponent{}~\vpett{\nameref{comp:ComponentsP1NotificationDispatcher}}, \iconcomponent{}~\vpett{\nameref{mod:ComponentsP1PhysicianCommandServiceP1OnDemandConsultationHandler}}, \iconcomponent{}~\vpett{\nameref{comp:ComponentsP1PatientQueryService}}, \iconcomponent{}~\vpett{\nameref{comp:ComponentsP1PhysicianCommandService}}, \iconcomponent{}~\vpett{\nameref{mod:ComponentsP1NotificationDispatcherP1RiskEventSubscriber}} 
      \item[Operations:] ~
    \begin{itemize}[noitemsep,nolistsep,leftmargin=-.25cm]
      \item \textsf{\ref{data:DatatypesSubscriptionId} subscribe(\ref{data:DatatypesSubscriberId} subscriberId, \ref{data:DatatypesFilterCriteria} filterCriteria)}
        \begin{itemize}[noitemsep,nolistsep]
           \item Effect: Register a subscriber to receive \ref{data:DatatypesRiskEvent} payloads matching the filter. An empty filter receives every event. Delivery happens by callback on the subscriber's P1IRiskEventListener. 
           \item Returns: A \ref{data:DatatypesSubscriptionId} the caller uses to unsubscribe later.
           \item Sequence Diagrams: None
        \end{itemize}
      \item \textsf{void unsubscribe(\ref{data:DatatypesSubscriptionId} subscriptionId)}
        \begin{itemize}[noitemsep,nolistsep]
           \item Effect: Remove the named subscription so no further RiskEvents are pushed to it. 
           \item Sequence Diagrams: None
        \end{itemize}
    \end{itemize}
      \item[Diagrams:] None
    \end{description}

  %%%%%%%% P1IRiskLevelCommand
  \subsection{P1IRiskLevelCommand}\label{int:InterfacesP1IRiskLevelCommand}
    \begin{description}[noitemsep,nolistsep]
      \item[Provided by:] \iconcomponent{}~\vpett{\nameref{mod:ComponentsP1PhysicianCommandServiceP1RiskLevelHandler}}
      \item[Required by:] \iconcomponent{}~\vpett{\nameref{mod:ComponentsP1PhysicianCommandServiceP1CommandRouter}} 
      \item[Operations:] ~
    \begin{itemize}[noitemsep,nolistsep,leftmargin=-.25cm]
      \item \textsf{void updatePatientRiskLevel(\ref{data:DatatypesPatientId} patientId, \ref{data:DatatypesPatientStatus} patientStatus)}
        \begin{itemize}[noitemsep,nolistsep]
           \item Effect: Internal delegate. \vpett{\nameref{mod:ComponentsP1PhysicianCommandServiceP1CommandRouter}} forwards the UC8 request to \vpett{\nameref{mod:ComponentsP1PhysicianCommandServiceP1RiskLevelHandler}}, which writes the new internal risk level via `OtherDataMgmt.setEstimatedPatientStatus` (firing P1IRiskEvents) and then forwards to the EHR via PatientRecordMgmt. 
           \item Sequence Diagrams: None
        \end{itemize}
    \end{itemize}
      \item[Diagrams:] None
    \end{description}

  %%%%%%%% PatientRecordMgmt
  \subsection{PatientRecordMgmt}\label{int:InterfacesPatientRecordMgmt}
    \begin{description}[noitemsep,nolistsep]
      \item[Provided by:] \iconcomponent{}~\vpett{\nameref{mod:ComponentsP1PatientDataServiceP1EHRProxyModule}}, \iconcomponent{}~\vpett{\nameref{comp:ComponentsP1PatientDataService}}
      \item[Required by:] \iconcomponent{}~\vpett{\nameref{comp:ComponentsP1PhysicianCommandService}}, \iconcomponent{}~\vpett{\nameref{mod:ComponentsP1PhysicianCommandServiceP1RiskLevelHandler}}, \iconcomponent{}~\vpett{\nameref{comp:ComponentseHealthPlatformRiskEstimationCombiner}} 
      \item[Operations:] ~
    \begin{itemize}[noitemsep,nolistsep,leftmargin=-.25cm]
      \item \textsf{\ref{data:DatatypesPatientRecord} getPatientRecord(\ref{data:DatatypesPatientId} patientId) throws \ref{ex:ExceptionsNoSuchPatientRecordException}}
        \begin{itemize}[noitemsep,nolistsep]
           \item Effect: Fetch the EHR record of the patient with given id from the \vpett{\nameref{comp:ComponentsHIS}} and return it. If the \vpett{\nameref{comp:ComponentsHIS}} is unavailable, return an older cached copy of the patient record when possible. 
           \item Returns: The requested patient record.
           \item Sequence Diagrams: None
        \end{itemize}
      \item \textsf{void updatePatientRecord(\ref{data:DatatypesPatientId} patientId, \ref{data:DatatypesPatientRecord} patientRecord)}
        \begin{itemize}[noitemsep,nolistsep]
           \item Effect: Write the patient record back to \vpett{\nameref{comp:ComponentsHIS}} via \vpett{\nameref{comp:ComponentsP1HISAdapter}} and invalidate the EHR cache entry so subsequent reads see the new value. 
           \item Sequence Diagrams: None
        \end{itemize}
    \end{itemize}
      \item[Diagrams:] None
    \end{description}

  %%%%%%%% Results
  \subsection{Results}\label{int:InterfacesResults}
    \begin{description}[noitemsep,nolistsep]
      \item[Provided by:] \iconcomponent{}~\vpett{\nameref{comp:ComponentseHealthPlatformRiskEstimationCombiner}}
      \item[Required by:] \iconcomponent{}~\vpett{\nameref{mod:ComponentseHealthPlatformRiskEstimationProcessorJobProcessor}}, \iconcomponent{}~\vpett{\nameref{comp:ComponentseHealthPlatformRiskEstimationProcessor}} 
      \item[Operations:] ~
    \begin{itemize}[noitemsep,nolistsep,leftmargin=-.25cm]
      \item \textsf{void addClinicalModelJobResults(\ref{data:DatatypesClinicalModelJobId} jobId, \ref{data:DatatypesClinicalModelJobResult} results)}
        \begin{itemize}[noitemsep,nolistsep]
           \item Effect: Sends the result of the performed clinical model computation (identified by the jobId) to the \vpett{\nameref{comp:ComponentseHealthPlatformRiskEstimationCombiner}} for combining with the other partial results belonging to the same risk estimation. 
           \item Sequence Diagrams: \cref{diag:Interaction:Computeclinicalmodelresult,diag:Interaction:CombiningClinicalModelJobresults}
        \end{itemize}
    \end{itemize}
      \item[Diagrams:] None
    \end{description}

  %%%%%%%% SensorDataAnalysis
  \subsection{SensorDataAnalysis}\label{int:InterfacesSensorDataAnalysis}
    \begin{description}[noitemsep,nolistsep]
      \item[Provided by:] \iconcomponent{}~\vpett{\nameref{mod:ComponentseHealthPlatformRiskEstimationProcessorSensorDataRiskAssessment}}
      \item[Required by:] \iconcomponent{}~\vpett{\nameref{mod:ComponentseHealthPlatformRiskEstimationProcessorJobProcessor}} 
      \item[Operations:] ~
    \begin{itemize}[noitemsep,nolistsep,leftmargin=-.25cm]
      \item \textsf{void performModelPrediction(\ref{data:DatatypesClinicalModelJob} clinicalModelJob, \ref{data:DatatypesMLModel} model)}
        \begin{itemize}[noitemsep,nolistsep]
           \item Effect: Perform a prediction using the provided \ref{data:DatatypesMLModel} and the data provided in the clinicalModelJob definition.
Afterwards the results are submitted for combination by the \vpett{\nameref{comp:ComponentseHealthPlatformRiskEstimationCombiner}}. 
           \item Sequence Diagrams: \cref{diag:Interaction:Computeclinicalmodelresult}
        \end{itemize}
    \end{itemize}
      \item[Diagrams:] None
    \end{description}

  %%%%%%%% SensorDataMgmt
  \subsection{SensorDataMgmt}\label{int:InterfacesSensorDataMgmt}
    \begin{description}[noitemsep,nolistsep]
      \item[Provided by:] \iconcomponent{}~\vpett{\nameref{comp:ComponentsAv2SensorDataRepository}}, \iconcomponent{}~\vpett{\nameref{comp:ComponentsP1PatientDataService}}, \iconcomponent{}~\vpett{\nameref{mod:ComponentsP1PatientDataServiceP1SensorIngestModule}}
      \item[Required by:] \iconcomponent{}~\vpett{\nameref{comp:ComponentsAv2DataIngestionService}}, \iconcomponent{}~\vpett{\nameref{comp:ComponentseHealthPlatformClinicalJobCreator}}, \iconcomponent{}~\vpett{\nameref{mod:ComponentsP1PhysicianCommandServiceP1OnDemandConsultationHandler}}, \iconcomponent{}~\vpett{\nameref{comp:ComponentsP1PhysicianCommandService}}, \iconcomponent{}~\vpett{\nameref{comp:ComponentseHealthPlatformRiskEstimationScheduler}} 
      \item[Operations:] ~
    \begin{itemize}[noitemsep,nolistsep,leftmargin=-.25cm]
      \item \textsf{void addSensorData(\ref{data:DatatypesPatientId} patientId, \ref{data:DatatypesSensorDataPackage} sensorData, \ref{data:DatatypesTimestamp} receivedAt, boolean triggerEstimation)}
        \begin{itemize}[noitemsep,nolistsep]
           \item Effect: Store the given sensor data and meta-data. When `triggerEstimation` is true (UC4 default) the ingest module launches a risk-estimation job on the new data. When false the call is store-only. UC9 callers pass `false` because they explicitly launch a priority risk job and would otherwise double the scheduler load for every on-demand consultation. 
           \item Sequence Diagrams: None
        \end{itemize}
      \item \textsf{Map\textless{}\ref{data:DatatypesTimestamp}, \ref{data:DatatypesSensorDataPackage}\textgreater{} getAllSensorDataOfPatient(\ref{data:DatatypesPatientId} patientId)}
        \begin{itemize}[noitemsep,nolistsep]
           \item Effect: Fetch and return all sensor data belonging to the patient identified by patientId. 
           \item Returns: The sensor data and the timestamp of arrival in the system.
           \item Sequence Diagrams: None
        \end{itemize}
      \item \textsf{Map\textless{}\ref{data:DatatypesTimestamp}, \ref{data:DatatypesSensorDataPackage}\textgreater{} getAllSensorDataOfPatientBefore(\ref{data:DatatypesPatientId} patientId, \ref{data:DatatypesTimestamp} stopTime)}
        \begin{itemize}[noitemsep,nolistsep]
           \item Effect: Fetch and return all sensor data belonging to the patient identified by patientId which was received before the specified stopTime. 
           \item Returns: The sensor data and the timestamp of arrival in the system.
           \item Sequence Diagrams: \cref{diag:Interaction:Riskestimationprocess}
        \end{itemize}
    \end{itemize}
      \item[Diagrams:] None
    \end{description}

  %%%%%%%% UserAuthentication
  \subsection{UserAuthentication}\label{int:InterfacesUserAuthentication}
    \begin{description}[noitemsep,nolistsep]
      \item[Provided by:] \iconcomponent{}~\vpett{\nameref{comp:ComponentsPatientRecordRepository}}
      \item[Required by:] \iconcomponent{}~\vpett{\nameref{comp:ComponentsUserAuthenticationService}} 
      \item[Operations:] ~
    \begin{itemize}[noitemsep,nolistsep,leftmargin=-.25cm]
      \item \textsf{boolean authenticate(string userId, string credentials)}
        \begin{itemize}[noitemsep,nolistsep]
           \item Effect: Verifies the provided credentials against the stored record in \vpett{\nameref{comp:ComponentsPatientRecordRepository}}. Returns true if authentication succeeds. Throws AuthenticationException if credentials are missing or malformed. 
           \item Returns: True if the provided credentials match the stored credentials for the specified user in the \vpett{\nameref{comp:ComponentsPatientRecordRepository}}. False is never returned — if credentials do not match, AuthenticationException is thrown instead, so a true return always means authentication succeeded.
           \item Sequence Diagrams: None
        \end{itemize}
      \item \textsf{void revokeSession(string userId, string sessionToken)}
        \begin{itemize}[noitemsep,nolistsep]
           \item Effect: Invalidates the specified session token for the given user immediately. Any subsequent request carrying the revoked token is rejected with an authentication error. Called on explicit logout or when an administrative session revocation is triggered. The token is added to a revocation list consulted by \vpett{\nameref{comp:ComponentsUserAuthenticationService}} on every incoming request. 
           \item Sequence Diagrams: None
        \end{itemize}
    \end{itemize}
      \item[Diagrams:] None
    \end{description}

% END INTERFACES

% NODES
\section{Nodes}\label{sec:nodes}
\subsection{AdminNode}\label{node:DeploymentAdminNode}
	\begin{description}[noitemsep,nolistsep]
		\item[Responsibility:]~Server node within the PMS backend that hosts administrative coordination components, in particular the RedeploymentCoordinator. Provides the control plane the AdminPortal connects to in order to trigger redeployment or rollback of failed subsystems within the 5-minute Av2 deadline.
		\item[Visible on diagrams:]~\cref{diag:Deployment:Primarydeploymentdiagram}		
	\end{description}

\subsection{AdminNode (Pilot Deployment)}\label{node:DeploymentAdminNodePilotDeployment}
	\begin{description}[noitemsep,nolistsep]
		\item[Responsibility:]~Pilot-deployment instance of the administrative server node hosting the RedeploymentCoordinator. Single instance shared across all pilot patients and clinical models.
		\item[Visible on diagrams:]~\cref{diag:Deployment:PilotDeployment200patients5clinmodels}		
	\end{description}

\subsection{AdminWorkstation}\label{node:DeploymentAdminWorkstation}
	\begin{description}[noitemsep,nolistsep]
		\item[Responsibility:]~Client device (workstation, laptop) used by the system administrator. Hosts the AdminPortal web client through which administrators receive notifications and trigger redeployments or rollbacks. Outside the PMS server-side deployment boundary.
		\item[Visible on diagrams:]~\cref{diag:Deployment:Contextdiagramforthedeploymentview}		
	\end{description}

\subsection{BackendGatewayNode}\label{node:DeploymentBackendGatewayNode}
	\begin{description}[noitemsep,nolistsep]
		\item[Responsibility:]~Server node hosting all backend-side gateway components: CommunicationGateway, MessageBroker, DataIngestionService, AcknowledgementService, and EmergencyBackupReceiver. Acts as the single backend ingress point for all primary-channel and SMS-backup patient-gateway traffic.
		\item[Visible on diagrams:]~\cref{diag:Deployment:Primarydeploymentdiagram}		
	\end{description}

\subsection{BackendGatewayNode (Pilot Deployment)}\label{node:DeploymentBackendGatewayNodePilotDeployment}
	\begin{description}[noitemsep,nolistsep]
		\item[Responsibility:]~Pilot-deployment instance of the backend gateway tier hosting CommunicationGateway, MessageBroker, DataIngestionService, AcknowledgementService, and EmergencyBackupReceiver. Single ingress point for all 200 pilot patient gateways over both the primary and SMS-backup channels.
		\item[Visible on diagrams:]~\cref{diag:Deployment:PilotDeployment200patients5clinmodels}		
	\end{description}

\subsection{DevBackendNode (Dev Deployment)}\label{node:DeploymentDevBackendNodeDevDeployment}
	\begin{description}[noitemsep,nolistsep]
		\item[Responsibility:]~Single-machine development node that collapses all backend-side responsibilities (gateway ingress, broker, ingestion, acknowledgement, SMS receiver, monitoring, recovery sync, redeployment coordination) onto one host to simplify local testing of the 20-patient dev deployment.
		\item[Visible on diagrams:]~\cref{diag:Deployment:DevelopmentTestDeployment20patients3clinmodels}		
	\end{description}

\subsection{DevGatewayNode (Dev Deployment)}\label{node:DeploymentDevGatewayNodeDevDeployment}
	\begin{description}[noitemsep,nolistsep]
		\item[Responsibility:]~Development-deployment instance of the patient-side device hosting the PatientGateway. Replicated 20 times in the dev deployment to simulate the dev-scale patient population.
		\item[Visible on diagrams:]~\cref{diag:Deployment:DevelopmentTestDeployment20patients3clinmodels}		
	\end{description}

\subsection{DevPatientRecordsNode (Dev Deployment)}\label{node:DeploymentDevPatientRecordsNodeDevDeployment}
	\begin{description}[noitemsep,nolistsep]
		\item[Responsibility:]~Development-deployment node hosting the patient-data persistence components (\vpett{\nameref{comp:ComponentsPatientRecordRepository}}, SensorDataRepository, and related management/status/authentication services) used by the local dev environment.
		\item[Visible on diagrams:]~\cref{diag:Deployment:DevelopmentTestDeployment20patients3clinmodels}		
	\end{description}

\subsection{eHealthPlatform}\label{node:DeploymenteHealthPlatform}
	\begin{description}[noitemsep,nolistsep]
		\item[Responsibility:]~Container node representing the deployment context boundary for the \vpett{\nameref{comp:ComponentseHealthPlatform}}.
Hosts NotificationDispatcher, PhysicianCommandService, PatientDataService, PatientQueryService, PatientStatusCache.
		\item[Visible on diagrams:]~\cref{diag:Deployment:Contextdiagramforthedeploymentview}		
	\end{description}

\subsection{GatewayNode}\label{node:DeploymentGatewayNode}
	\begin{description}[noitemsep,nolistsep]
		\item[Responsibility:]~\ref{data:Datatypesorghl7fhirr5modelPatient}-side device node hosting the PatientGateway component. One node instance is deployed per enrolled patient; the multiplicity is captured on the deployment-association line to BackendGatewayNode, not on the node itself.
		\item[Visible on diagrams:]~\cref{diag:Deployment:Primarydeploymentdiagram}		
	\end{description}

\subsection{GatewayNode (Pilot Deployment)}\label{node:DeploymentGatewayNodePilotDeployment}
	\begin{description}[noitemsep,nolistsep]
		\item[Responsibility:]~Pilot-deployment instance of the patient-side device hosting the PatientGateway. Replicated 200 times in the pilot deployment, once per enrolled patient.
		\item[Visible on diagrams:]~\cref{diag:Deployment:PilotDeployment200patients5clinmodels}		
	\end{description}

\subsection{HISServer}\label{node:DeploymentHISServer}
	\begin{description}[noitemsep,nolistsep]
		\item[Responsibility:]~External infrastructure hosting the Hospital Information System. Exposes the healthAPI (FHIR-based) consumed by PMS backend components that need to read or write patient records, observations, and risk assessments. Outside the PMS deployment boundary.
		\item[Visible on diagrams:]~\cref{diag:Deployment:Contextdiagramforthedeploymentview}		
	\end{description}

\subsection{MLaaS Conn (Pilot Deployment)}\label{node:DeploymentMLaaSConnPilotDeployment}
	\begin{description}[noitemsep,nolistsep]
		\item[Responsibility:]~This node is responsible for the connection with the \vpett{\nameref{comp:ComponentsMLaaS}} Service.
		\item[Visible on diagrams:]~\cref{diag:Deployment:PilotDeployment200patients5clinmodels}		
	\end{description}

\subsection{MLaaS Service (Dev Deployment)}\label{node:DeploymentMLaaSServiceDevDeployment}
	\begin{description}[noitemsep,nolistsep]
		\item[Responsibility:]~This represents the \vpett{\nameref{comp:ComponentsMLaaS}} cloud service.
		\item[Visible on diagrams:]~\cref{diag:Deployment:DevelopmentTestDeployment20patients3clinmodels}		
	\end{description}

\subsection{MLaaS Service (Pilot Deployment)}\label{node:DeploymentMLaaSServicePilotDeployment}
	\begin{description}[noitemsep,nolistsep]
		\item[Responsibility:]~This represents the cloud \vpett{\nameref{comp:ComponentsMLaaS}} service.
		\item[Visible on diagrams:]~\cref{diag:Deployment:PilotDeployment200patients5clinmodels}		
	\end{description}

\subsection{MLaaSCloud}\label{node:DeploymentMLaaSCloud}
	\begin{description}[noitemsep,nolistsep]
		\item[Responsibility:]~Cloud infrastructure of the \vpett{\nameref{comp:ComponentsMLaaS}} Service provider.
		\item[Visible on diagrams:]~\cref{diag:Deployment:Contextdiagramforthedeploymentview}		
	\end{description}

\subsection{MLaaSConnectorNode}\label{node:DeploymenteHealthPlatformMLaaSConnectorNode}
	\begin{description}[noitemsep,nolistsep]
		\item[Responsibility:]~Node responsible for handling connections with external \vpett{\nameref{comp:ComponentsMLaaS}} providers.
		\item[Visible on diagrams:]~\cref{diag:Deployment:Contextdiagramforthedeploymentview,diag:Deployment:Primarydeploymentdiagram}		
	\end{description}

\subsection{MonitoringNode}\label{node:DeploymentMonitoringNode}
	\begin{description}[noitemsep,nolistsep]
		\item[Responsibility:]~Server node hosting the monitoring and recovery tier of the PMS backend: MonitoringService and RecoverySyncService. Responsible for detecting missing sensor data and internal failures within the Av2 detection deadlines, classifying failures by patient risk, dispatching notifications, recording communication gaps, and coordinating post-failure recovery.
		\item[Visible on diagrams:]~\cref{diag:Deployment:Primarydeploymentdiagram}		
	\end{description}

\subsection{MonitoringNode (Pilot Deployment)}\label{node:DeploymentMonitoringNodePilotDeployment}
	\begin{description}[noitemsep,nolistsep]
		\item[Responsibility:]~Pilot-deployment instance of the monitoring tier hosting MonitoringService and RecoverySyncService. Responsible for failure detection, classification, notification, and post-failure recovery coordination for the 200-patient pilot population.
		\item[Visible on diagrams:]~\cref{diag:Deployment:PilotDeployment200patients5clinmodels}		
	\end{description}

\subsection{PatientDataNode}\label{node:DeploymentPatientDataNode}
	\begin{description}[noitemsep,nolistsep]
		\item[Responsibility:]~Server node hosting the patient data service layer: PatientDataService (owns raw sensor data, patient status, and \vpett{\nameref{comp:ComponentsHIS}} EHR proxy), PatientStatusCache (read-through cache for fast status lookups), and HISAdapter (outbound adapter for the external \vpett{\nameref{comp:ComponentsHIS}} healthAPI). Bridges the backend ingestion tier and the external Hospital Information System.
		\item[Visible on diagrams:]~\cref{diag:Deployment:Primarydeploymentdiagram}		
	\end{description}

\subsection{PatientDataNode (Dev Deployment)}\label{node:DeploymentPatientDataNodeDevDeployment}
	\begin{description}[noitemsep,nolistsep]
		\item[Responsibility:]~Development-deployment instance of the patient data tier. HISAdapter points to a mock \vpett{\nameref{comp:ComponentsHIS}} stub; no real \vpett{\nameref{comp:ComponentsHIS}} connection.
		\item[Visible on diagrams:]~\cref{diag:Deployment:DevelopmentTestDeployment20patients3clinmodels}		
	\end{description}

\subsection{PatientDataNode (Pilot Deployment)}\label{node:DeploymentPatientDataNodePilotDeployment}
	\begin{description}[noitemsep,nolistsep]
		\item[Responsibility:]~Pilot-deployment instance of the patient data tier.
		\item[Visible on diagrams:]~\cref{diag:Deployment:PilotDeployment200patients5clinmodels}		
	\end{description}

\subsection{PatientGatewayDevice}\label{node:DeploymentPatientGatewayDevice}
	\begin{description}[noitemsep,nolistsep]
		\item[Responsibility:]~This node represents the physical wearable/handheld device carried by the patient.
		\item[Visible on diagrams:]~\cref{diag:Deployment:Contextdiagramforthedeploymentview}		
	\end{description}

\subsection{PatientRecordsNode}\label{node:DeploymentPatientRecordsNode}
	\begin{description}[noitemsep,nolistsep]
		\item[Responsibility:]~Server node hosting all patient-data persistence and identity components: \vpett{\nameref{comp:ComponentsPatientRecordRepository}}, \vpett{\nameref{comp:ComponentsPatientManagementService}}, \vpett{\nameref{comp:ComponentsPatientStatusService}}, \vpett{\nameref{comp:ComponentsUserAuthenticationService}}, and SensorDataRepository. Acts as the central data tier of the PMS backend.
		\item[Visible on diagrams:]~\cref{diag:Deployment:Primarydeploymentdiagram}		
	\end{description}

\subsection{PatientRecordsNode (Pilot Deployment)}\label{node:DeploymentPatientRecordsNodePilotDeployment}
	\begin{description}[noitemsep,nolistsep]
		\item[Responsibility:]~Pilot-deployment instance of the patient-data persistence tier. Single instance hosting \vpett{\nameref{comp:ComponentsPatientRecordRepository}} and SensorDataRepository (plus the management/status/authentication services) for the 200-patient pilot.
		\item[Visible on diagrams:]~\cref{diag:Deployment:PilotDeployment200patients5clinmodels}		
	\end{description}

\subsection{PhysicianAccessNode}\label{node:DeploymentPhysicianAccessNode}
	\begin{description}[noitemsep,nolistsep]
		\item[Responsibility:]~Server node hosting all physician-facing interface components: PhysicianGateway (the single external entry point for all physician API calls, with no business logic), PhysicianCommandService (handles physician write and command flows for UC7 configure risk assessment, UC8 update risk level, and UC9 on-demand consultation), PatientQueryService (handles UC6 patient-status reads with tiered SLA priority queue for high, medium, and low risk patients), and NotificationDispatcher (UC5 dispatcher that subscribes to risk events, maintains a red-yellow-green priority queue, and fans notifications out to registered physicians). Communicates with BackendGatewayNode to reach sensor data and gateway services, with PatientRecordsNode for patient status and record reads, and with RiskEstimationMgmtNode to launch risk estimations.
		\item[Visible on diagrams:]~\cref{diag:Deployment:Primarydeploymentdiagram}		
	\end{description}

\subsection{PhysicianNode (Pilot Deployment)}\label{node:DeploymentPhysicianNodePilotDeployment}
	\begin{description}[noitemsep,nolistsep]
		\item[Responsibility:]~Pilot-deployment instance of the physician interface tier. Single instance serving all physicians accessing the 200-patient pilot.
		\item[Visible on diagrams:]~\cref{diag:Deployment:PilotDeployment200patients5clinmodels}		
	\end{description}

\subsection{PhysicianWorkstation}\label{node:DeploymentPhysicianWorkstation}
	\begin{description}[noitemsep,nolistsep]
		\item[Responsibility:]~Client device used by physicians to access the PMS. Hosts the PhysicianGateway web or desktop client through which cardiologists and other registered physicians consult patient status (UC6), perform on-demand consultations (UC9), configure risk assessments (UC7), update risk levels (UC8), and receive notifications about patient status changes (UC5). Outside the PMS server-side deployment boundary.
		\item[Visible on diagrams:]~\cref{diag:Deployment:Contextdiagramforthedeploymentview}		
	\end{description}

\subsection{RiskEstimationManagementNode (Pilot Deployment)}\label{node:DeploymentRiskEstimationManagementNodePilotDeployment}
	\begin{description}[noitemsep,nolistsep]
		\item[Responsibility:]~The RiskEstimationManagement node hosts the creation, scheduling and combining of the clinical jobs.
		\item[Visible on diagrams:]~\cref{diag:Deployment:PilotDeployment200patients5clinmodels}		
	\end{description}

\subsection{RiskEstimationManagmentNode (Dev Deployment)}\label{node:DeploymentRiskEstimationManagmentNodeDevDeployment}
	\begin{description}[noitemsep,nolistsep]
		\item[Responsibility:]~The development deployment is a single RiskEstimationManagement node that contains all other functionality, job creation, scheduling, and combination.
		\item[Visible on diagrams:]~\cref{diag:Deployment:DevelopmentTestDeployment20patients3clinmodels}		
	\end{description}

\subsection{RiskEstimationMgmtNode}\label{node:DeploymenteHealthPlatformRiskEstimationMgmtNode}
	\begin{description}[noitemsep,nolistsep]
		\item[Responsibility:]~The RiskEstimationMgmtNode keeps track of the risk estimation jobs, and schedules and combines the results.
		\item[Visible on diagrams:]~\cref{diag:Deployment:Primarydeploymentdiagram}		
	\end{description}

\subsection{RiskEstimationProcessorNode}\label{node:DeploymenteHealthPlatformRiskEstimationProcessorNode}
	\begin{description}[noitemsep,nolistsep]
		\item[Responsibility:]~Multiple \vpett{\nameref{comp:ComponentseHealthPlatformRiskEstimationProcessor}} nodes enable the processing of multiple sensorData inputs in parallel.
		\item[Visible on diagrams:]~\cref{diag:Deployment:Primarydeploymentdiagram}		
	\end{description}

\subsection{RiskEstimationProcessorNode (Dev Deployment)}\label{node:DeploymentRiskEstimationProcessorNodeDevDeployment}
	\begin{description}[noitemsep,nolistsep]
		\item[Responsibility:]~The development deployment only uses two RiskEstimationProcessorNodes to run \ref{data:DatatypesClinicalModelJob}s in parallel.
		\item[Visible on diagrams:]~\cref{diag:Deployment:DevelopmentTestDeployment20patients3clinmodels}		
	\end{description}

\subsection{RiskEstimationProcessorNode (Pilot Deployment)}\label{node:DeploymentRiskEstimationProcessorNodePilotDeployment}
	\begin{description}[noitemsep,nolistsep]
		\item[Responsibility:]~There are five instances of the RiskEstimationProcessorNode to enable the processing of multiple \ref{data:DatatypesClinicalModelJob}s in parallel.
		\item[Visible on diagrams:]~\cref{diag:Deployment:PilotDeployment200patients5clinmodels}		
	\end{description}

\subsection{SMSGatewayProvider}\label{node:DeploymentSMSGatewayProvider}
	\begin{description}[noitemsep,nolistsep]
		\item[Responsibility:]~Represents the third-party SMS service used as the backup channel.
		\item[Visible on diagrams:]~\cref{diag:Deployment:Contextdiagramforthedeploymentview}		
	\end{description}

\subsection{TelecomNetwork}\label{node:DeploymentTelecomNetwork}
	\begin{description}[noitemsep,nolistsep]
		\item[Responsibility:]~represents the cellular network governed by the SLA (99.5\% / 80\% / 64Kb/s).
		\item[Visible on diagrams:]~\cref{diag:Deployment:Contextdiagramforthedeploymentview}		
	\end{description}

\subsection{TODO Node (Pilot Deployment)}\label{node:DeploymentTODONodePilotDeployment}
	\begin{description}[noitemsep,nolistsep]
		\item[Responsibility:]~This TODO node for the pilot deployment contains the \vpett{\nameref{comp:ComponentseHealthPlatformOtherFunctionality}} of the system that has not been worked out in detail yet.
		\item[Visible on diagrams:]~\cref{diag:Deployment:PilotDeployment200patients5clinmodels}		
	\end{description}

\subsection{TODONode}\label{node:DeploymenteHealthPlatformTODONode}
	\begin{description}[noitemsep,nolistsep]
		\item[Responsibility:]~Temporary node that requires further decomposition.
		\item[Visible on diagrams:]~\cref{diag:Deployment:Primarydeploymentdiagram}		
	\end{description}


% END NODES

% EXCEPTIONS
\section{Exceptions}\label{sec:exceptions}
\begin{itemize}[nolistsep,noitemsep]
\item \vpeexception{ex:ExceptionsAuthenticationException}{AuthenticationException} (\iconinher{}~\ref{ex:ExceptionsException}): 
\begin{itemize}[noitemsep,nolistsep]

\item[] Signals a problem with the credentials used for authentication. These may be invalid or expired.
\end{itemize}

\item \vpeexception{ex:ExceptionsAuthenticationException2}{AuthenticationException2} (\iconinher{}~\ref{ex:ExceptionsSecurityException}): 
\begin{itemize}[noitemsep,nolistsep]

\item[] Missing or invalid credentials.
\end{itemize}

\item \vpeexception{ex:ExceptionsAuthorizationException}{AuthorizationException} (\iconinher{}~\ref{ex:ExceptionsException}): 
\begin{itemize}[noitemsep,nolistsep]

\item[] Signals that the principal has insufficient access rights to perform the operation.
\end{itemize}

\item \vpeexception{ex:ExceptionsAuthorizationException2}{AuthorizationException2} (\iconinher{}~\ref{ex:ExceptionsSecurityException}): 
\begin{itemize}[noitemsep,nolistsep]

\item[] The requestor is not authorized to perform the requested action.
\end{itemize}

\item \vpeexception{ex:ExceptionsException}{Exception}: 
\begin{itemize}[noitemsep,nolistsep]
\item[] Subtypes: \iconinher{}~\ref{ex:ExceptionsAuthenticationException}, \iconinher{}~\ref{ex:ExceptionsAuthorizationException}, \iconinher{}~\ref{ex:ExceptionsInvalidAPIKeyException}, \iconinher{}~\ref{ex:ExceptionsInvalidMLModelException}, \iconinher{}~\ref{ex:ExceptionsIOException}, \iconinher{}~\ref{ex:ExceptionsNoSuchMLModelException}, \iconinher{}~\ref{ex:ExceptionsNoSuchPatientException}, \iconinher{}~\ref{ex:ExceptionsNoSuchPatientRecordException}
\item[] Base exception class.
\end{itemize}

\item \vpeexception{ex:ExceptionsInvalidAPIKeyException}{InvalidAPIKeyException} (\iconinher{}~\ref{ex:ExceptionsException}): 
\begin{itemize}[noitemsep,nolistsep]

\item[] Thrown when no correctly formatted API key is configured for constructing authentication Credential|s.
\end{itemize}

\item \vpeexception{ex:ExceptionsInvalidMLModelException}{InvalidMLModelException} (\iconinher{}~\ref{ex:ExceptionsException}): 
\begin{itemize}[noitemsep,nolistsep]

\item[] Thrown when a provided MLModel is invalid.
\end{itemize}

\item \vpeexception{ex:ExceptionsIOException}{IOException} (\iconinher{}~\ref{ex:ExceptionsException}): 
\begin{itemize}[noitemsep,nolistsep]

\item[] Thrown when an IO operation fails.
\end{itemize}

\item \vpeexception{ex:ExceptionsNoSuchMLModelException}{NoSuchMLModelException} (\iconinher{}~\ref{ex:ExceptionsException}): 
\begin{itemize}[noitemsep,nolistsep]

\item[] Thrown when the requested MLModel cannot be found.
\end{itemize}

\item \vpeexception{ex:ExceptionsNoSuchPatientException}{NoSuchPatientException} (\iconinher{}~\ref{ex:ExceptionsException}): 
\begin{itemize}[noitemsep,nolistsep]

\item[] Thrown if no patient with the specified identifier exists.
\end{itemize}

\item \vpeexception{ex:ExceptionsNoSuchPatientRecordException}{NoSuchPatientRecordException} (\iconinher{}~\ref{ex:ExceptionsException}): 
\begin{itemize}[noitemsep,nolistsep]

\item[] Thrown if no patient record for the patient with the given identifier exists in the HIS.
\end{itemize}

\item \vpeexception{ex:ExceptionsSecurityException}{SecurityException}: 
\begin{itemize}[noitemsep,nolistsep]
\item[] Subtypes: \iconinher{}~\ref{ex:ExceptionsAuthenticationException2}, \iconinher{}~\ref{ex:ExceptionsAuthorizationException2}
\item[] Parent class for security exceptions.
\end{itemize}

\end{itemize}
% END EXCEPTIONS

% DATA TYPES
\section{Data types}\label{sec:datatypes}
\begin{itemize}[nolistsep,noitemsep]
\item \vpedatatype{data:Datatypesorghl7fhirr5modelAddress}{Address}: 
\begin{itemize}[noitemsep,nolistsep]
\item[] Attributes: \ref{data:Datatypesorghl7fhirr5modelAddressAddressUse} use, \ref{data:Datatypesorghl7fhirr5modelAddressAddressType} type, String text, List\textless{}String\textgreater{} line, String city, String district, String state, String postalCode, String country, \ref{data:Datatypesorghl7fhirr5modelPeriod} period, long serialVersionUID
\item[] Base StructureDefinition for Address Type: An address expressed using postal conventions (as opposed to GPS or other location definition formats).  This data type may be used to convey addresses for use in delivering mail as well as for visiting locations which might not be valid for mail delivery.  There are a variety of postal address formats defined around the world.
The line attribute contains contains the house number, apartment number, street name, street direction, P.O.
The text attribute specifies the entire address as it should be displayed e.g. on a postal label. This may be provided instead of or as well as the specific parts.
\end{itemize}

\item \vpedatatype{data:Datatypesorghl7fhirr5modelAddressAddressType}{AddressType}: 
\begin{itemize}[noitemsep,nolistsep]

\item[] Enum: \textsc{postal}, \textsc{physical}, \textsc{both}, \textsc{null}.\\Distinguishes between physical addresses (those you can visit) and mailing addresses (e.g. PO Boxes and care-of addresses). Most addresses are both.
\end{itemize}

\item \vpedatatype{data:Datatypesorghl7fhirr5modelAddressAddressUse}{AddressUse}: 
\begin{itemize}[noitemsep,nolistsep]

\item[] Enum: \textsc{home}, \textsc{work}, \textsc{temp}, \textsc{old}, \textsc{billing}, \textsc{null}.\\The purpose of the address.
\end{itemize}

\item \vpedatatype{data:Datatypesorghl7fhirr5modelAnnotation}{Annotation}: 
\begin{itemize}[noitemsep,nolistsep]
\item[] Attributes: \ref{data:DatatypesDate} time, String text, long serialVersionUID, String author
\item[] Base StructureDefinition for Annotation Type: A  text note which also  contains information about who made the statement and when.
\end{itemize}

\item \vpedatatype{data:DatatypesbackupChannelType}{backupChannelType}: 
\begin{itemize}[noitemsep,nolistsep]

\item[] Enum: SMS.
  The type of backup communication channel available to the 
  patient gateway when the primary channel is unavailable.
\end{itemize}

\item \vpedatatype{data:DatatypesBackupDeliveryResult}{BackupDeliveryResult}: 
\begin{itemize}[noitemsep,nolistsep]
\item[] Attributes: boolean succes, \ref{data:DatatypesTimestamp} sendAt, string channelUsed
\item[] The result of an attempt to send a message via the backup 
  communication channel.
\end{itemize}

\item \vpedatatype{data:DatatypesBigDecimal}{BigDecimal}: 
\begin{itemize}[noitemsep,nolistsep]

\item[] Immutable, arbitrary-precision signed decimal numbers. A BigDecimal consists of an arbitrary precision integer unscaled value and a 32-bit integer scale
\end{itemize}

\item \vpedatatype{data:DatatypesClassifiedFailure}{ClassifiedFailure}: 
\begin{itemize}[noitemsep,nolistsep]
\item[] Attributes: \ref{data:DatatypesFailureEvent} failure, \ref{data:DatatypesEscalationLevel} level, \ref{data:DatatypesTimestamp} notificationDeadline
\item[] A failure event enriched with its escalation level and the 
  deadline by which the relevant stakeholder must be notified, 
  derived from the affected patient risk level per Av2 response 
  measures.
\end{itemize}

\item \vpedatatype{data:DatatypesClinicalModel}{ClinicalModel}: 
\begin{itemize}[noitemsep,nolistsep]
\item[] Attributes: String id, String name, String description, List\textless{}String\textgreater{} requiredParameters, List\textless{}String\textgreater{} optionalParameters, List\textless{}String\textgreater{} applicableMLModelIds, boolean requiresHistoricalSensorData
\item[] A clinical model. This is model has an id, name, description.
It contains a list of  parameters for using the model and a set of MLModelIds to specify which types of MLModels can be used in the context of the ClinicalModel.

\end{itemize}

\item \vpedatatype{data:DatatypesClinicalModelConfiguration}{ClinicalModelConfiguration}: 
\begin{itemize}[noitemsep,nolistsep]
\item[] Attributes: Map\textless{}String, String\textgreater{} modelParameters
\item[] The configuration for a clinical model. Represented as a set key/value pair for the parameters of the clinical model.
\end{itemize}

\item \vpedatatype{data:DatatypesClinicalModelJob}{ClinicalModelJob}: 
\begin{itemize}[noitemsep,nolistsep]
\item[] Attributes: \ref{data:DatatypesClinicalModel} clinicalModel, \ref{data:DatatypesClinicalModelJobId} jobId, \ref{data:DatatypesSensorDataPackage} sensorData, \ref{data:DatatypesPatientId} patientId, \ref{data:DatatypesClinicalModelConfiguration} clinicalModelConfiguration, \ref{data:DatatypesPatientRecord} patientRecord, Map\textless{}\ref{data:DatatypesTimestamp}, \ref{data:DatatypesSensorDataPackage}\textgreater{} historicalSensorData
\item[] A single job for the computation of a clinical model. This contains a ClinicalModelJobId, ClinicalModelId and PatientId respectively identifying the job itself, the corresponding clinical model and the patient. If the corresponding risk estimation the job belongs to is triggered by the arrival of a sensor data package, this sensor data is also contained.
\end{itemize}

\item \vpedatatype{data:DatatypesClinicalModelJobId}{ClinicalModelJobId}: 
\begin{itemize}[noitemsep,nolistsep]

\item[]  A piece of data uniquely identifying a certain clinical model job in the system. This architecture does not specify the exact format of this identifier,
but possibilities are a long integer, a string, a URI etc.
\end{itemize}

\item \vpedatatype{data:DatatypesClinicalModelJobResult}{ClinicalModelJobResult}: 
\begin{itemize}[noitemsep,nolistsep]

\item[] The result of the computation of a clinical model containing the estimated
risk level for the patient.
The resulting risk score for ClinicalModels are on the same unit scale so the the results of multiple ClinicalModels can be combined.
\end{itemize}

\item \vpedatatype{data:Datatypesorghl7fhirr5modelCodeableConcept}{CodeableConcept}: 
\begin{itemize}[noitemsep,nolistsep]
\item[] Attributes: List\textless{}\ref{data:Datatypesorghl7fhirr5modelCoding}\textgreater{} coding, String text, long serialVersionUID
\item[] Base StructureDefinition for CodeableConcept Type: A concept that may be defined by a formal reference to a terminology or ontology or may be provided by text.
\end{itemize}

\item \vpedatatype{data:Datatypesorghl7fhirr5modelCodeType}{CodeType}: 
\begin{itemize}[noitemsep,nolistsep]
\item[] Attributes: long serialVersionUID, String system, String value
\item[] Primitive type "code" in FHIR, when not bound to an enumerated list of codes
\end{itemize}

\item \vpedatatype{data:Datatypesorghl7fhirr5modelCoding}{Coding}: 
\begin{itemize}[noitemsep,nolistsep]
\item[] Attributes: String system, String version, \ref{data:Datatypesorghl7fhirr5modelCodeType} code, String display, boolean userSelected, long serialVersionUID
\item[] Base StructureDefinition for Coding Type: A reference to a code defined by a terminology system.
\end{itemize}

\item \vpedatatype{data:DatatypesConsultationId}{ConsultationId}: 
\begin{itemize}[noitemsep,nolistsep]

\item[] TODO maybe delete/merge into CorrelationId

\end{itemize}

\item \vpedatatype{data:Datatypesorghl7fhirr5modelPatientContactComponent}{ContactComponent}: 
\begin{itemize}[noitemsep,nolistsep]
\item[] Attributes: List\textless{}\ref{data:Datatypesorghl7fhirr5modelCodeableConcept}\textgreater{} relationship, String name, List\textless{}\ref{data:Datatypesorghl7fhirr5modelContactPoint}\textgreater{} telecom, \ref{data:Datatypesorghl7fhirr5modelAddress} address, String administrativeGender, String organization, \ref{data:Datatypesorghl7fhirr5modelPeriod} period, long serialVersionUID
\item[] A contact party (e.g. guardian, partner, friend) for the patient.
The relationship attribute species the nature of the relationship between the patient and the contact person.
It is a CodeableConcept (which can be formal reference to terminology or an ontology).
For example, in the system \url{http://hl7.org/fhir/ValueSet/relatedperson-relationshiptype}, it can be coded as CHILD for a child, or MGRMTH for maternal grandmother, etc.
\end{itemize}

\item \vpedatatype{data:Datatypesorghl7fhirr5modelContactPoint}{ContactPoint}: 
\begin{itemize}[noitemsep,nolistsep]
\item[] Attributes: \ref{data:Datatypesorghl7fhirr5modelContactPointContactPointSystem} system, String value, \ref{data:Datatypesorghl7fhirr5modelContactPointContactPointUse} use, int rank, long serialVersionUID, \ref{data:Datatypesorghl7fhirr5modelPeriod} period
\item[] Base StructureDefinition for ContactPoint Type: Details for all kinds of technology mediated contact points for a person or organization, including telephone, email, etc.
The value attribute contains the actual contact point details, in a form that is meaningful to the designated communication system (i.e. phone number or email address).
Rank specifies a preferred order in which to use a set of contacts. ContactPoints with lower rank values are more preferred than those with higher rank values.
Period specifies the tiime period when the contact point was/is in use.

\end{itemize}

\item \vpedatatype{data:Datatypesorghl7fhirr5modelContactPointContactPointSystem}{ContactPointSystem}: 
\begin{itemize}[noitemsep,nolistsep]

\item[] Enum: \textsc{phone}, \textsc{fax}, \textsc{email}, \textsc{pager}, \textsc{url}, \textsc{sms}, \textsc{other}, \textsc{null}.\\Enum denotying the type of communication channel for the ContactPoint.
\end{itemize}

\item \vpedatatype{data:Datatypesorghl7fhirr5modelContactPointContactPointUse}{ContactPointUse}: 
\begin{itemize}[noitemsep,nolistsep]

\item[] Enum: \textsc{home}, \textsc{work}, \textsc{temp}, \textsc{old}, \textsc{mobile}, \textsc{null}.\\Enum denoting the usage type of the ContactPoint.
\end{itemize}

\item \vpedatatype{data:DatatypesCorrelationId}{CorrelationId}: 
\begin{itemize}[noitemsep,nolistsep]

\item[] Cross-component tag that ties a UC9 on-demand consultation to its eventual risk-estimation result.
\end{itemize}

\item \vpedatatype{data:DatatypesCredential}{Credential}: 
\begin{itemize}[noitemsep,nolistsep]

\item[] Authentication object for making authenticated requests to the MLaaS service provider.
\end{itemize}

\item \vpedatatype{data:DatatypesDate}{Date}: 
\begin{itemize}[noitemsep,nolistsep]

\item[] The class Date represents a specific instant in time, with millisecond precision. 
\end{itemize}

\item \vpedatatype{data:DatatypesDegradedModeReason}{DegradedModeReason}: 
\begin{itemize}[noitemsep,nolistsep]

\item[] Enum: NO\_NETWORK, SERVER\_UNREACHABLE, INTERNAL\_FAILURE.
The reason the patient gateway entered degraded mode.
\end{itemize}

\item \vpedatatype{data:DatatypesDuration}{Duration}: 
\begin{itemize}[noitemsep,nolistsep]

\item[] Represents a length of time.
\end{itemize}

\item \vpedatatype{data:DatatypesEmergencyMessage}{EmergencyMessage}: 
\begin{itemize}[noitemsep,nolistsep]
\item[] Attributes: \ref{data:DatatypesPatientId} patientId, \ref{data:DatatypesGatewayId} gatewayId, {\colorbox{red!30}{\underline{EmergencyType}}} type, \ref{data:DatatypesTimestamp} detectedAt, string payload
\item[] An emergency alert to be dispatched via the backup channel 
  when the primary communication channel is unavailable.
\end{itemize}

\item \vpedatatype{data:DatatypesEscalationLevel}{EscalationLevel}: 
\begin{itemize}[noitemsep,nolistsep]

\item[] Enum: HIGH, MEDIUM, LOW.
The urgency level of a detected failure, derived from the risk level of the affected patient. HIGH corresponds to high-risk patients (notify within 2 minutes), MEDIUM to medium-risk patients (notify within 5 minutes), LOW to low-risk patients (notify within 10 minutes).
\end{itemize}

\item \vpedatatype{data:DatatypesFailureEvent}{FailureEvent}: 
\begin{itemize}[noitemsep,nolistsep]
\item[] Attributes: \ref{data:DatatypesFailureEventId} id, \ref{data:DatatypesFailureType} type, \ref{data:DatatypesGatewayId} gatewayId, int subsystemId, \ref{data:DatatypesTimestamp} detectedAt, \ref{data:DatatypesPatientId} patientId
\item[] Represents a detected failure in the system. The gatewayId 
  is present when the failure concerns a specific patient gateway, 
  and may be null for internal subsystem failures. Similarly 
  patientId is null when the failure does not concern a specific 
  patient. The detectedAt timestamp records when the system became 
  aware of the failure, while startedAt estimates when the failure 
  actually began.
\end{itemize}

\item \vpedatatype{data:DatatypesFailureEventId}{FailureEventId}: 
\begin{itemize}[noitemsep,nolistsep]

\item[] A unique identifier for a failure event.
\end{itemize}

\item \vpedatatype{data:DatatypesFailureType}{FailureType}: 
\begin{itemize}[noitemsep,nolistsep]

\item[] Enum: MISSING\_SENSOR\_DATA, INTERNAL\_SUBSYSTEM\_CRASH, 
        COMMUNICATION\_CHANNEL\_DOWN, GATEWAY\_UNREACHABLE.
  Classifies the kind of failure detected by the MonitoringService.
\end{itemize}

\item \vpedatatype{data:DatatypesFilterCriteria}{FilterCriteria}: 
\begin{itemize}[noitemsep,nolistsep]
\item[] Attributes: \ref{data:DatatypesCorrelationId} correlationId
\item[] Optional filter passed to IRiskEvents.subscribe. Empty for NotificationDispatcher (wants every event); carries a correlationId for PhysicianCommandService UC9 result lookup.
\end{itemize}

\item \vpedatatype{data:DatatypesGatewayId}{GatewayId}: 
\begin{itemize}[noitemsep,nolistsep]

\item[] A piece of data uniquely identifying a patient gateway device. 
\end{itemize}

\item \vpedatatype{data:DatatypesHealthStatus}{HealthStatus}: 
\begin{itemize}[noitemsep,nolistsep]

\item[] Enum: OK, DEGRADED, DOWN.
  The operational status of an internal subsystem as reported 
  by its health check endpoint.
\end{itemize}

\item \vpedatatype{data:Datatypesorghl7fhirr5modelIdentifier}{Identifier}: 
\begin{itemize}[noitemsep,nolistsep]
\item[] Attributes: \ref{data:Datatypesorghl7fhirr5modelIdentifierIdentifierUse} use, String type, String system, String value, \ref{data:Datatypesorghl7fhirr5modelPeriod} period, String assigner, long serialVersionUID
\item[] Base StructureDefinition for Identifier Type: An identifier - identifies some entity uniquely and unambiguously. Typically this is used for business identifiers.
The type is a a coded type for the identifier that can be used to determine which identifier to use for a specific purpose.
The system attribute establishes the namespace for the value - that is, a URL that describes a set values that are unique.
The value attribute specifies the portion of the identifier typically relevant to the user and which is unique within the context of the system.
The period specifies the time period during which identifier is/was valid for use.
The assigner specifies the organization that issued/manages the identifier.
\end{itemize}

\item \vpedatatype{data:Datatypesorghl7fhirr5modelIdentifierIdentifierUse}{IdentifierUse}: 
\begin{itemize}[noitemsep,nolistsep]

\item[] Enum: \textsc{usual}, \textsc{official}, \textsc{temp}, \textsc{secondary}, \textsc{old}, \textsc{null}.\\Enum denoting the usage type of a particular Identifier.
\end{itemize}

\item \vpedatatype{data:DatatypesMLModel}{MLModel}: 
\begin{itemize}[noitemsep,nolistsep]
\item[] Attributes: String modelId, \ref{data:DatatypesMLModelDescription} modelDescription
\item[] Machine learning model for making predictions based on sensordata.
The actual internal encoding of the model itself is left open.
\end{itemize}

\item \vpedatatype{data:DatatypesMLModelDescription}{MLModelDescription}: 
\begin{itemize}[noitemsep,nolistsep]
\item[] Attributes: String uuid, String name, String description, String providerName
\item[] Object providing the MLModel meta data. It does not contain the model itself.
\end{itemize}

\item \vpedatatype{data:DatatypesMLModelUpdate}{MLModelUpdate}: 
\begin{itemize}[noitemsep,nolistsep]
\item[] Attributes: \ref{data:DatatypesMLModelDescription} modelDescription
\item[] Represents a model update that can be applied on an MLModel.
These model updates enable exchanging changes to the MLModel between different parties without sending the actual MLModel to the other parties.
\end{itemize}

\item \vpedatatype{data:DatatypesNetworkType}{NetworkType}: 
\begin{itemize}[noitemsep,nolistsep]

\item[] NetworkType: Enum: NONE, CELLULAR, WIFI.
The type of network connection currently available to the patient gateway. CELLULAR refersto mobile network connectivity governed by the SLA with the telecom operator.
\end{itemize}

\item \vpedatatype{data:DatatypesNotification}{Notification}: 
\begin{itemize}[noitemsep,nolistsep]
\item[] Attributes: \ref{data:DatatypesNotificationId} id, \ref{data:DatatypesPatientId} patientId, \ref{data:DatatypesPatientStatus} status, \ref{data:DatatypesNotificationSeverity} severity, \ref{data:DatatypesTimestamp} ts
\item[] UC5 payload returned to physicians via INotificationInbox.getPendingNotifications.
\end{itemize}

\item \vpedatatype{data:DatatypesNotificationId}{NotificationId}: 
\begin{itemize}[noitemsep,nolistsep]

\item[] Identifies one queued notification (used for ack/drop).
\end{itemize}

\item \vpedatatype{data:DatatypesNotificationSeverity}{NotificationSeverity}: 
\begin{itemize}[noitemsep,nolistsep]

\item[] Enum: \textsc{red}, \textsc{yellow}, \textsc{green}, \textsc{literal}.\\Severity of a Notification. Drives the dispatcher's priority queue + drop policy (Response measure 5: red \textgreater{} yellow \textgreater{} green; green dropped first under overload).
\end{itemize}

\item \vpedatatype{data:DatatypesObject}{Object}: 
\begin{itemize}[noitemsep,nolistsep]

\item[] Root class for objects that can be stored in the Key-Value storage service.
Data types that need to be saved need to inherit from this class.
\end{itemize}

\item \vpedatatype{data:Datatypesorghl7fhirr5modelObservation}{Observation}: 
\begin{itemize}[noitemsep,nolistsep]
\item[] Attributes: List\textless{}\ref{data:Datatypesorghl7fhirr5modelIdentifier}\textgreater{} identifier, List\textless{}\ref{data:Datatypesorghl7fhirr5modelIdentifier}\textgreater{} partOf, \ref{data:Datatypesorghl7fhirr5modelObservationStatus} status, List\textless{}\ref{data:Datatypesorghl7fhirr5modelCodeableConcept}\textgreater{} category, \ref{data:Datatypesorghl7fhirr5modelCodeableConcept} code, \ref{data:Datatypesorghl7fhirr5modelPatient} subject, \ref{data:Datatypesorghl7fhirr5modelPeriod} effective, \ref{data:DatatypesDate} issued, List\textless{}\ref{data:Datatypesorghl7fhirr5modelIdentifier}\textgreater{} performer, String value, \ref{data:Datatypesorghl7fhirr5modelCodeableConcept} dataAbsentReason, \ref{data:Datatypesorghl7fhirr5modelCodeableConcept} interpretation, List\textless{}\ref{data:Datatypesorghl7fhirr5modelAnnotation}\textgreater{} note, \ref{data:Datatypesorghl7fhirr5modelCodeableConcept} bodySite, \ref{data:Datatypesorghl7fhirr5modelCodeableConcept} method, String device, List\textless{}\ref{data:Datatypesorghl7fhirr5modelObservation}\textgreater{} hasMember, List\textless{}\ref{data:Datatypesorghl7fhirr5modelObservation}\textgreater{} derivedFrom, long serialVersionUID, \ref{data:Datatypesorghl7fhirr5modelObservationObservationReferenceRangeComponent} referenceRange, \ref{data:Datatypesorghl7fhirr5modelObservationObservationComponentComponent} component
\item[] Measurements and simple assertions made about a patient, device or other subject.
The partOf attribute specifies the larger event of which this particular Observation is a component or step. For example, an observation as part of a procedure.

The category specifies a code that classifies the general type of observation being made.
For example, in the system \url{http://terminology.hl7.org/CodeSystem/observation-category}, vital-signs is used for clinical observations such as boold pressure, heart rate, etc.

The code attribute describes what was observed. Sometimes this is called the observation `name'.
For example, in the system \url{http://loinc.org}, code `8867-4' is used for expressing the heart rate, `85354-9' for a blood pressure panel (which must have components `8480-6' and `8462-4' for, respectively, systolic and diastolic pressure.

The performer is who was responsible for asserting the observed value as `true'. This can refer to a Patient, Practitioner, organization, etc.
The value is the information determined as a result of making the observation, if the information has a simple value.

The dataAbsentReason provides a reason why the expected value in the element Observation.value[x] is missing.
For example, in the system \url{http://terminology.hl7.org/CodeSystem/data-absent-reason}, `not-applicable'.

The interpretation is a categorical assessment of an observation value.  For example, high, low, normal.
For example, in the system \url{http://terminology.hl7.org/CodeSystem/v3-ObservationInterpretation}, `H' for high, or `ND' for not detected.

The bodySite indicates the site on the subject's body where the observation was made (i.e. the target site).
For example, in the system \url{http://snomed.info/sct}, code `344001' indicates the ankle.

The method indicates the mechanism used to perform the observation.
For example, in the system \url{http://snomed.info/sct}, code `46973005' indicates blood pressure taking procedure.

The device attribute specifies an identifier of the device used to generate the observation data.
\end{itemize}

\item \vpedatatype{data:Datatypesorghl7fhirr5modelObservationObservationComponentComponent}{ObservationComponentComponent}: 
\begin{itemize}[noitemsep,nolistsep]
\item[] Attributes: String code, \ref{data:Datatypesorghl7fhirr5modelRange} value, String dataAbsentReason, List\textless{}String\textgreater{} interpretation, long serialVersionUID, \ref{data:Datatypesorghl7fhirr5modelObservationObservationReferenceRangeComponent} referenceRange
\item[] Some observations have multiple component observations.  These component observations are expressed as separate code value pairs that share the same attributes.  Examples include systolic and diastolic component observations for blood pressure measurement and multiple component observations for genetics observations.
The code describes what was observed. Sometimes this is called the observation `code'.
For example, in the system \url{http://loinc.org}, `8480-6' is used for the systolic component of the blood pressure.,

The dataAbsentReason provides a reason why the expected value in the element Observation.component.value[x] is missing.
For example, in the system \url{http://terminology.hl7.org/CodeSystem/data-absent-reason}, `not-applicable'.

The interpretation provides a categorical assessment of an observation value. For example, high, low, normal.
For example, in the system \url{http://terminology.hl7.org/CodeSystem/v3-ObservationInterpretation}, `H' for high, or `ND' for not detected.
\end{itemize}

\item \vpedatatype{data:Datatypesorghl7fhirr5modelObservationObservationReferenceRangeComponent}{ObservationReferenceRangeComponent}: 
\begin{itemize}[noitemsep,nolistsep]
\item[] Attributes: \ref{data:Datatypesorghl7fhirr5modelQuantity} low, \ref{data:Datatypesorghl7fhirr5modelQuantity} high, \ref{data:Datatypesorghl7fhirr5modelCodeableConcept} type, List\textless{}\ref{data:Datatypesorghl7fhirr5modelCodeableConcept}\textgreater{} appliesTo, \ref{data:Datatypesorghl7fhirr5modelRange} age, String text, long serialVersionUID
\item[] Guidance on how to interpret the value by comparison to a normal or recommended range.  Multiple reference ranges are interpreted as an "OR".   In other words, to represent two distinct target populations, two `referenceRange` elements would be used.

The appliesTo specifies codes to indicate the target population this reference range applies to.
For example, in the system \url{http://snomed.info/sct}, code `248152002' indicates female.

The type specifies codes to indicate the what part of the targeted reference population it applies to.
For example, in the system \url{http://terminology.hl7.org/CodeSystem/referencerange-meaning}, `normal' for the 95\% normal range, or `recommended' for the recommended range by a relevant professional body.

The text attribute contains a text based reference range in an observation which may be used when a quantitative range is not appropriate for an observation.
\end{itemize}

\item \vpedatatype{data:Datatypesorghl7fhirr5modelObservationStatus}{ObservationStatus}: 
\begin{itemize}[noitemsep,nolistsep]

\item[] Enum: \textsc{registered}, \textsc{preliminary}, \textsc{final}, \textsc{amended}, \textsc{corrected}, \textsc{cancelled}, \textsc{enteredinerror}, \textsc{unknown}, \textsc{null}.\\The status of the observation result value or the RiskAssessment.
\end{itemize}

\item \vpedatatype{data:Datatypesorghl7fhirr5modelPatient}{Patient}: 
\begin{itemize}[noitemsep,nolistsep]
\item[] Attributes: List\textless{}\ref{data:Datatypesorghl7fhirr5modelIdentifier}\textgreater{} identifier, boolean active, List\textless{}String\textgreater{} name, String administrativeGender, \ref{data:DatatypesDate} birthDate, \ref{data:DatatypesDate} deceased, List\textless{}\ref{data:Datatypesorghl7fhirr5modelAddress}\textgreater{} address, boolean multipleBirth, List\textless{}\ref{data:Datatypesorghl7fhirr5modelIdentifier}\textgreater{} generalPractitioner, \ref{data:Datatypesorghl7fhirr5modelIdentifier} managingOrganization, long serialVersionUID, \ref{data:Datatypesorghl7fhirr5modelPatientContactComponent} contact
\item[] Demographics and other administrative information about an individual or animal receiving care or other health-related services.
\end{itemize}

\item \vpedatatype{data:DatatypesPatientId}{PatientId}: 
\begin{itemize}[noitemsep,nolistsep]

\item[] A piece of data uniquely identifying a patient in the system. This architecture does not
specify the exact format of this identifier, but possibilities are a long integer, a string, a URI, etc.
\end{itemize}

\item \vpedatatype{data:DatatypesPatientRecord}{PatientRecord}: 
\begin{itemize}[noitemsep,nolistsep]

\item[] The structured contents of the EHR record of a certain patient, including medication history, recent treatments, allergies, etc.
\end{itemize}

\item \vpedatatype{data:DatatypesPatientStatus}{PatientStatus}: 
\begin{itemize}[noitemsep,nolistsep]

\item[] The status of a patient, i.e., their risk level.
\end{itemize}

\item \vpedatatype{data:Datatypesorghl7fhirr5modelPeriod}{Period}: 
\begin{itemize}[noitemsep,nolistsep]
\item[] Attributes: \ref{data:DatatypesDate} start, \ref{data:DatatypesDate} end, long serialVersionUID
\item[] Base StructureDefinition for Period Type: A time period defined by a start and end date and optionally time.
\end{itemize}

\item \vpedatatype{data:DatatypesPhysicianId}{PhysicianId}: 
\begin{itemize}[noitemsep,nolistsep]

\item[] Identifies a physician for notification routing.
\end{itemize}

\item \vpedatatype{data:DatatypesPrediction}{Prediction}: 
\begin{itemize}[noitemsep,nolistsep]

\item[] A prediction from the MLModel.
\end{itemize}

\item \vpedatatype{data:DatatypesPriority}{Priority}: 
\begin{itemize}[noitemsep,nolistsep]

\item[] Enum: \textsc{high}, \textsc{normal}.\\Priority of a launched risk-estimation job. UC9 uses `HIGH`; `NORMAL` is the default for legacy schedule-driven jobs.
\end{itemize}

\item \vpedatatype{data:Datatypesorghl7fhirr5modelQuantity}{Quantity}: 
\begin{itemize}[noitemsep,nolistsep]
\item[] Attributes: \ref{data:DatatypesBigDecimal} value, String unit, String system, \ref{data:Datatypesorghl7fhirr5modelCodeType} code, long serialVersionUID
\item[] Base StructureDefinition for Quantity Type: A measured amount (or an amount that can potentially be measured). Note that measured amounts include amounts that are not precisely quantified, including amounts involving arbitrary units and floating currencies.
Unit specifies a human-readable form of the unit.
System contains the identification of the system that provides the coded form of the unit.
The code contains a computer processable form of the unit in some unit representation system.
For example, blood pressure can be expressed in the system \url{http://unitsofmeasure.org} and the unit mm[Hg].
\end{itemize}

\item \vpedatatype{data:DatatypesQuery}{Query}\textbf{\textless{}T\textgreater{}}: 
\begin{itemize}[noitemsep,nolistsep]
\item[] Attributes: List\textless{}T\textgreater{} values, List\textless{}String\textgreater{} qualifier
\item[] The Query object is used to express search restrictions in a string-based format for interacting with the REST-based health API.
For example, before a specified date is encoded as `ltyyyy-MM-dd'.
The list of qualifiers specify how the corresponding value should be used (e.g., lt for less than a date, matches if it should be present in a string, exact for an exact string match).
\end{itemize}

\item \vpedatatype{data:Datatypesorghl7fhirr5modelRange}{Range}: 
\begin{itemize}[noitemsep,nolistsep]
\item[] Attributes: long serialVersionUID, \ref{data:Datatypesorghl7fhirr5modelQuantity} low, \ref{data:Datatypesorghl7fhirr5modelQuantity} high
\item[] Base StructureDefinition for Range Type: A set of ordered Quantities defined by a low and high limit.
\end{itemize}

\item \vpedatatype{data:DatatypesRedeploymentJob}{RedeploymentJob}: 
\begin{itemize}[noitemsep,nolistsep]
\item[] Attributes: int subsystemId, \ref{data:DatatypesRedeploymentJobId} id, \ref{data:DatatypesTimestamp} startedAt, \ref{data:DatatypesRedeploymentStatus} status
\item[] Represents an active or completed redeployment or rollback 
  operation. Must complete within 5 minutes per Av2 response 
  measures.
\end{itemize}

\item \vpedatatype{data:DatatypesRedeploymentJobId}{RedeploymentJobId}: 
\begin{itemize}[noitemsep,nolistsep]

\item[] A unique identifier for a redeployment job.
\end{itemize}

\item \vpedatatype{data:DatatypesRedeploymentStatus}{RedeploymentStatus}: 
\begin{itemize}[noitemsep,nolistsep]

\item[] Enum: IN\_PROGRESS, COMPLETED, FAILED, ROLLED\_BACK.
  The current status of a redeployment or rollback operation 
  initiated by the system administrator.
\end{itemize}

\item \vpedatatype{data:DatatypesRetryState}{RetryState}: 
\begin{itemize}[noitemsep,nolistsep]
\item[] Attributes: \ref{data:DatatypesGatewayId} gatewayId, int attemtNumber, \ref{data:DatatypesTimestamp} nextRetryAt, \ref{data:DatatypesDuration} currentBackoffInterval
\item[]  Tracks the current exponential backoff retry state for a 
  patient gateway attempting to re-establish communication. 
  The currentBackoffInterval doubles with each failed attempt.
\end{itemize}

\item \vpedatatype{data:Datatypesorghl7fhirr5modelRiskAssessment}{RiskAssessment}: 
\begin{itemize}[noitemsep,nolistsep]
\item[] Attributes: List\textless{}\ref{data:Datatypesorghl7fhirr5modelIdentifier}\textgreater{} identifier, \ref{data:Datatypesorghl7fhirr5modelObservationStatus} status, \ref{data:Datatypesorghl7fhirr5modelCodeableConcept} method, \ref{data:Datatypesorghl7fhirr5modelCodeableConcept} code, \ref{data:Datatypesorghl7fhirr5modelPatient} subject, \ref{data:DatatypesDate} occurrence, String condition, \ref{data:Datatypesorghl7fhirr5modelIdentifier} performer, List\textless{}\ref{data:Datatypesorghl7fhirr5modelCodeableConcept}\textgreater{} reason, List\textless{}\ref{data:Datatypesorghl7fhirr5modelObservation}\textgreater{} basis, String mitigation, List\textless{}\ref{data:Datatypesorghl7fhirr5modelAnnotation}\textgreater{} note, long serialVersionUID, \ref{data:Datatypesorghl7fhirr5modelRiskAssessmentRiskAssessmentPredictionComponent} prediction
\item[] An assessment of the likely outcome(s) for a patient or other subject as well as the likelihood of each outcome.
Method specifies the algorithm, process or mechanism used to evaluate the risk.
Code specifies the type of the risk assessment performed.
For example, in the system  \url{http://snomed.info/sct}, `709510001' represents Assessment of risk for disease (procedure).
For assessments or prognosis specific to a particular condition, the condition attribute indicates the condition being assessed. It should be a reference to an existing condition but is a simplified as a string here.
Reason specifies the reason the risk assessment was performed.
Mitigation contains a description of the steps that might be taken to reduce the identified risk(s).
\end{itemize}

\item \vpedatatype{data:Datatypesorghl7fhirr5modelRiskAssessmentRiskAssessmentPredictionComponent}{RiskAssessmentPredictionComponent}: 
\begin{itemize}[noitemsep,nolistsep]
\item[] Attributes: String outcome, \ref{data:Datatypesorghl7fhirr5modelRange} probability, String qualitativeRisk, \ref{data:DatatypesBigDecimal} relativeRisk, \ref{data:DatatypesDate} when, String rationale, long serialVersionUID
\item[] Describes the expected outcome for the subject.
Outcome speciifes one of the potential outcomes for the patient (e.g. remission, death,  a particular condition).
QualitativeRisk indicates how likely the outcome is (in the specified timeframe), expressed as a qualitative value (e.g. low, medium, or high).
Rationale contains additional information explaining the basis for the prediction.
Both outcome and qualitativeRisk are CodeableConcepts (which can be formal reference to terminology or an ontology), but simplified as a string.
\end{itemize}

\item \vpedatatype{data:DatatypesRiskEstimationId}{RiskEstimationId}: 
\begin{itemize}[noitemsep,nolistsep]

\item[] A piece of data uniquely identifying a risk estimation performed by the PMS. This architecture does not specify the exact format of this identifier,
but possibilities are a long integer, a string, a URL etc.
\end{itemize}

\item \vpedatatype{data:DatatypesRiskEvent}{RiskEvent}: 
\begin{itemize}[noitemsep,nolistsep]
\item[] Attributes: \ref{data:DatatypesPatientId} patientId, \ref{data:DatatypesPatientStatus} status, \ref{data:DatatypesTimestamp} ts, \ref{data:DatatypesCorrelationId} correlationId
\item[] Payload pushed to IRiskEvents subscribers when setEstimatedPatientStatus actually changes the stored status.
\end{itemize}

\item \vpedatatype{data:DatatypesSensorDataElement}{SensorDataElement}: 
\begin{itemize}[noitemsep,nolistsep]
\item[] Attributes: String sensorId, String sensorType, String measurementType, String measurementUnit, \ref{data:DatatypesBigDecimal} measurementValue, String measurementValueString
\item[] Each sub-package lists the id of the sensor, the type of the sensor, the
type of the measurements and the measurements itself. These
measurements range from a single value (e.g., in case of the blood
pressure) to a complex data structure (e.g., in case of an ECG).
Non-numerical types of measurements can be captured in the measurementValueString attribute.
\end{itemize}

\item \vpedatatype{data:DatatypesSensorDataPackage}{SensorDataPackage}: 
\begin{itemize}[noitemsep,nolistsep]
\item[] Attributes: List\textless{}\ref{data:DatatypesSensorDataElement}\textgreater{} sensorDataElements
\item[] A package of sensor data, i.e., a list of sensorDataElements.

\end{itemize}

\item \vpedatatype{data:DatatypesSignalStrength}{SignalStrength}: 
\begin{itemize}[noitemsep,nolistsep]

\item[] Enum: NONE, WEAK, MODERATE, STRONG.
Represents the quality of the network signal available to the patient gateway. NONE indicates no signal, WEAK indicates signal insufficient for reliable data transmission, MODERATE indicates functional but potentially lossy connectivity, STRONG indicates reliable connectivity.
\end{itemize}

\item \vpedatatype{data:DatatypesSubscriberId}{SubscriberId}: 
\begin{itemize}[noitemsep,nolistsep]

\item[] Identifies a component-level subscriber (NotificationDispatcher, PhysicianCommandService). Distinct from PhysicianId.
\end{itemize}

\item \vpedatatype{data:DatatypesSubscriptionId}{SubscriptionId}: 
\begin{itemize}[noitemsep,nolistsep]

\item[] Handle for an active subscription.
\end{itemize}

\item \vpedatatype{data:DatatypesTimestamp}{Timestamp}: 
\begin{itemize}[noitemsep,nolistsep]

\item[] The representation of a time (i.e., date and time of the day) in the system.
\end{itemize}

\end{itemize}
% END TYPES

% CHECKS
\section{Unresolved issues}\label{sec:checks}
\emph{SA Plugin v10.0.0 (VP OpenAPI v17.2)}
\begin{itemize}[nolistsep,noitemsep]
\item \texttt{[Comp.03] }: 3 (\emph{The provided interface x::x is not required by any sibling components, but is also not provided by its parent component x})

\item \texttt{[Comp.04] }: 7 (\emph{The required interface x::x is not provided by any sibling components, but is also not required by its parent component x})

\item \texttt{[Comp.05] }: 2 (\emph{Component x has no provided or required interfaces})

\item \texttt{[Comp.07] }: 1 (\emph{The provided interface x::x is not provided by any child module})

\item \texttt{[Comp.08] }: 3 (\emph{The required interface x::x is not required by any child module})

\item \texttt{[Diag.comp.06] }: 24 (\emph{The component x may be missing on the C/S context diagram.})

\item \texttt{[Diag.depl.05] }: 17 (\emph{Component x at x cannot communicate with component providing x through a deployment connection})

\item \texttt{[Diag.depl.06] }: 4 (\emph{The deployment connection between x and x is not used by any of the components deployed in the two nodes.})

\item \texttt{[Diag.seq.01] }: 2 (\emph{x: x (x) calls x but does not require its interface.})

\item \texttt{[If.02] }: 1 (\emph{The interface x is not provided by any component or module})

\item \texttt{[If.03] }: 1 (\emph{The interface x is not required by any component or module})

\item \texttt{[Op.01] }: 1 (\emph{The interface x has no operations})

\item \texttt{[Op.03] }: 2 (\emph{Operation x::x uses non-existing data type x.})

\end{itemize}
% END CHECKS

%\end{document}